% ======================================================================
\chapter{Параметрические одномерные контрольные карты по количественным признакам}
% ======================================================================

\section*{Обзор главы}
В этой главе мы изучаем параметрические контрольные карты по количественным признакам. Это карты, которые основаны на предположении о лежащем в основе распределении процесса (таком как нормальное) или на некотором приближении (таком как нормальность, посредством центральной предельной теоремы) распределения статистики для нанесения на карту. Существуют три основных класса параметрических контрольных карт по количественным признакам: карта Шухарта, карта кумулятивных сумм (CUSUM) и карта экспоненциально взвешенного скользящего среднего (EWMA), каждая из которых обычно используется для обнаружения сдвигов определённого типа. Мы подробно описываем некоторые из этих карт, а также относительные преимущества и недостатки этих карт, которые хорошо задокументированы в литературе и затрагиваются далее. Этот материал полезен, поскольку аналоги многих из этих параметрических карт по количественным признакам рассматриваются в Главе 4 для непараметрической постановки.

Карты по количественным признакам основаны на статистиках для нанесения на карту, имеющих непрерывное распределение, и мы сосредоточимся на них здесь. Помимо карт по количественным признакам, существуют также карты по альтернативным признакам, основанные на статистиках для нанесения на карту с дискретным распределением. К ним относятся, например, карта доли несоответствующих единиц (называемая \(p\)-картой) и контрольная карта для числа несоответствий (называемая \(c\)-картой), как описано у Montgomery (2009, с. 290, 308). Эти контрольные карты, хотя и весьма полезны на практике, в данной книге не рассматриваются. В этой главе также рассматриваются различные способы повышения чувствительности контрольной карты (такие как правила серий или правила чувствительности, предупреждающие границы, интерпретация паттернов на контрольной карте и т.д.), и в завершение обсуждается робастность параметрических контрольных карт.

\section{Введение}
Пусть \(X_{i1}, X_{i2}, \ldots, X_{in}\) обозначают случайную выборку (измерения некоторой количественной непрерывной характеристики качества или переменной, такой как температура, диаметр и т.д.) объёма \(n \geq 1\), взятую из процесса в момент времени \(i = 1, 2, 3, \ldots\). Для параметрических контрольных карт по количественным признакам предполагается, что (форма) лежащего в основе распределения процесса известна, и наиболее распространённым предположением на практике является нормальность, хотя возможны и изучались другие распределения, такие как экспоненциальное, Вейбулла и т.д. Однако, поскольку большинство контрольных карт в литературе предназначены для нормального распределения, мы рассматриваем только их и с этого момента называем эти карты параметрическими картами или картами нормальной теории. Предположим, что процесс следует нормальному распределению со средним \(\mu\) и дисперсией \(\sigma^2\), которые могут быть известны (заданы) или неизвестны. Если параметры процесса в состоянии IC известны или заданы, эта ситуация называется \textbf{Случаем K} (Case K) и обычно обозначается добавлением нижнего индекса ноль к символам параметров (например, \(\mu_0\) и \(\sigma_0^2\)) для среднего и дисперсии соответственно. Если параметры процесса в состоянии IC неизвестны, их необходимо оценить до начала мониторинга. Это обычно делается в ходе ретроспективного анализа имеющихся данных, называемого исследованием Фазы I, где процесс стабилизируется или приводится под контроль, и таким образом формируется набор эталонных данных до того, как может начаться проспективный мониторинг процесса в Фазе II (см. Раздел 2.1.12 для более подробной информации о Фазе I и Фазе II). Эта ситуация, когда параметры процесса неизвестны, называется \textbf{Случаем U} (Case U), а среднее и дисперсия процесса, оценённые по эталонным данным Фазы I, обычно обозначаются как \(\hat{\mu}\) и \(\hat{\sigma}^2\) соответственно.

Сначала мы обсуждаем три базовые параметрические контрольные карты по количественным признакам для случая известных параметров. Случай неизвестных параметров обсуждается позже, в Разделе 3.8.

\section{Параметрические контрольные карты по количественным признакам в случае K}
Напомним, что существуют три основных класса параметрических контрольных карт по количественным признакам: карта Шухарта, карта CUSUM и карта EWMA, каждая из которых обычно используется для обнаружения сдвигов определённого типа. Мы более подробно описываем эти карты в каждом из трёх последующих разделов.

\subsection{Контрольные карты Шухарта}
Среди множества контрольных карт, используемых на практике, карты Шухарта являются наиболее популярными благодаря своей простоте, лёгкости применения и тому, что эти универсальные карты достаточно эффективны для обнаружения умеренных и больших сдвигов. Эти карты были впервые предложены Уолтером Шухартом в 1926 году. Чтобы описать карту Шухарта в общем виде, предположим, что параметр положения процесса \(\theta\), такой как среднее значение, должен отслеживаться с помощью статистики для нанесения на карту \(T\), которая, с статистической точки зрения, является хорошей точечной оценкой \(\theta\). Предположим далее, что ожидаемое значение и стандартное отклонение \(T\) равны \(\mu_T\) и \(\sigma_T\) соответственно. Статистические соображения часто приводят нас к выбору \(T\) в качестве несмещённой оценки \(\theta\), так что \(\mu_T = \theta\). Тогда общая формула для центральной линии (\(CL\)) и контрольных границ карты Шухарта имеет вид

\begin{equation}
	\begin{aligned}
		UCL &= \theta + k\sigma_T \\
		CL &= \theta \\
		LCL &= \theta - k\sigma_T
	\end{aligned}
	\label{eq:3.1}
\end{equation}

где \(k > 0\) — это константа для нанесения на карту, которая является параметром проектирования карты и определяет «расстояние» контрольных границ от \(CL\), выраженное в единицах стандартного отклонения. Следовательно, эти контрольные границы часто называют \(k\)-сигма границами. \(k\)-сигма контрольная карта Шухарта — это график, который отображает эти три границы в виде прямых линий вместе с реализованными (вычисленными) значениями статистики для нанесения на карту \(T\) для ряда выборок или во времени. Обратите внимание, что на карте Шухарта верхняя и нижняя контрольные границы расположены симметрично относительно \(CL\). Такие контрольные границы более осмысленны, когда распределение \(T\) симметрично или приблизительно симметрично, что хорошо согласуется с предположением в карте Шухарта о том, что либо распределение процесса является нормальным, либо \(T\) имеет распределение, которое приблизительно нормально со средним \(\theta\). Предположим, например, что \(\theta\) — это среднее процесса \(\mu\), которое нужно отслеживать, а значение \(\mu\) в состоянии IC равно \(\mu_0\). В этом случае \(T\) берётся в качестве среднего \(\overline{X}\) выборки, и тогда \(k\)-сигма границы задаются как \(\mu_0 \pm k\sigma_0/\sqrt{n}\), где \(\sigma_0\) — известное стандартное отклонение процесса, поскольку \(\sigma_{\overline{X}} = \sigma_0/\sqrt{n}\). Обоснование \(k\)-сигма границ состоит в том, что \(\overline{X}\) является точно (или приблизительно) нормально распределённой величиной, когда распределение процесса нормально (или в силу центральной предельной теоремы). Когда статистика для нанесения на карту попадает на или за пределы верхней или нижней контрольной границы, мы говорим, что наблюдался сигнал или произошло сигнальное событие, и процесс объявляется вышедшим из-под контроля (OOC).

В зависимости от цели обнаружения увеличения и/или уменьшения параметра процесса \(\mu_0\), контрольные карты часто проектируются в односторонней или двусторонней форме. Односторонняя верхняя (или нижняя) карта используется для обнаружения восходящих (или нисходящих) изменений, в то время как двусторонняя контрольная карта используется для обнаружения любого изменения процесса, которое может быть как восходящим, так и нисходящим.

Таким образом, использование односторонней контрольной карты подразумевает, что нас интересует сдвиг или изменение в известном направлении, то есть увеличение или уменьшение. Нижняя односторонняя карта Шухарта сигнализирует о том, что произошёл нисходящий сдвиг, если \(T_i \leq LCL\) (карта имеет только \(LCL\)), тогда как верхняя односторонняя карта Шухарта сигнализирует о том, что произошёл восходящий сдвиг, если \(T_i \geq UCL\) (карта имеет только \(UCL\)). В дальнейшем мы рассматриваем двусторонние карты, поскольку они более общие (так как практик часто может не знать направления сдвига) и поскольку именно они получили наибольшее внимание в литературе. Более того, двусторонние карты можно легко адаптировать для получения односторонних карт.

В литературе были рассмотрены различные усовершенствования и улучшения карты Шухарта; некоторые из них мы обсуждаем в последующих разделах. Иллюстративный пример контрольной карты Шухарта также приводится ниже.

\noindent\textbf{Замечание 3.1. Константа для нанесения на карту: принцип 3-сигма}

Шухарт (1926, 1931) рекомендовал брать константу для нанесения на карту \(k\) равной 3, что даёт хорошо известную 3-сигма контрольную карту Шухарта \(\overline{X}\). Очевидно, что такой выбор предполагает, явно или неявно, лежащее в основе нормальное распределение. Следовательно, 3-сигма контрольная карта Шухарта \(\overline{X}\) является параметрической контрольной картой. Выбор \(k\) обычно зависит от некоторого критерия, связанного с эффективностью карты. В своём обсуждении Шухарт обосновал выбор \(k = 3\), используя комбинацию как эмпирических, так и теоретических данных. Предполагая, что лежащий в основе процесс распределён нормально, используя свойства нормального распределения, имеем

\[
\begin{aligned}
	P\left(\mu_0 - 3\frac{\sigma_0}{\sqrt{n}} < \overline{X} < \mu_0 + 3\frac{\sigma_0}{\sqrt{n}}\right) &= P\left(-3 < \frac{\overline{X} - \mu_0}{\frac{\sigma_0}{\sqrt{n}}} < 3\right) \\
	&= \Phi(3) - \Phi(-3) = 0.9973
\end{aligned}
\]

где \(\Phi\) обозначает функцию распределения (cdf) стандартного нормального распределения. Это означает, что \(99.73\%\) значений выборочных средних попадут в диапазон \(\pm 3\sigma_{\overline{X}}\) от среднего в состоянии IC; принцип 3-сигма проиллюстрирован на рисунке \ref{fig-3.1}. Другими словами, вероятность ложной тревоги, то есть вероятность того, что статистика для нанесения на карту \(\overline{X}\) попадёт на или за пределы контрольных границ, когда процесс фактически находится в состоянии IC, равна \(1 - 0.9973 = 0.0027\) точно, когда процесс нормален, или приблизительно, когда объём выборки достаточно велик. Вероятность ложной тревоги называется частотой ложных тревог (FAR). Значение 0.0027 является отраслевым стандартом и может быть найдено из таблиц D1 и D2 в Приложении A. Значение FAR 0.0027 означает, что ложная тревога ожидается в 27 из 10 000 или 3 из 1000 выборок, но возможны и другие значения, которые могут быть желательными в зависимости от конкретной ситуации.

Некоторые авторы, см., например, Nelson (2003), задавались вопросом: «Когда границы на контрольной карте Шухарта должны быть иными, чем \(CL \pm 3\)-сигма?» Nelson (2003) отвечает на этот вопрос так: «Всякий раз, когда для этого есть веская причина». Здесь мы развиваем это. Когда константа для нанесения на карту \(k\) меньше 3, контрольная карта имеет более узкие контрольные границы, что приводит к более высокой FAR. И наоборот, когда константа для нанесения на карту \(k\) больше 3, контрольная карта имеет более широкие контрольные границы, что приводит к более низкой FAR. В литературе есть и другие аргументы относительно того, когда \(k \neq 3\) не следует использовать.

Известно, что карты Шухарта более эффективны для обнаружения внезапных или резких больших сдвигов (см., например, Montgomery, 2009, с. 400).

\begin{figure}[!ht]
	\centering
	\includegraphics[width=0.7\linewidth]{"./images/fig-3.1.png"}
	\caption[]{Принцип 3-сигма}
	\label{fig-3.1}
\end{figure}

\subsection{Контрольные карты CUSUM}
Хотя карты Шухарта широко известны и часто используются на практике благодаря своей простоте и эффективности в обнаружении умеренных и больших сдвигов, другие карты, такие как карты CUSUM, могут быть более полезны в определённых ситуациях для обнаружения меньших, устойчивых сдвигов. Эти карты, иногда называемые картами с взвешиванием по времени, более естественно подходят для среды контроля процесса с учётом последовательного характера сбора данных. Контрольные карты CUSUM были впервые введены Пейджем (1954) (хотя и не в их нынешнем виде) и изучались многими авторами на протяжении последних 60 лет. См., например, Barnard (1959), Ewan and Kemp (1960), Johnson (1961), Goldsmith and Whitfield (1961), Page (1961), Ewan (1963), Hawkins (1992, 1993) и Hawkins and Olwell (1998). Это лишь некоторые примеры, и с момента введения карт CUSUM в 1954 году Пейджем было проделано невероятно много работы по картам CUSUM (см. обзор в Encyclopaedia of Statistics in Quality and Reliability под редакцией Ruggeri, Kenett и Faltin (2007a) и цитируемые там источники, например). Карты CUSUM обычно основаны на кумулятивных суммах (CUSUM) некоторой статистики или разностей статистики и её ожидаемого значения в состоянии IC, и рассчитываются последовательно по мере накопления данных во времени. Например, карта CUSUM для среднего обычно основана на кумулятивной сумме отклонений отдельных наблюдений (или средних подгрупп) от заданного значения целевого среднего в состоянии IC.

Чтобы описать карту CUSUM более подробно, предположим, как и ранее, что \(X_{i1}, X_{i2}, X_{i3}, \ldots, X_{in}\) обозначают \(i\)-ю выборку или выборку (подгруппу) объёма \(n \geq 1\) в момент выборки (время) \(i\), \(i = 1, 2, \ldots\), из процесса с известным или заданным средним процесса в состоянии IC \(\mu_0\) и известным или заданным стандартным отклонением процесса в состоянии IC \(\sigma_0\). Пусть

\begin{equation}
	T_i = \psi(x_{i1}, x_{i2}, \ldots, x_{in})
	\label{eq:3.2}
\end{equation}

— статистика, построенная по данным \(i\)-й выборки, \(i = 1, 2, \ldots\), где \(\psi\) — некоторая функция. Статистика в уравнении \eqref{eq:3.2} называется базовой (опорной) статистикой, которая была бы оценкой среднего процесса; например, \(T_i\) может быть выборочным средним. Как мы отмечали ранее, можно формировать либо CUSUM выборочных средних (для сгруппированных в подгруппы данных), либо CUSUM отдельных наблюдений (\(n = 1\)), либо CUSUM разностей выборочных средних или отдельных наблюдений и \(\mu_0\). Ниже описывается карта CUSUM для отдельных наблюдений. Здесь \(n = 1\) и \(T_i = X_i\). Адаптация к выборочным средним проста и будет прокомментирована позже. Карта CUSUM формируется путём нанесения на график \(C_i\), где

\[
C_i = \sum_{j=1}^{i}(x_j - \mu_0)
\]

что можно записать как

\[
\begin{aligned}
	C_i &= (x_i - \mu_0) + \sum_{j=1}^{i-1}(x_j - \mu_0) \\
	C_i &= (x_i - \mu_0) + C_{i-1}.
\end{aligned}
\]

Существуют два способа представления CUSUM, а именно: (i) табличная (или алгоритмическая) форма CUSUM и (ii) форма CUSUM с V-маской. Мы рассматриваем только табличную форму CUSUM, поскольку форма CUSUM с V-маской больше не используется на практике. Montgomery (2009, с. 403) утверждает: «Из двух представлений табличная форма CUSUM предпочтительнее». Поскольку мы рассматриваем только табличную форму CUSUM, слово «табличная» с этого момента будет опускаться, и мы будем называть её просто картой CUSUM.

Верхняя односторонняя карта CUSUM используется для обнаружения увеличений среднего и работает путём накопления разностей между \(X_i\) и \(\mu_0\). Следовательно, для верхней односторонней карты CUSUM мы используем статистику для нанесения на карту

\begin{equation}
	C_i^{+} = \max[0, X_i - \mu_0 - K + C_{i-1}^{+}] \quad \text{для } i = 1, 2, 3, \ldots
	\label{eq:3.3}
\end{equation}

для обнаружения положительных сдвигов (увеличений) среднего от \(\mu_0\), с начальным значением \(C_0^+ = 0\) и где \(K \geq 0\) называется опорным или свободным значением. Обычно \(K\) выражается через стандартное отклонение процесса, так что \(K = k\sigma_0\). Контрольная карта подаёт сигнал для первого \(i\), такого что \(C_i^+ \geq H\), где константа \(H > 0\) называется решающим интервалом.

Аналогично, нижняя односторонняя карта CUSUM представляет интерес для обнаружения уменьшений среднего и работает путём накопления разностей от \(\mu_0\), которые ниже целевого значения

\begin{equation}
	C_i^{-} = \min[0, X_i - \mu_0 + K + C_{i-1}^{-}] \quad \text{для } i = 1, 2, 3, \ldots
	\label{eq:3.4}
\end{equation}

или

\begin{equation}
	C_i^{-*} = \max[0, \mu_0 - K - X_i + C_{i-1}^{-*}] \quad \text{для } i = 1, 2, 3, \ldots
	\label{eq:3.5}
\end{equation}

и используется для обнаружения отрицательных сдвигов или отклонений от \(\mu_0\) с начальным значением \(C_0^- = C_0^{-*} = 0\). Здесь сигнальное событие происходит для первого \(i\), такого что \(C_i^- \leq -H\) (если используется уравнение \eqref{eq:3.4}) или \(C_i^{-*} \geq H\) (если используется уравнение \eqref{eq:3.5}). Для более визуально привлекательной карты уравнение \eqref{eq:3.4} будет использоваться в данной книге для построения нижней односторонней карты CUSUM с этого момента. Двусторонняя карта CUSUM представляет интерес, когда важны как увеличения, так и уменьшения среднего. Эта карта подаёт сигнал для первого \(i\), при котором выполняется любое из двух неравенств: \(C_i^+ \geq H\) или \(C_i^- \leq -H\). Обратите внимание, что для карты CUSUM существуют счётчики \(N^+\) и \(N^-\), которые указывают количество последовательных периодов, в течение которых CUSUM \(C_i^+\) и \(C_i^-\) были ненулевыми до подачи сигнала, что помогает определить, в какой момент времени мог произойти сдвиг. Если сигнал подан верхней CUSUM (т.е. \(C_i^+\)) на номере выборки 29 и, скажем, соответствующий счётчик равен 7 (т.е. \(N^+ = 7\)), то мы заключим, что процесс в последний раз был в состоянии IC на номере выборки 29 - 7 = 22, поэтому сдвиг, вероятно, произошёл между номерами выборок 22 и 23. Карты CUSUM имеют дополнительное преимущество: можно получить оценку нового среднего процесса \(\hat{\mu}\) после сдвига. Величина \(\frac{C_i^+}{N^+}\) является оценкой того, насколько текущее среднее после сдвига превышает \(\mu_0 + K\), когда сигнал подаётся с \(C_i^+\). Величина \(\frac{C_i^-}{N^-}\) является оценкой того, насколько текущее среднее после сдвига ниже \(\mu_0 - K\), когда сигнал подаётся с \(C_i^-\). Таким образом, если верхняя CUSUM подала сигнал, мы оцениваем новое среднее процесса как

\begin{equation}
	\hat{\mu} = \mu_0 + K + \frac{C_i^+}{N^+}
	\label{eq:3.6}
\end{equation}

тогда как, если нижняя CUSUM подала сигнал, мы оцениваем новое среднее процесса как

\begin{equation}
	\hat{\mu} = \mu_0 - K - \frac{C_i^-}{N^-}
	\label{eq:3.7}
\end{equation}

Для реализации карты CUSUM необходимы константы \(K\) и \(H\); это обсуждается далее.

Константы \(K\) и \(H\) называются параметрами проектирования карты CUSUM и обычно выбираются так, чтобы карта имела заданную номинальную \(ARL_{\text{IC}}\) (например, 370) и была способна обнаружить сдвиг среднего, особенно малый сдвиг, как можно скорее. Первым шагом в этом направлении является выбор \(K\), который называется опорным значением. Поскольку \(K = k\sigma_0\), а \(\sigma_0\) предполагается известным, задавая \(k\), мы фактически задаём \(K\). Соответственно, с этого момента мы рассматриваем только выбор \(k\). Отметим, что последнее замечание также относится к параметру проектирования \(H\), поскольку \(H\) также обычно определяется как \(H = h\sigma_0\) с \(h = 4\) или 5, обеспечивающими хорошие свойства \(ARL\) (см. Montgomery, 2009, с. 408); таким образом, задавая \(h\), мы фактически задаём \(H\). Для параметра проектирования \(k\) рассмотрим параметрическую карту CUSUM для мониторинга известного среднего нормального распределения (\(\mu_0 = 0\), без потери общности) со стандартным отклонением \(\sigma_0 = 1\) (опять же, без потери общности), на основе отдельных данных (\(n = 1\)). В этом случае распределение \(X_i\) в состоянии IC следует нормальному распределению со средним \(\mu_0 = 0\) и стандартным отклонением \(\sigma_0 = 1\). Чтобы изучить влияние выбора \(k\), мы рассматриваем график некоторых значений \(ARL_{\text{OOC}}\), рассчитанных для нормального распределения методом имитационного моделирования, на рисунке \ref{fig-3.2}, устанавливая номинальную \(ARL_{\text{IC}} = 500\), для \(\mu_1 = 0.1\), 0.25, 0.5 и 1.0. Отметим, что \(\mu_1\) представляет увеличенные значения \(\mu\), которые должны быть обнаружены «быстро» от \(\mu_0 = 0\); следовательно, \(\mu_1\) представляет истинный сдвиг среднего, где \(\mu_1 = \mu_0 + \delta\sigma_0\) с \(\delta = 0.1\), 0.25, 0.5 и 1.0 соответственно.

Из рисунка \ref{fig-3.2} можно сделать несколько интересных наблюдений. С одной стороны, когда сдвиг мал (см. панели (a) и (b)) и выбрано большее значение \(k\), значения \(ARL_{\text{OOC}}\) становятся неприемлемо высокими. С другой стороны, если сдвиг велик (см. панели (c) и (d)) и выбрано меньшее значение \(k\), значения \(ARL_{\text{OOC}}\) также высоки, но не настолько, как в предыдущем случае. Это говорит о том, что когда имеется мало априорной информации о размере сдвига или она отсутствует, меньшее значение \(k\) является более безопасным выбором (чтобы защититься от любой ненужной задержки в обнаружении).

\begin{figure}[!ht]
	\centering
	\includegraphics[width=0.7\linewidth]{"./images/fig-3.2.png"}
	\caption[]{Значения \(ARL_{\text{OOC}}\) традиционной карты CUSUM с номинальной \(ARL_{\text{IC}} = 500\) для различных значений \(k\) и \(\mu_1 = 0.1\), 0.25, 0.5 и 1.0}
	\label{fig-3.2}
\end{figure}

Выбор \(k\) для параметрической карты CUSUM для нормального среднего обсуждался многими авторами; см., например, Lucas (1985), Hawkins and Olwell (1998), Kim et al. (2007) и Montgomery (2009). Lucas (1985) утверждает: «Параметр CUSUM \(k\) определяется приемлемым средним уровнем (\(\mu_a\)) и неприемлемым средним уровнем (\(\mu_d\)), который схема CUSUM должна быстро обнаружить. Для нормально распределённых переменных значение \(k\) выбирается на полпути между приемлемым средним уровнем и неприемлемым средним уровнем». В более поздней литературе, см., например, Montgomery (2009), принято, что в рамках нормальной теории \(k\) обычно выбирается относительно размера сдвига, который мы хотим обнаружить, то есть \(k = \frac{1}{2}\delta\), где \(\delta\) — размер сдвига среднего, выраженный в единицах стандартного отклонения.

После выбора значения \(k\) следующим шагом является нахождение решающего интервала \(h\) в сочетании с выбранным \(k\) так, чтобы достигалась заданная номинальная \(ARL_{\text{IC}}\). Отметим, однако, что для дискретной случайной величины вероятность того, что \(h\) не всегда может быть найдено таким образом, чтобы заданная номинальная \(ARL_{\text{IC}}\) достигалась точно, и поэтому, используя консервативный подход, \(h\) находится так, чтобы достигнутая \(ARL_{\text{IC}}\) была меньше или равна заданной номинальной \(ARL_{\text{IC}}\). Решающий интервал \(h\) может быть найден с помощью алгоритма поиска по сетке, используя, скажем, 100 000 имитаций Монте-Карло с помощью некоторого статистического программного обеспечения, такого как SAS\(^{\text{\textregistered}}\) или R\(^{\text{\textregistered}}\).

На этом этапе следует отметить, что способность обнаружения карты CUSUM зависит от правильного проектирования (настройки) карты. Правильное проектирование карты CUSUM включает получение оптимальной комбинации параметров карты CUSUM (\(k\) и \(h\)) путём минимизации средней длины серии в состоянии OOC (обозначаемой \(ARL_{\text{OOC}}\)) для заданного значения размера сдвига \(\delta\) и номинальной средней длины серии в состоянии IC (обозначаемой \(ARL_{\text{IC}}\)). Другими словами, можно получить много комбинаций \(k\) и \(h\), которые достигают заданной номинальной \(ARL_{\text{IC}}\), но оптимальной парой является та, для которой \(ARL_{\text{OOC}}\) наименьшая.

Отметим, что это общее обсуждение карты CUSUM относится к случаю K, то есть когда параметры процесса известны. В случае U параметры процесса неизвестны и должны быть оценены. Естественно, это повлияет на выбор \(k\) и \(h\). Читатель отсылается к Jones, Champ и Rigdon (2004), которые оценили эффективность карты CUSUM с оценёнными параметрами. Авторы также предоставляют метод аппроксимации распределения длины серии и его моментов.

\begin{table}[htbp]
	\centering
	\caption{Параметры карты CUSUM: значения \((h, k)\) для номинальных значений \(ARL_{\text{IC}}\) 370 и 500 соответственно.}
	\label{tab:3.1}
	\renewcommand{\arraystretch}{1.2} % Увеличиваем межстрочный интервал
	\begin{tabular}{c c *{6}{r}}
		\toprule
		\multirow{2}{*}{370} & \(k\) & 0.25 & 0.50 & 0.75 & 1.00 & 1.25 & 1.50 \\
		& \(h\) & 8.01 & 4.77 & 3.34 & 2.52 & 1.99 & 1.61 \\
		\addlinespace % Добавляет небольшой вертикальный отступ между блоками
		\multirow{2}{*}{500} & \(k\) & 0.25 & 0.50 & 0.75 & 1.00 & 1.25 & 1.50 \\
		& \(h\) & 7.27 & 4.39 & 3.08 & 2.32 & 1.83 & 1.47 \\
		\bottomrule
	\end{tabular}
	\end{table}

В этом разделе мы обсудили карту CUSUM для отдельных измерений. Однако если берутся рациональные подгруппы объёма \(n > 1\), то просто замените \(X_i\) на \(\overline{X}_i\) и \(\sigma_0\) на \(\sigma_{\overline{X}} = \frac{\sigma_0}{\sqrt{n}}\) в приведённых выше уравнениях.

Читатель отсылается к Hawkins and Olwell (1998) для подробного обзора параметрических карт CUSUM. Они также предоставляют программное обеспечение и некоторые очень полезные таблицы со значениями \(ARL_{\text{IC}}\) для различных значений \(h\) и \(k\); см. их таблицы 3.1 и 3.2 на страницах 48 и 49 соответственно. Hawkins (1993) дополнительно предоставляет таблицу значений \((h, k)\) для номинального значения \(ARL_{\text{IC}}\) 370. Используя эти таблицы, мы приводим значения \((h, k)\) для двусторонней карты CUSUM в таблице \ref{tab:3.1} для номинальных значений \(ARL_{\text{IC}}\) 370 и 500 соответственно. Например, для сдвига примерно \(1\sigma_0\) в среднем процесса, взяв \(k = 0.5\) и \(h = 4.77\), получаем \(ARL_{\text{IC}} = 370\).

Следует отметить, что, как и в случае с картой Шухарта, было проделано много работы по улучшению эффективности карты CUSUM. Одним из таких улучшений является функция быстрого начального отклика (FIR). Улучшение FIR было первоначально предложено в Lucas and Crosier (1982) для параметрической карты CUSUM. Функция FIR, или «форы» (head-start, HS), используется, когда хотят повысить чувствительность CUSUM при запуске процесса. Это делается путём установки начальных значений \(C_0^+\) и \(C_0^-\) равными некоторому ненулевому значению, обычно мы устанавливаем \(C_0^+ = \frac{H}{2}\) и \(C_0^- = -\frac{H}{2}\). Это называется 50\%-ной форой.

При таком объёме работы, проделанной с параметрическими контрольными картами, естественно рассмотреть аналоги с использованием непараметрических статистик для нанесения на карту. Этот подход привёл к непараметрическим картам CUSUM (обозначаемым NPCUSUM), которые обсуждаются в следующей главе.

\subsection{Контрольные карты EWMA}
Другим популярным классом контрольных карт являются карты EWMA. Карты EWMA также используют преимущество последовательно (упорядоченного во времени) накапливающегося характера данных, возникающего в типичной среде SPC, и известны своей эффективностью в обнаружении малых сдвигов, но при этом могут быть проще в реализации, чем карты CUSUM (см., например, Montgomery, 2009, с. 419). Классические карты EWMA для среднего были введены Робертсом (1959), и они содержат карты Шухарта как частный случай. Литература по картам EWMA огромна и продолжает расти значительными темпами (см., например, обзор в Encyclopaedia of Statistics in Quality and Reliability под редакцией Ruggeri, Kenett и Faltin (2007a) и цитируемые там источники). Некоторые более поздние ссылки включают Capizzi and Masarotto (2012) и Ross et al. (2012).

Чтобы описать карту EWMA более подробно, предположим, как и ранее, что \(X_{i1}, X_{i2}, X_{i3}, \ldots, X_{in}\) обозначают случайную выборку (подгруппу) объёма \(n \geq 1\) по выходу процесса в каждый момент выборки \(i = 1, 2, \ldots\), из процесса с известным средним процесса в состоянии IC \(\mu_0\) и известным стандартным отклонением процесса в состоянии IC \(\sigma_0\). Как и в случае карты CUSUM, мы рассматриваем случай отдельных наблюдений, так что \(n = 1\), и отслеживаем \(X_1, X_2, X_3, \ldots\) с \(T_i = X_i\) для \(i = 1, 2, \ldots\). Адаптация к выборочным средним проста и будет прокомментирована позже.

Статистика для нанесения на карту для контрольной карты EWMA определяется как

\begin{equation}
	Z_i = \lambda X_i + (1 - \lambda)Z_{i-1} \quad \text{для } i = 1, 2, 3, \ldots
	\label{eq:3.8}
\end{equation}

где \(0 < \lambda \leq 1\) — константа, называемая параметром сглаживания. Начальное значение \(Z_0\) обычно берётся равным значению среднего процесса в состоянии IC, то есть \(Z_0 = \mu_0\).

Чтобы показать, что статистика для нанесения на карту \(Z_i\) является взвешенным средним всех предыдущих статистик, \(Z_{i-1}\) можно заменить на \(Z_{i-1} = \lambda X_{i-1} + (1 - \lambda)Z_{i-2}\) в уравнении \eqref{eq:3.8}, чтобы получить

\[
\begin{aligned}
	Z_i &= \lambda X_i + (1 - \lambda)(\lambda X_{i-1} + (1 - \lambda)Z_{i-2}) \\
	&= \lambda X_i + \lambda(1 - \lambda)X_{i-1} + (1 - \lambda)^2 Z_{i-2} \\
	&= \lambda X_i + \lambda(1 - \lambda)X_{i-1} + (1 - \lambda)^2(\lambda X_{i-2} + (1 - \lambda)Z_{i-3}) \\
	&= \lambda X_i + \lambda(1 - \lambda)X_{i-1} + \lambda(1 - \lambda)^2 X_{i-2} + (1 - \lambda)^3 Z_{i-3}.
\end{aligned}
\label{eq:3.8b}
\]

Этот метод подстановки называется рекурсивной подстановкой. Продолжая процесс рекурсивной подстановки для \(Z_{i-j}\), \(j = 2, 3, \ldots\), получаем

\[
Z_i = \lambda \sum_{j=0}^{i-1}(1 - \lambda)^j X_{i-j} + (1 - \lambda)^i Z_0.
\]

Ожидаемое значение и дисперсия статистики для нанесения на карту \(Z_i\) задаются как

\begin{equation}
	E(Z_i) = \mu_0
	\label{eq:3.9}
\end{equation}

и

\begin{equation}
	var(Z_i) = \sigma_0^2\left(\frac{\lambda}{2 - \lambda}\right)(1 - (1 - \lambda)^{2i}) \quad \text{для } i = 1, 2, \ldots
	\label{eq:3.10}
\end{equation}

соответственно (см. Приложение 3.1 для выводов). Следовательно, точные контрольные границы и \(CL\) контрольной карты EWMA задаются как

\begin{equation}
	\begin{aligned}
		UCL &= E(Z_i) + L \times SD(Z_i) = \mu_0 + L\sigma_0\sqrt{\left(\frac{\lambda}{2 - \lambda}\right)(1 - (1 - \lambda)^{2i})} \\
		CL &= E(Z_i) = \mu_0 \\
		LCL &= E(Z_i) - L \times SD(Z_i) = \mu_0 - L\sigma_0\sqrt{\left(\frac{\lambda}{2 - \lambda}\right)(1 - (1 - \lambda)^{2i})}
	\end{aligned}
	\label{eq:3.11}
\end{equation}

где \(L > 0\) — константа для нанесения на карту. Обратите внимание, что контрольные границы EWMA устанавливаются на расстоянии \(L\) стандартных отклонений от среднего в состоянии IC (ожидаемого значения статистики для нанесения на карту в состоянии IC). Когда \(\lambda = 1\), контрольные границы сводятся к границам карты Шухарта \(\mu_0 \pm L\sigma_0\), и карта EWMA сводится к карте Шухарта.

Расчёт и реализация точных контрольных границ EWMA могут быть несколько громоздкими. В качестве альтернативы можно использовать так называемые контрольные границы стационарного состояния (которые обычно используются, когда карта EWMA работает в течение нескольких периодов времени, так что член \((1 - (1 - \lambda)^{2i})\) в уравнении \eqref{eq:3.11} стремится к единице), задаваемые как

\begin{equation}
	\begin{aligned}
		UCL &= \mu_0 + L\sigma_0\sqrt{\left(\frac{\lambda}{2 - \lambda}\right)} \\
		\text{и} \\
		LCL &= \mu_0 - L\sigma_0\sqrt{\left(\frac{\lambda}{2 - \lambda}\right)}.
	\end{aligned}
	\label{eq:3.12}
\end{equation}

Двусторонняя карта EWMA строится путём нанесения на график \(Z_i\) в зависимости от номера выборки \(i\) (или времени). Если статистика для нанесения на карту \(Z_i\) попадает между двумя контрольными границами, то есть \(LCL < Z_i < UCL\), процесс считается находящимся в состоянии IC. Если статистика для нанесения на карту \(Z_i\) попадает на или за пределы одной из контрольных границ, то есть \(Z_i \leq LCL\) или \(Z_i \geq UCL\), процесс считается находящимся в состоянии OOC, и необходимо провести поиск назначаемых причин.

Двусторонняя карта EWMA может быть легко модифицирована для построения односторонней карты EWMA почти так же, как это было сделано для карты CUSUM. Например, верхняя односторонняя статистика для нанесения на карту EWMA задаётся как \(Z_i^{+} = \max[\mu_0, \lambda X_i + (1 - \lambda) Z_{i-1}]\) для \(i = 1, 2, 3, \ldots\) с начальным значением \(Z_0^{+} = \mu_0\). Если \(Z_i\) попадает на или выше \(UCL\), процесс объявляется находящимся в состоянии OOC, и необходимо провести поиск назначаемых причин.

Также, как и в случае карты CUSUM, параметры проектирования карты EWMA, \(L\) и \(\lambda\), выбираются так, чтобы карта имела заданную номинальную \(ARL_{\text{IC}}\) и была способна обнаружить заданную величину сдвига, особенно малого сдвига, как можно скорее, в терминах кратчайшей \(ARL_{\text{OOC}}\). Montgomery (2009, с. 422) утверждает, что «Оптимальная процедура проектирования должна состоять в задании желаемых средних длин серий в состоянии IC и OOC и величины ожидаемого сдвига процесса, а затем в выборе комбинации \(\lambda\) и \(L\), обеспечивающей желаемые характеристики \(ARL\)». Константа \(\lambda\) (\(0 < \lambda \leq 1\)) является параметром сглаживания и выбирается в зависимости от величины сдвига, который нужно обнаружить. Константа \(L > 0\) — это расстояние контрольных границ от \(CL\) (чем больше значение \(L\), тем шире контрольные границы и наоборот), и она выбирается в сочетании со значением параметра сглаживания \(\lambda\). Что касается реализации карты EWMA, первым шагом является выбор \(\lambda\). Рекомендуется выбирать малое \(\lambda\), скажем, равное 0.05, когда представляют интерес малые сдвиги. Если большее беспокойство вызывают умеренные сдвиги, выбирайте \(\lambda = 0.10\), тогда как \(\lambda = 0.20\) следует выбирать, если представляют интерес большие сдвиги (см., например, Montgomery, 2009, с. 423). После выбора \(\lambda\) второй шаг включает выбор \(L\) так, чтобы достигалась желаемая \(ARL_{\text{IC}}\). Это можно сделать с помощью алгоритма поиска по сетке, используя, скажем, 100 000 имитаций Монте-Карло с помощью некоторого статистического программного обеспечения, такого как SAS® или R®.

В таблице \ref{tab:3.2} приведены некоторые комбинации \((\lambda, L)\), которые дают значения \(ARL_{\text{IC}}\), близкие к желаемым значениям 370 и 500 соответственно, для отдельных данных. Они были найдены с использованием имитаций Монте-Карло, как описано выше.

\begin{table}[htbp]
	\centering
	\caption{Комбинации \((\lambda, L)\), дающие значения \(ARL_{\text{IC}}\), близкие к желаемым номинальным значениям 370 и 500.}
	\label{tab:3.2}
	\renewcommand{\arraystretch}{1.1} % Небольшое увеличение высоты строк
	\begin{tabular}{lll}
		\toprule
		Размер сдвига & Номинальная \(ARL_{\text{IC}} = 370\) & Номинальная \(ARL_{\text{IC}} = 500\) \\
		\midrule
		Малый    & (0.05, 2.492) & (0.05, 2.615) \\
		Умеренный & (0.10, 2.703) & (0.10, 2.814) \\
		Большой    & (0.20, 2.860) & (0.20, 2.962) \\
		&  Не указано для \(\lambda = 0.25\) & (0.25, 2.998) \\
		& Не указано для \(\lambda = 0.40\) & (0.40, 3.054) \\
		\bottomrule
	\end{tabular}
\end{table}	

Способность обнаружения сдвига карты EWMA, как и карты CUSUM, зависит от правильного проектирования (настройки) карты. Правильное проектирование карты EWMA включает получение (оптимальной) комбинации параметров карты EWMA (\(\lambda\) и \(L\)) путём минимизации \(ARL_{\text{OOC}}\) для заданного значения размера сдвига \(\delta\) и для заданной номинальной \(ARL_{\text{IC}}\). Другими словами, можно получить много комбинаций \(\lambda\) и \(L\), которые дают заданную номинальную \(ARL_{\text{IC}}\). Оптимальной парой из этих пар является та, для которой \(ARL_{\text{OOC}}\) наименьшая.

Отметим, что это общее обсуждение карты EWMA относится к случаю K, то есть когда значения параметров процесса в состоянии IC известны. В случае U параметры процесса неизвестны и должны быть оценены. Процесс оценивания влияет на эффективность карты и параметры проектирования карты, и поэтому оптимальный выбор должен быть пересмотрен в этом случае. Читатель отсылается к Jones, Champ и Rigdon (2001), которые оценили эффективность карты EWMA с оценёнными параметрами. Мы затронем вопрос оценки параметров позже в этой главе.

В этом разделе мы обсудили карту EWMA для отдельных измерений. Однако если берутся рациональные подгруппы объёма \(n > 1\), то просто замените \(X_i\) на \(\overline{X}_i\) и \(\sigma_0\) на \(\sigma_{\overline{X}} = \frac{\sigma_0}{\sqrt{n}}\) в приведённых выше уравнениях.

Было проделано много исследований для улучшения эффективности карты EWMA путём добавления улучшений. Одним из них является функция FIR. Улучшение FIR было первоначально предложено в Lucas and Crosier (1982) для параметрической карты CUSUM. Затем Lucas and Saccucci (1990) предложили аналогичную функцию для параметрической карты EWMA. Функция FIR используется как средство против проблем запуска и тех процессов, которые не имеют корректирующих действий после предыдущего сигнала OOC. Lucas and Saccucci (1990) утверждали, что для карты EWMA функция FIR наиболее полезна, когда \(\lambda\) мало. В таблице \ref{tab:3.3} мы приводим различные типы функций FIR для карты EWMA, найденные в литературе на сегодняшний день.

\begin{table}[htbp]
	\centering
	\caption{Различные типы улучшений FIR для карт EWMA и соответствующие статьи.}
	\label{tab:3.3}
	\renewcommand{\arraystretch}{1.2}
	\begin{tabular}{p{4.5cm}p{5.5cm}p{3.5cm}}
		\toprule
		Тип улучшения FIR & Описание & Статья \\
		\midrule
		(i) Фиксированные контрольные границы с форой (HS) & \(0\% < \text{HS} \leq 100\%\) & Lucas and Saccucci (1990) \\
		(ii) Зависящие от времени (TV) контрольные границы с HS & \(0\% < \text{HS} \leq 100\%\) & Rhoads, Montgomery, and Mastrangelo (1996) \\
		(iii) TV контрольные границы с FIR экспоненциального типа & \(\text{FIR}_{\text{radj}} = 1 - (1 - f)^{1+a(t-1)}\) & Steiner (1999) \\
		(iv) TV контрольные границы с модифицированным FIR экспоненциального типа & \(\text{MFIR}_{\text{radj}} = (1 - (1 - f)^{1+a(t-1)})^{1+\frac{1}{t}}\) & Haq, Brown, and Moltchanova (2014) \\
		\bottomrule
	\end{tabular}
\end{table}

При таком объёме работы, доступной по параметрическим контрольным картам, явно естественно рассмотреть аналоги этих карт с использованием непараметрических статистик для нанесения на карту. Эта линия мышления привела к нескольким непараметрическим картам EWMA (NPEWMA), которые обсуждаются в следующей главе.

Прежде чем перейти к более подробному рассмотрению и приведению иллюстративных примеров по параметрическим картам CUSUM и EWMA, мы сначала сосредоточимся на параметрических контрольных картах Шухарта.

\section{Типы параметрических карт по количественным признакам в случае K: иллюстративные примеры}

\subsection{Контрольные карты Шухарта}

\subsubsection{Контрольные карты Шухарта для мониторинга среднего процесса}
Для мониторинга среднего процесса мы обычно используем выборочное среднее (\(\overline{X}\)). Далее следует пример.

\noindent\textbf{Пример 3.1. Параметрическая контрольная карта Шухарта \(\overline{X}\)}

Карты Шухарта обычно применяются с данными подгрупп. Столбец (a) таблицы \ref{tab:3.4} содержит некоторые смоделированные данные из нормального распределения, которые представляют собой измерения, взятые из 25 независимых выборок, каждая объёма 5 (\(n = 5\)), по некоторому типу пластины. Предположим, что из инженерных соображений средний размер пластин в состоянии IC \(\mu_0\) задан равным \(1.5\) см, а стандартное отклонение процесса в состоянии IC \(\sigma_0\) известно и равно \(0.1\) см. Среднее значение каждой выборки показано в столбце (b) таблицы \ref{tab:3.4}. Чтобы проиллюстрировать расчёты, рассмотрим выборку номер 1. Первая статистика для нанесения на карту, среднее выборки 1, рассчитывается следующим образом:

\[
\overline{X}_1 = \frac{1.3235 + 1.4128 + 1.6744 + 1.4573 + 1.6914}{5} = \frac{7.5594}{5} = 1.5119.
\]

Чтобы применить формулы для контрольных границ из уравнения \eqref{eq:3.1}, отметим, что ожидаемое значение \(\overline{X}\) (т.е. \(\mu_{\overline{X}}\)) — это просто заданное среднее процесса \(\mu_0\), тогда как стандартное отклонение \(\overline{X}\), а именно \(\sigma_{\overline{X}}\), равно \(\frac{\sigma_0}{\sqrt{n}}\). Таким образом, \(CL\) и \(k\)-сигма контрольные границы для контрольной карты Шухарта \(\overline{X}\) задаются как

\begin{equation}
	\begin{aligned}
		UCL &= \mu_0 + k\frac{\sigma_0}{\sqrt{n}} \\
		CL &= \mu_0 \\
		LCL &= \mu_0 - k\frac{\sigma_0}{\sqrt{n}}.
	\end{aligned}
	\label{eq:3.13}
\end{equation}

\begin{table}[htbp]
	\centering
	\caption{Измерения, использованные для построения контрольной карты Шухарта \(\overline{X}\), и средние подгрупп.}
	\label{tab:3.4}
	\renewcommand{\arraystretch}{1.05} % Уменьшаем высоту строк, так как их много
	\begin{tabular}{c *{5}{r} r}
		\toprule
		Выборка & \multicolumn{5}{c}{(a)} & (b) \\
		\cmidrule(lr){2-6}
		\(i\) & \(X_1\) & \(X_2\) & \(X_3\) & \(X_4\) & \(X_5\) & \(\overline{X}_i\) \\
		\midrule
		1  & 1.3235 & 1.4128 & 1.6744 & 1.4573 & 1.6914 & 1.5119 \\
		2  & 1.4314 & 1.3592 & 1.6075 & 1.4666 & 1.6109 & 1.4951 \\
		3  & 1.4284 & 1.4871 & 1.4932 & 1.4324 & 1.5674 & 1.4817 \\
		4  & 1.5028 & 1.6352 & 1.3841 & 1.2831 & 1.5507 & 1.4712 \\
		5  & 1.5604 & 1.2735 & 1.5265 & 1.4363 & 1.6441 & 1.4882 \\
		6  & 1.5955 & 1.5451 & 1.3574 & 1.3281 & 1.4198 & 1.4492 \\
		7  & 1.6274 & 1.5064 & 1.8366 & 1.4177 & 1.5144 & 1.5805 \\
		8  & 1.4190 & 1.4303 & 1.6637 & 1.6067 & 1.5519 & 1.5343 \\
		9  & 1.3884 & 1.7277 & 1.5355 & 1.5176 & 1.3688 & 1.5076 \\
		10 & 1.4039 & 1.6697 & 1.5089 & 1.4627 & 1.5220 & 1.5134 \\
		11 & 1.4158 & 1.7667 & 1.4278 & 1.5928 & 1.4181 & 1.5242 \\
		12 & 1.5821 & 1.3355 & 1.5777 & 1.3908 & 1.7559 & 1.5284 \\
		13 & 1.2856 & 1.4106 & 1.4447 & 1.6398 & 1.1928 & 1.3947 \\
		14 & 1.4951 & 1.4036 & 1.5893 & 1.6458 & 1.4969 & 1.5261 \\
		15 & 1.3589 & 1.2863 & 1.5996 & 1.2497 & 1.5471 & 1.4083 \\
		16 & 1.5747 & 1.5301 & 1.5171 & 1.1839 & 1.8662 & 1.5344 \\
		17 & 1.3680 & 1.7269 & 1.3957 & 1.5014 & 1.4449 & 1.4874 \\
		18 & 1.4163 & 1.3864 & 1.3057 & 1.6210 & 1.5573 & 1.4573 \\
		19 & 1.5796 & 1.4185 & 1.6541 & 1.5116 & 1.7247 & 1.5777 \\
		20 & 1.7106 & 1.4412 & 1.2361 & 1.3820 & 1.7601 & 1.5060 \\
		21 & 1.4371 & 1.5051 & 1.3485 & 1.5670 & 1.4880 & 1.4691 \\
		22 & 1.4738 & 1.5936 & 1.6583 & 1.4973 & 1.4720 & 1.5390 \\
		23 & 1.5917 & 1.4333 & 1.5551 & 1.5295 & 1.6866 & 1.5592 \\
		24 & 1.6399 & 1.5243 & 1.5705 & 1.5563 & 1.5530 & 1.5688 \\
		25 & 1.5797 & 1.3663 & 1.6240 & 1.3732 & 1.6887 & 1.5264 \\
		\bottomrule
	\end{tabular}
\end{table}

Подставляя \(\mu_0 = 1.5\), \(\sigma_0 = 0.1\), \(n = 5\) и \(k = 3\) в уравнение \eqref{eq:3.13}, получаем \(CL\) и контрольные границы для карты Шухарта \(\overline{X}\):

\[
\begin{aligned}
	UCL &= 1.5 + 3\frac{0.1}{\sqrt{5}} = 1.634 \\
	CL &= 1.5 \\
	LCL &= 1.5 - 3\frac{0.1}{\sqrt{5}} = 1.366.
\end{aligned}
\]

На рисунке \ref{fig-3.3} показаны выборочные средние \(\overline{X}_i\) из столбца (b) таблицы \ref{tab:3.4}, которые являются статистиками для нанесения на карту для контрольной карты Шухарта \(\overline{X}\), нанесённые на карту вместе с контрольными границами. Это отображение, созданное с помощью статистического программного пакета Minitab®, версия 17, обычно называется контрольной картой Шухарта \(\overline{X}\) или картой «X-bar».

После построения контрольной карты принимается решение о состоянии процесса. На рисунке \ref{fig-3.3} ни одно из выборочных средних не выходит за пределы контрольных границ, и точки, по-видимому, не демонстрируют какого-либо неслучайного паттерна. Следовательно, мы заключаем, что процесс функционирует в состоянии статистической управляемости (IC) относительно целевого среднего.

\begin{figure}[!ht]
	\centering
	\includegraphics[width=0.7\linewidth]{"./images/fig-3.3.png"}
	\caption[]{Контрольная карта Шухарта \(\overline{X}\) для мониторинга положения процесса}
	\label{fig-3.3}
\end{figure}

\subsubsection{Контрольные карты Шухарта для мониторинга вариации процесса}
Вариация является важным аспектом любого анализа, и поэтому необходимо отслеживать вариацию или разброс процесса и гарантировать, что она находится в состоянии IC. Более того, как мы видели в уравнении \eqref{eq:3.1}, контрольные границы Шухарта для среднего процесса зависят от стандартного отклонения процесса. Таким образом, если стандартное отклонение не остаётся в состоянии IC, контрольная карта для среднего не будет информативной. Поэтому нам нужно отслеживать дисперсию или стандартное отклонение с помощью контрольной карты.

Существует несколько возможных статистик, которые можно использовать для мониторинга вариации. Наиболее популярными вариантами являются выборочный размах (\(R\)), выборочное стандартное отклонение (\(S\)) и выборочная дисперсия (\(S^2\)).

Обычно мы используем контрольную карту для мониторинга среднего процесса вместе с контрольной картой для мониторинга вариации процесса. Если вариация находится в состоянии IC, мы переходим к рассмотрению контрольной карты для среднего. Например, карта Шухарта \(\overline{X}\) для среднего часто используется вместе с картой Шухарта \(R\) для разброса. Отметим, что для иллюстрации мы рассматриваем карту Шухарта \(R\), даже несмотря на то, что недавняя литература рекомендует использовать другую карту разброса, такую как карта Шухарта \(S\); см., например, Mahmoud et al. (2010). Мы делаем это потому, что карта Шухарта \(R\) проста и продолжает использоваться в промышленности.

В случае K значения \(\mu\) и \(\sigma\) известны или заданы, так что они могут быть использованы для построения соответствующих контрольных карт. Мы иллюстрируем карты Шухарта \(R\) и \(S\) для известного стандартного отклонения \(\sigma_0\).

\noindent\textbf{Пример 3.2. Параметрическая контрольная карта Шухарта \(R\)}

Те же данные, которые использовались для иллюстрации карты Шухарта \(\overline{X}\), используются для иллюстрации карты Шухарта \(R\). Столбец (a) таблицы \ref{tab:3.5} содержит измерения, взятые из 25 независимых выборок по пластинам, каждая объёма \((n = 5)\), из нормального распределения. Размах каждой выборки показан в столбце (b) таблицы \ref{tab:3.5}. Стандартное отклонение каждой выборки также приведено в столбце (c) таблицы \ref{tab:3.5}. Однако оно будет использоваться только в примере 3.3, где иллюстрируется карта Шухарта \(S\). Для карты Шухарта \(R\), чтобы проиллюстрировать расчёты, рассмотрим выборку номер 1. Первая статистика для нанесения на карту, размах выборки 1, рассчитывается следующим образом:

\[
R_1 = \max_1 - \min_1 = 1.6914 - 1.3235 = 0.3679.
\]

Для контрольной карты Шухарта \(R\) \(CL\) и контрольные границы задаются как

\begin{equation}
	\begin{aligned}
		UCL &= D_2\sigma_0 \\
		CL &= d_2\sigma_0 \\
		LCL &= D_1\sigma_0
	\end{aligned}
	\label{eq:3.14}
\end{equation}

где \(D_1\), \(d_2\) и \(D_2\) — константы, зависящие от объёма выборки \(n\) (см. таблицу C в Приложении A для этих констант). Для \(n = 5\) получаем \(D_1 = 0\), \(d_2 = 2.326\) и \(D_2 = 4.918\). Предполагая заданное значение \(\sigma_0 = 0.30\) и подставляя его в уравнение \eqref{eq:3.14}, получаем \(CL\) и контрольные границы для данных из таблицы \ref{tab:3.5}:

\[
\begin{aligned}
	UCL &= D_2\sigma_0 = (4.918)(0.3) = 1.4754 \\
	CL &= d_2\sigma_0 = (2.326)(0.3) = 0.6978 \\
	LCL &= D_1\sigma_0 = (0)(0.3) = 0.
\end{aligned}
\]

На рисунке \ref{fig-3.4} показаны выборочные размахи \(R_i\) из таблицы \ref{tab:3.5}, столбец (b), которые являются статистиками для нанесения на карту для контрольной карты Шухарта \(R\), нанесённые на карту вместе с контрольными границами.

На рисунке \ref{fig-3.4} все точки отображаются между контрольными границами, но поскольку все точки находятся на уровне \(CL\) или ниже, это вызывает сильное подозрение, что вариация процесса не находится в состоянии IC.

\begin{table}[htbp]
	\centering
	\caption{Измерения, использованные для построения контрольных карт Шухарта \(R\) и \(S\), и статистики для нанесения на карту.}
	\label{tab:3.5}
	\renewcommand{\arraystretch}{1.05} % Уменьшаем высоту строк, так как их много
	\begin{tabular}{c *{5}{r} r r}
		\toprule
		Выборка & \multicolumn{5}{c}{(a)} & (b) & (c) \\
		\cmidrule(lr){2-6}
		\(i\) & \(X_1\) & \(X_2\) & \(X_3\) & \(X_4\) & \(X_5\) & \(R_i\) & \(S_i\) \\
		\midrule
		1  & 1.3235 & 1.4128 & 1.6744 & 1.4573 & 1.6914 & 0.3679 & 0.1635 \\
		2  & 1.4314 & 1.3592 & 1.6075 & 1.4666 & 1.6109 & 0.2517 & 0.1111 \\
		3  & 1.4284 & 1.4871 & 1.4932 & 1.4324 & 1.5674 & 0.1390 & 0.0565 \\
		4  & 1.5028 & 1.6352 & 1.3841 & 1.2831 & 1.5507 & 0.3521 & 0.1389 \\
		5  & 1.5604 & 1.2735 & 1.5265 & 1.4363 & 1.6441 & 0.3706 & 0.1412 \\
		6  & 1.5955 & 1.5451 & 1.3574 & 1.3281 & 1.4198 & 0.2674 & 0.1168 \\
		7  & 1.6274 & 1.5064 & 1.8366 & 1.4177 & 1.5144 & 0.4189 & 0.1614 \\
		8  & 1.4190 & 1.4303 & 1.6637 & 1.6067 & 1.5519 & 0.2447 & 0.1077 \\
		9  & 1.3884 & 1.7277 & 1.5355 & 1.5176 & 1.3688 & 0.3589 & 0.1439 \\
		10 & 1.4039 & 1.6697 & 1.5089 & 1.4627 & 1.5220 & 0.2658 & 0.0988 \\
		11 & 1.4158 & 1.7667 & 1.4278 & 1.5928 & 1.4181 & 0.3509 & 0.1548 \\
		12 & 1.5821 & 1.3355 & 1.5777 & 1.3908 & 1.7559 & 0.4204 & 0.1682 \\
		13 & 1.2856 & 1.4106 & 1.4447 & 1.6398 & 1.1928 & 0.4470 & 0.1699 \\
		14 & 1.4951 & 1.4036 & 1.5893 & 1.6458 & 1.4969 & 0.2422 & 0.0937 \\
		15 & 1.3589 & 1.2863 & 1.5996 & 1.2497 & 1.5471 & 0.3499 & 0.1568 \\
		16 & 1.5747 & 1.5301 & 1.5171 & 1.1839 & 1.8662 & 0.6823 & 0.2423 \\
		17 & 1.3680 & 1.7269 & 1.3957 & 1.5014 & 1.4449 & 0.3589 & 0.1432 \\
		18 & 1.4163 & 1.3864 & 1.3057 & 1.6210 & 1.5573 & 0.3153 & 0.1289 \\
		19 & 1.5796 & 1.4185 & 1.6541 & 1.5116 & 1.7247 & 0.3062 & 0.1195 \\
		20 & 1.7106 & 1.4412 & 1.2361 & 1.3820 & 1.7601 & 0.5240 & 0.2230 \\
		21 & 1.4371 & 1.5051 & 1.3485 & 1.5670 & 1.4880 & 0.2185 & 0.0819 \\
		22 & 1.4738 & 1.5936 & 1.6583 & 1.4973 & 1.4720 & 0.1863 & 0.0832 \\
		23 & 1.5917 & 1.4333 & 1.5551 & 1.5295 & 1.6866 & 0.2533 & 0.0922 \\
		24 & 1.6399 & 1.5243 & 1.5705 & 1.5563 & 1.5530 & 0.1156 & 0.0431 \\
		25 & 1.5797 & 1.3663 & 1.6240 & 1.3732 & 1.6887 & 0.3224 & 0.1482 \\
		\bottomrule
	\end{tabular}
\end{table}

\noindent\textbf{Пример 3.3. Параметрическая контрольная карта Шухарта \(S\)}

Те же данные используются для иллюстрации карты Шухарта \(S\). Напомним, что столбец (a) таблицы \ref{tab:3.5} содержит измерения, взятые из 25 независимых выборок по пластинам, каждая объёма \((n = 5)\), из нормального распределения. Стандартное отклонение каждой выборки показано в столбце (c) таблицы \ref{tab:3.5}. Чтобы проиллюстрировать расчёты, рассмотрим выборку номер 1. Первая статистика для нанесения на карту, стандартное отклонение выборки 1, равна 0.1635. Для контрольной карты Шухарта \(S\) \(CL\) и контрольные границы задаются (Montgomery, 2009, с. 253)

\begin{equation}
	\begin{aligned}
		UCL &= B_6\sigma_0 \\
		CL &= c_4\sigma_0 \\
		LCL &= B_5\sigma_0
	\end{aligned}
	\label{eq:3.15}
\end{equation}

где \(B_5 = c_4 - 3\sqrt{1 - c_4^2}\) и \(B_6 = c_4 + 3\sqrt{1 - c_4^2}\), а \(c_4\) — константа, зависящая от объёма выборки \(n\) (см. таблицу C Приложения A для этих констант).

Для \(n = 5\) получаем \(B_6 = 1.964\), \(c_4 = 0.94\) и \(B_5 = 0\) из таблицы C (Приложение A), так что, снова предполагая заданное значение \(\sigma_0 = 0.30\) и используя уравнение \eqref{eq:3.15}, \(CL\) и контрольные границы для карты \(S\) для данных из таблицы \ref{tab:3.5}:

\[
\begin{aligned}
	UCL &= (1.964)(0.3) = 0.5892 \\
	CL &= (0.94)(0.3) = 0.282 \\
	LCL &= (0)(0.3) = 0.
\end{aligned}
\]

На рисунке \ref{fig-3.5} показаны выборочные стандартные отклонения \(S_i\) из столбца (c) таблицы \ref{tab:3.5}, которые являются статистиками для нанесения на карту для контрольной карты Шухарта \(S\), нанесённые на карту вместе с контрольными границами.

Рассматривая рисунок \ref{fig-3.5}, мы приходим к тому же выводу, что и с картой \(R\): вариация процесса не находится в состоянии IC, даже несмотря на то, что все точки отображаются между контрольными границами и в данных нет аномальных паттернов.

\begin{figure}[!ht]
	\centering
	\includegraphics[width=0.7\linewidth]{"./images/fig-3.4.png"}
	\caption[]{Контрольная карта Шухарта \(R\) для мониторинга вариации процесса}
	\label{fig-3.4}
\end{figure}

\begin{figure}[!ht]
	\centering
	\includegraphics[width=0.7\linewidth]{"./images/fig-3.5.png"}
	\caption[]{Контрольная карта Шухарта \(S\) для мониторинга вариации процесса}
	\label{fig-3.5}
\end{figure}

Часто на практике контрольные границы Шухарта выражаются в виде кратных стандартного отклонения (или стандартной ошибки) статистики для нанесения на карту. Множитель 3, конечно же, является наиболее распространённым. Представленные выше карты Шухарта \(R\) и \(S\) используют множитель 3 и основаны на том факте, что распределение статистики для нанесения на карту (\(R\) или \(S\)) может быть аппроксимировано нормальным распределением с помощью центральной предельной теоремы. Хотя этот подход может быть разумным в некоторых ситуациях и использовался в течение многих лет, для обычно рекомендуемых объёмов выборки, таких как 5, нижняя контрольная граница оказывается отрицательной и поэтому устанавливается равной 0, поскольку стандартное отклонение \(\sigma\) не может быть отрицательным. Это практическое «исправление» создаёт небольшую проблему, поскольку контрольная карта с нижней контрольной границей, равной нулю, не способна обнаруживать какие-либо уменьшения вариации процесса. Обе представленные здесь 3-сигма карты \(R\) и \(S\) имеют этот недостаток для объёмов выборки меньше или равных 6. Одним из решений этой проблемы является использование вероятностных границ, где константы для нанесения на карту находятся для заданной FAR. Таким образом, для номинальной FAR, скажем, 0.002, нам нужно найти 0.001 и 0.999 процентные точки распределения \(R/\sigma\) (или \(S/\sigma\)) в состоянии IC и построить контрольные границы на их основе. Предполагая нормальное распределение, такие процентили были рассчитаны и приведены в Grant and Leavenworth (1980) для некоторых номинальных значений FAR. Montgomery (2009, с. 242) утверждает, что в США наиболее распространено использование \(k\)-сигма контрольных границ, тогда как в Великобритании и некоторых частях Западной Европы более распространены вероятностные границы. В нашем предстоящем разделе о контрольных картах дисперсии для случая неизвестных параметров мы используем вероятностные границы.

Важным вопросом в отношении мониторинга среднего процесса является тот факт, что, поскольку стандартное отклонение фигурирует в контрольных границах для среднего, мы сначала исследуем стандартное отклонение, и когда будет установлено, что оно находится в состоянии IC, мы исследуем карту для среднего. Таким образом, можно утверждать, что на самом деле и среднее, и стандартное отклонение отслеживаются совместно (или одновременно) при мониторинге среднего. Следовательно, как видно из приведённых выше примеров, две карты (\(\overline{X}, R\)) или (\(\overline{X}, S\)) используются одновременно, по сути как одна схема нанесения на карту, для принятия решения о состоянии (IC или OOC) процесса. Процесс считается находящимся в состоянии IC всякий раз, когда обе карты отображаются в пределах своих соответствующих контрольных границ и демонстрируют случайные паттерны. И наоборот, процесс объявляется находящимся в состоянии OOC, когда хотя бы одна из карт показывает ситуацию OOC. Если стандартное отклонение находится в состоянии OOC, нет смысла рассматривать карту среднего; если стандартное отклонение находится в состоянии IC, карта среднего исследуется, чтобы определить, находится ли среднее процесса в состоянии IC или нет. Согласно текущей практике, карта Шухарта \(\overline{X}\) для среднего и карта Шухарта \(R\) для стандартного отклонения строятся независимо, каждая с использованием 3-сигма границ. Это распространённая практика, рекомендуемая в большинстве учебников. Однако это может быть проблематичным, поскольку FAR схемы будет значительно выше (почти вдвое), чем у каждой из двух карт по отдельности, что является проблемой. Здесь возникает вопрос эффекта неучтённого множественного тестирования. Поскольку в процессе принятия решений используются две карты одновременно, если не будут внесены надлежащие корректировки в контрольные границы, схема нанесения на карту не обеспечит заданную номинальную \(ARL_{\text{IC}}\). Фактически, некоторые недавние работы показали, что достигнутые значения могут быть намного ниже номинального значения, что означает, что ложных тревог будет гораздо больше, чем номинально ожидается. Это может значительно снизить эффективность контрольных карт. Читатель отсылается к Diko, Chakraborti и Graham (2016) для получения более подробной информации и некоторых исправленных (скорректированных) констант для нанесения на карту.

\subsection{Контрольные карты CUSUM}

Далее следует пример параметрической контрольной карты CUSUM для индивидуальных данных.

\noindent\textbf{Пример 3.4. Параметрическая контрольная карта CUSUM}

Столбец (a) таблицы \ref{tab:3.6} содержит 30 измерений, сгенерированных из процесса, который имеет нормальное распределение. Первые 20 из этих наблюдений были случайно извлечены из нормального распределения со средним \(\mu_0 = 10\) и стандартным отклонением \(\sigma_0 = 1\). Последние 10 наблюдений были извлечены из нормального распределения со средним \(\mu_1 = 11\) и стандартным отклонением \(\sigma_0 = 1\). Таким образом, \(\mu_1 = \mu_0 + \delta\sigma_0 = 10 + (1)(1)\), так что \(\delta = 1\). Чтобы проиллюстрировать карту CUSUM, предположим, что нас интересует обнаружение увеличения (сдвига) среднего на величину \(1\sigma_0\), что является сдвигом среднего размера. Отметим, что здесь мы можем использовать \(K\) и \(k\) взаимозаменяемо, поскольку \(\sigma_0 = 1\). То же самое рассуждение справедливо для \(H\) и \(h\). Hawkins (1993; таблица 1, с. 465) рекомендует, чтобы для сдвига примерно на \(1\sigma_0\) в среднем процесса значения \(k = 0.5\) и \(h = 4.77\) давали \(ARL_{\text{IC}} = 370\). Соответственно, мы применяем параметрическую CUSUM с \(k = 0.5\) и \(h = 4.77\).

Статистики для нанесения на карту CUSUM, \(C_i^+\) и \(C_i^-\), приведены в столбцах (b) и (d) таблицы \ref{tab:3.6} соответственно, тогда как соответствующие величины \(N^+\) и \(N^-\) приведены в столбцах (c) и (e) соответственно. Чтобы проиллюстрировать расчёты, рассмотрим период 1. Статистики для нанесения на карту для первого периода рассчитываются следующим образом:

\[
\begin{aligned}
	C_1^+ &= \max[0, X_1 - (\mu_0 + K) + C_0^+] = \max[0, 9.45 - (10 + 0.5) + 0] \\
	&= \max[0, -1.05] = 0
\end{aligned}
\]

и

\[
\begin{aligned}
	C_1^- &= \min[0, X_1 - \mu_0 + K + C_0^-] = \min[0, 9.45 - 10 + 0.5 + 0] \\
	&= \min[0, -0.05] = 0.
\end{aligned}
\]

На рисунке \ref{fig-3.6} показаны статистики CUSUM \(C_i^+\) и \(C_i^-\) из столбцов (a) и (c) соответственно, которые нанесены на параметрическую контрольную карту CUSUM вместе с решающим интервалом (\(H = 4.77\)).

Из рисунка \ref{fig-3.6} видно, что карта CUSUM подаёт сигнал на номере выборки/периоде 29. Карта CUSUM также указывает, когда, вероятно, произошёл сдвиг. Поскольку \(N^+ = 7\) на периоде 29, мы заключаем, что процесс в последний раз был в состоянии IC на периоде \(29 - 7 = 22\), поэтому сдвиг, вероятно, произошёл между периодами 22 и 23.

\begin{table}[htbp]
	\centering
	\caption{Данные и статистики для нанесения на карту, использованные для построения параметрической контрольной карты CUSUM.}
	\label{tab:3.6}
	\small % Уменьшенный шрифт, как в книге
	\setlength{\tabcolsep}{5pt} % Плотные отступы между колонками
	\renewcommand{\arraystretch}{1.0} % Минимальная высота строк
	
	\begin{tabular}{c r r r r r}
		\toprule
		Номер & \multicolumn{1}{c}{(a)} & \multicolumn{1}{c}{(b)} & \multicolumn{1}{c}{(c)} & \multicolumn{1}{c}{(d)} & \multicolumn{1}{c}{(e)} \\
		наблюдения, & & & & & \\
		\(i\) & \(x_i\) & \(C_i^+\) & \(N^+\) & \(C_i^-\) & \(N^-\) \\
		\midrule
		& & 0 & & 0 & \\
		1  & 9.45  & 0 & 0 & -0.05 & 1 \\
		2  & 7.99  & 0 & 0 & -1.56 & 2 \\
		3  & 9.29  & 0 & 0 & -1.77 & 3 \\
		4  & 11.66 & 1.16 & 1 & 0 & 0 \\
		5  & 12.16 & 2.82 & 2 & 0 & 0 \\
		6  & 10.18 & 2.50 & 3 & 0 & 0 \\
		7  & 8.04  & 0.04 & 4 & -1.46 & 1 \\
		8  & 11.46 & 1.00 & 5 & 0 & 0 \\
		9  & 9.20  & 0 & 0 & -0.30 & 1 \\
		10 & 10.34 & 0 & 0 & 0 & 0 \\
		11 & 9.03  & 0 & 0 & -0.47 & 1 \\
		12 & 11.47 & 0.97 & 1 & 0 & 0 \\
		13 & 10.51 & 0.98 & 2 & 0 & 0 \\
		14 & 9.40  & 0 & 0 & -0.10 & 1 \\
		15 & 10.08 & 0 & 0 & 0 & 0 \\
		16 & 9.37  & 0 & 0 & -0.13 & 1 \\
		17 & 10.62 & 0.12 & 1 & 0 & 0 \\
		18 & 10.31 & 0 & 0 & 0 & 0 \\
		19 & 8.52  & 0 & 0 & -0.98 & 1 \\
		20 & 10.84 & 0.34 & 1 & 0 & 0 \\
		21 & 10.90 & 0.74 & 2 & 0 & 0 \\
		22 & 9.33  & 0 & 0 & -0.17 & 1 \\
		23 & 12.29 & 1.79 & 1 & 0 & 0 \\
		24 & 11.50 & 2.79 & 2 & 0 & 0 \\
		25 & 10.60 & 2.89 & 3 & 0 & 0 \\
		26 & 11.08 & 3.47 & 4 & 0 & 0 \\
		27 & 10.38 & 3.35 & 5 & 0 & 0 \\
		28 & 11.62 & 4.47 & 6 & 0 & 0 \\
		29 & 11.31 & 5.28 & 7 & 0 & 0 \\
		30 & 10.52 & 5.30 & 8 & 0 & 0 \\
		\bottomrule
	\end{tabular}
\end{table}

\begin{figure}[!ht]
	\centering	\includegraphics[width=0.7\linewidth]{"./images/fig-3.6.png"}
	\caption[]{Параметрическая контрольная карта CUSUM}
	\label{fig-3.6}
\end{figure}

В этом примере мы проиллюстрировали параметрическую CUSUM для отдельных измерений. Однако если берутся рациональные подгруппы объёма \(n > 1\), то мы просто заменяем \(X_i\) на \(\overline{X}_i\) и \(\sigma_0\) на \(\sigma_{\overline{X}} = \frac{\sigma_0}{\sqrt{n}}\) в предыдущих уравнениях. В этом случае оптимальная пара \(k\) и \(h\) может отличаться от пары, использованной в этом примере для индивидуальных данных; например, для \(n = 1\) Hawkins (1993) рекомендует, чтобы для сдвига примерно на \(1\sigma_0\) в среднем процесса значения \(k = 0.5\) и \(h = 4.77\) давали \(ARL_{\text{IC}} = 370\). Используя алгоритм поиска, мы находим, что для \(n = 5\) при тех же значениях \(k\) и \(h\), то есть \(k = 0.5\) и \(h = 4.77\), получаем \(ARL_{\text{IC}} = 295\). Используя алгоритм поиска, мы находим, что для объёма выборки 5 следует использовать пару \((k, h) = (0.5, 4.8)\), чтобы получить \(ARL_{\text{IC}} = 370\). Однако отметим, что эта карта CUSUM использует смоделированные данные из нормального распределения, что может быть трудно обосновать на практике, особенно для индивидуальных данных.

\subsection{Контрольные карты EWMA}

Далее следует пример параметрической контрольной карты EWMA для индивидуальных данных.

\noindent\textbf{Пример 3.5. Параметрическая контрольная карта EWMA}

Те же данные, которые использовались для иллюстрации параметрической карты CUSUM, используются для иллюстрации параметрической карты EWMA. Столбец (a) таблицы \ref{tab:3.7} содержит 30 измерений, взятых для процесса, который имеет нормальное распределение. Первые 20 из этих наблюдений были случайно извлечены из нормального распределения со средним \(\mu_0 = 10\) и стандартным отклонением \(\sigma_0 = 1\). Как и ранее, последние 10 наблюдений были извлечены из нормального распределения со средним \(\mu_1 = \mu_0 + \delta\sigma_0 = 10 + (1)(1) = 11\) и стандартным отклонением \(\sigma_0 = 1\). Предположим, нас интересует обнаружение сдвига размером \(1\sigma_0\), то есть среднего сдвига. Статистика для нанесения на карту EWMA показана в столбце (b), тогда как зависящие от времени нижняя и верхняя контрольные границы показаны в столбцах (c) и (d) соответственно. Мы применим параметрическую контрольную карту EWMA с \(\lambda = 0.1\), \(L = 2.703\) (см. таблицу \ref{tab:3.2}; это даёт значение \(ARL_{\text{IC}} = 370\)) и начальным значением \(Z_0 = 10\) к данным в таблице \ref{tab:3.7}. Чтобы проиллюстрировать расчёты, рассмотрим наблюдение номер 1. Первая статистика для нанесения на карту рассчитывается следующим образом:

\[
Z_1 = \lambda X_1 + (1 - \lambda)Z_0 = (0.1)(9.45) + (1 - 0.1)(10) = 9.945.
\]

\begin{table}[htbp]
	\centering
	\caption{Построение параметрической контрольной карты EWMA: данные, статистики для нанесения на карту и контрольные границы.}
	\label{tab:3.7}
	\renewcommand{\arraystretch}{1.1}
	\begin{tabular}{c r r r r}
		\toprule
		Номер & \multicolumn{1}{c}{(a)} & \multicolumn{1}{c}{(b)} & \multicolumn{1}{c}{(c)} & \multicolumn{1}{c}{(d)} \\
		наблюдения, & & & & \\
		\(i\) & \(X_i\) & \(Z_i\) & \(LCL_i\) & \(UCL_i\) \\
		\midrule
		&& 10 &&\\
		1  & 9.45  & 9.945  & 9.730 & 10.270 \\
		2  & 7.99  & 9.750  & 9.637 & 10.363 \\
		3  & 9.29  & 9.704  & 9.576 & 10.424 \\
		4  & 11.66 & 9.899  & 9.533 & 10.467 \\
		5  & 12.16 & 10.125 & 9.500 & 10.500 \\
		6  & 10.18 & 10.131 & 9.475 & 10.525 \\
		7  & 8.04  & 9.922  & 9.456 & 10.544 \\
		8  & 11.46 & 10.076 & 9.441 & 10.559 \\
		9  & 9.20  & 9.988  & 9.429 & 10.571 \\
		10 & 10.34 & 10.023 & 9.419 & 10.581 \\
		11 & 9.03  & 9.924  & 9.412 & 10.588 \\
		12 & 11.47 & 10.078 & 9.406 & 10.594 \\
		13 & 10.51 & 10.122 & 9.401 & 10.599 \\
		14 & 9.40  & 10.049 & 9.397 & 10.603 \\
		15 & 10.08 & 10.053 & 9.394 & 10.606 \\
		16 & 9.37  & 9.984  & 9.391 & 10.609 \\
		17 & 10.62 & 10.048 & 9.389 & 10.611 \\
		18 & 10.31 & 10.074 & 9.388 & 10.612 \\
		19 & 8.52  & 9.919  & 9.386 & 10.614 \\
		20 & 10.84 & 10.011 & 9.385 & 10.615 \\
		21 & 10.90 & 10.100 & 9.384 & 10.616 \\
		22 & 9.33  & 10.023 & 9.384 & 10.616 \\
		23 & 12.29 & 10.249 & 9.383 & 10.617 \\
		24 & 11.50 & 10.375 & 9.383 & 10.617 \\
		25 & 10.60 & 10.397 & 9.382 & 10.618 \\
		26 & 11.08 & 10.465 & 9.382 & 10.618 \\
		27 & 10.38 & 10.457 & 9.382 & 10.618 \\
		28 & 11.62 & 10.573 & 9.381 & 10.619 \\
		29 & 11.31 & 10.647 & 9.381 & 10.619 \\
		30 & 10.52 & 10.634 & 9.381 & 10.619 \\
		\bottomrule
	\end{tabular}
\end{table}

Подставляя \(\mu_0 = 10\), \(\sigma_0 = 1\), \(\lambda = 0.1\), \(L = 2.703\) и \(i = 1\) в уравнение \eqref{eq:3.11}, получаем \(CL\) и точные контрольные границы для данных (для \(i = 1\)) из таблицы \ref{tab:3.7}:

\[
\begin{aligned}
	UCL &= 10 + 2.703(1)\sqrt{\left(\frac{0.1}{2 - 0.1}\right)(1 - (1 - 0.1)^{2(1)})} = 10.270 \\
	CL &= 10 \\
	LCL &= 10 - 2.703(1)\sqrt{\left(\frac{0.1}{2 - 0.1}\right)(1 - (1 - 0.1)^{2(1)})} = 9.730.
\end{aligned}
\]

Контрольные границы стационарного состояния могут быть рассчитаны с использованием уравнения \eqref{eq:3.12} и задаются как \(LCL = 9.380\) и \(UCL = 10.620\) соответственно.

На рисунке \ref{fig-3.7} показаны статистики EWMA \(Z_i\), приведённые в столбце (b) таблицы \ref{tab:3.7}, которые нанесены на параметрическую контрольную карту EWMA вместе с зависящими от времени контрольными границами. Видно, что карта EWMA подаёт сигнал на периоде 29, что аналогично карте CUSUM, рассмотренной в примере 2.4.

В этом примере мы проиллюстрировали параметрическую карту EWMA для отдельных измерений. Однако если берутся рациональные подгруппы объёма \(n > 1\), то просто замените \(X_i\) на \(\overline{X}_i\) и \(\sigma_0\) на \(\sigma_{\overline{X}} = \frac{\sigma_0}{\sqrt{n}}\) в предыдущих уравнениях.

\begin{figure}[!ht]
	\centering
	\includegraphics[width=0.7\linewidth]{"./images/fig-3.7.png"}
	\caption[]{Параметрическая контрольная карта EWMA}
	\label{fig-3.7}
\end{figure}

\section{Карты Шухарта, EWMA и CUSUM: когда что использовать}

Как мы отмечали, карты Шухарта наиболее популярны на практике благодаря своей простоте, лёгкости применения и тому, что эти универсальные карты достаточно эффективны для обнаружения умеренных и больших сдвигов. Хотя карты Шухарта являются наиболее широко известными и используемыми контрольными картами на практике благодаря своей простоте и глобальной эффективности, другие классы карт, такие как карты CUSUM и EWMA, полезны и иногда более естественно подходят для среды контроля процесса с учётом последовательного характера сбора данных. Эти карты, обычно основанные на кумулятивных суммах статистики для нанесения на карту, получаемых по мере накопления данных, известны как более эффективные для обнаружения малых и умеренных по величине сдвигов в процессе. Таким образом, контрольные карты CUSUM и EWMA были разработаны как альтернативы карте Шухарта для обнаружения малых, устойчивых сдвигов среднего. Естественно возникает вопрос: «Какая контрольная карта, CUSUM или EWMA, является наиболее эффективной и когда её следует использовать?» Большая часть доступной литературы говорит о том, что эффективность карт CUSUM и EWMA очень схожа. Однако у каждой из них есть свои преимущества и недостатки.

Карта EWMA имеет так называемую проблему инерции, которая может быть практической проблемой. Прежде чем объяснить понятие инерции, сначала уточним, что и карта EWMA, и карта CUSUM имеют проблему инерции. Карта CUSUM не имеет столь значительной проблемы инерции, поскольку использует сбросы (см. Yashchin, 1987, 1993). Более подробное объяснение следует после определения термина «инерция». Термин «инерция» относится к мере сопротивления карты подаче сигнала о конкретном сдвиге процесса; например, если статистика для нанесения на карту EWMA в момент возникновения восходящего сдвига оказалась ближе к \(LCL\), то время, необходимое для достижения \(UCL\), будет больше, чем если бы статистика EWMA была ближе к \(CL\). Было показано, что карты EWMA имеют более серьёзную проблему инерции, чем карты CUSUM (см., например, Woodall and Mahmoud, 2005). Это особенно характерно для случая, когда нас интересуют как восходящие, так и нисходящие сдвиги (т.е. двусторонние контрольные карты), поскольку EWMA реализуется с помощью одной статистики для нанесения на карту, в отличие от процедуры CUSUM, которая использует две отдельные (верхнюю и нижнюю) статистики для нанесения на карту. Поэтому, хотя карта EWMA проще в реализации на практике, её эффективность в «наихудшем случае» для состояния OOC хуже, чем у карты CUSUM. Хотя в литературе есть некоторые рекомендации о том, как преодолеть проблему инерции, такие как использование односторонней процедуры EWMA со сбросом (см. Spliid, 2010), использование карт EWMA в сочетании с границами Шухарта (см. Woodall and Mahmoud, 2005) или использование адаптивного подхода EWMA (AEWMA) (см. Capizzi and Masarotto, 2003), эти усовершенствования обычно не используются на практике. Woodall and Mahmoud (2005) предложили меру инерции, называемую сопротивлением сигналу (signal resistance), которая определяется как наибольшее стандартное отклонение от целевого значения, не приводящее к немедленному сигналу OOC. Настоятельно рекомендуется учитывать свойства инерции контрольных карт как важный фактор при выборе контрольной карты и рассчитывать сопротивление сигналу именно по этой причине. Рассчитав сопротивление сигналу для карт EWMA и CUSUM, Woodall and Mahmoud (2005) пришли к выводу, что карта EWMA имеет худшие инерционные свойства, чем карта CUSUM, в том смысле, что значения сопротивления сигналу могут быть значительно выше. Ещё одним преимуществом карты CUSUM является то, что момент сдвига можно точно определить с помощью величин \(N^{+}\) и \(N^{-}\), которые указывают количество последовательных периодов, в течение которых CUSUM \(C_{i}^{+}\) и \(C_{i}^{-}\) были ненулевыми. Например, если сигнальное событие произошло на номере выборки 29 при соответствующей величине \(N^{+} = 7\) на номере выборки 29, мы заключим, что процесс в последний раз был в состоянии IC на номере выборки \(29 - 7 = 22\), поэтому сдвиг, вероятно, произошёл между номерами выборок 22 и 23. Ещё одним преимуществом использования карты CUSUM является то, что можно получить оценку нового среднего процесса после сдвига; например, продолжая предыдущий пример, если среднее процесса в состоянии IC равно 10 (\(\mu_{0} = 10\)), опорное значение взято равным 0.5 (\(K = 0.5\)), а значение \(C_{29}^{+} = 5.28\), то новое среднее процесса оценивается с использованием \(\hat{\mu} = \mu_{0} + K + \frac{C_{29}^{+}}{N^{+}} = 10 + 0.5 + \frac{5.28}{7} = 11.25\). Таким образом, среднее процесса в состоянии IC, равное 10, сдвинулось вверх до значения 11.25 где-то между номерами выборок 22 и 23. Обладая всеми этими преимуществами карты CUSUM, почему же тогда использовать карту EWMA? Hawkins and Wu (2014), пожалуй, выразились лучше всего, и мы цитируем их выводы: "Одно ясно — если сдвиг происходит в начале процесса или вблизи него (в „начальном состоянии“), то EWMA является лучшим выбором, чем карта CUSUM. Независимо от того, на какой размер сдвига рассчитан процесс мониторинга или какой сдвиг фактически произошёл, EWMA всегда реагирует быстрее, чем CUSUM".

\section{Усовершенствования контрольных карт}

\subsection{Правила чувствительности}

Хотя контрольная карта Шухарта эффективна для обнаружения больших сдвигов процесса, было показано, что ей не хватает чувствительности для обнаружения малых сдвигов процесса. В работах Koutras, Bersimis и Maravelakis (2007) и Park и Seo (2012) представлены обзоры литературы по картам Шухарта с дополнительными сенсибилизирующими правилами, основанными на сериях и сканах, для повышения эффективности контрольных карт Шухарта в обнаружении малых сдвигов. Сенсибилизирующие правила — это сигнальные правила, предназначенные для обнаружения некоторого маловероятного и/или неслучайного паттерна статистик для нанесения на карту на контрольной карте. Мы начнём с обсуждения хорошо известных правил Western Electric, а затем обсудим другие дополнительные сенсибилизирующие правила, основанные на сериях и сканах.

Правила Western Electric используют предупреждающие границы, обычно устанавливаемые на расстоянии 1-сигма или 2-сигма от \(CL\), для определения того, находится ли процесс в состоянии IC или OOC. Правила Western Electric (Western Electric Company, 1956) перечислены ниже. Карта подаст сигнал, если выполнено любое из этих условий.

\begin{enumerate}
	\item Одна или более точек на или за пределами контрольных границ.
	\item Две из трёх последовательных точек выходят за пределы 2-сигма предупреждающих границ, но всё ещё находятся внутри контрольных границ.
	\item Четыре из пяти последовательных точек выше 1-сигма границ.
	\item Серия из восьми последовательных точек по одну сторону от \(CL\).
\end{enumerate}

Правила Western Electric проиллюстрированы на рисунках \ref{fig-3.8}–\ref{fig-3.11}. Может возникнуть вопрос, как эти правила устанавливаются. Например, почему в правиле Western Electric номер 4 мы рассматриваем серию из восьми последовательных точек по одну сторону от \(CL\), а не девять последовательных точек по одну сторону от \(CL\)? Все эти правила разработаны так, чтобы они имели приблизительно одинаковую вероятность ложной тревоги, которая также близка к номинальному значению, обычно принимаемому равным 0.0027. Для правила Western Electric номер 4, например, если предположить, что \(CL\) контрольной карты равна медиане распределения в состоянии IC, то вероятности того, что статистика для нанесения на карту отобразится выше и ниже \(CL\), равны, таким образом, \(P(\text{Статистика для нанесения на карту отображается выше } CL) = P(\text{Статистика для нанесения на карту отображается ниже } CL) = 0.5\). Следовательно,

\[
\begin{aligned}
	&P(\text{8 точек отображаются выше } CL) \\
	&= P(\text{1-я точка выше})P(\text{2-я точка выше}) \cdots P(\text{8-я точка выше}) \\
	&= (0.5)(0.5) \cdots (0.5) \\
	&= (0.5)^8 \\
	&= 0.0039
\end{aligned}
\]

что близко к типичному желаемому номинальному значению 0.0027.

Правило Western Electric номер 1 проиллюстрировано на рисунке \ref{fig-3.8}. Это правило подаёт сигнал, если одна или более точек отображаются на или за пределами \(LCL\) или \(UCL\). Контрольная карта на рисунке \ref{fig-3.8} подаёт сигнал на номере выборки 9. Хотя это правило проиллюстрировано для возможного восходящего сдвига процесса, оно также подаёт сигнал, если бы статистики для нанесения на карту имели точную зеркальную копию относительно \(CL\), указывая на возможный нисходящий сдвиг процесса.

Правило Western Electric номер 2 проиллюстрировано на рисунке \ref{fig-3.9}. Это правило подаёт сигнал, если две из трёх последовательных точек выходят за пределы 2-сигма предупреждающих границ, но всё ещё находятся внутри контрольных границ. Контрольная карта на рисунке \ref{fig-3.9} подаёт сигнал на номере выборки 9. Хотя это правило проиллюстрировано для возможного восходящего сдвига процесса, оно также подаёт сигнал, если бы статистики для нанесения на карту имели точную зеркальную копию относительно \(CL\), указывая на возможный нисходящий сдвиг процесса.

Правило Western Electric номер 3 проиллюстрировано на рисунке \ref{fig-3.10}. Это правило подаёт сигнал, если четыре из пяти последовательных точек выходят за пределы 1-сигма предупреждающих границ, но всё ещё находятся внутри контрольных границ. Контрольная карта на рисунке \ref{fig-3.10} подаёт сигнал на номере выборки 10. Хотя это правило проиллюстрировано для возможного восходящего сдвига процесса, оно также подаёт сигнал, если бы статистики для нанесения на карту имели точную зеркальную копию относительно \(CL\), указывая на возможный нисходящий сдвиг процесса.

Правило Western Electric номер 4 проиллюстрировано на рисунке \ref{fig-3.11}. Это правило подаёт сигнал, если возникает серия из восьми последовательных точек по одну сторону от \(CL\). Контрольная карта на рисунке \ref{fig-3.11} подаёт сигнал на номере выборки 8. Хотя это правило проиллюстрировано для возможного восходящего сдвига процесса, оно также подаёт сигнал, если бы статистики для нанесения на карту имели точную зеркальную копию относительно \(CL\), указывая на возможный нисходящий сдвиг процесса.

\begin{figure}[!ht]
	\centering
	\includegraphics[width=0.7\linewidth]{"./images/fig-3.8.png"}
	\caption[]{Правило Western Electric номер 1}
	\label{fig-3.8}
\end{figure}

\begin{figure}[!ht]
	\centering
	\includegraphics[width=0.7\linewidth]{"./images/fig-3.9.png"}
	\caption[]{Правило Western Electric номер 2}
	\label{fig-3.9}
\end{figure}

\begin{figure}[!ht]
	\centering
	\includegraphics[width=0.7\linewidth]{"./images/fig-3.10.png"}
	\caption[]{Правило Western Electric номер 3}
	\label{fig-3.10}
\end{figure}

\begin{figure}[!ht]
	\centering
	\includegraphics[width=0.7\linewidth]{"./images/fig-3.11.png"}
	\caption[]{Правило Western Electric номер 4}
	\label{fig-3.11}
\end{figure}

В литературе были предложены некоторые другие сенсибилизирующие правила, к которым относятся:

\begin{itemize}
	\item Шесть точек подряд, монотонно возрастающих или убывающих.
	\item Пятнадцать точек подряд в зоне C (как выше, так и ниже \(CL\)); см. рисунки \ref{fig-3.8}–\ref{fig-3.11} для иллюстрации различных зон.
	\item Четырнадцать точек подряд, попеременно возрастающих и убывающих.
	\item Восемь точек подряд по обе стороны от \(CL\), ни одна из которых не находится в зоне C; см. рисунки \ref{fig-3.8}–\ref{fig-3.11} для иллюстрации различных зон.
	\item Необычный или неслучайный паттерн в данных.
	\item Одна или более точек вблизи предупреждающей или контрольной границы.
\end{itemize}

В таблице \ref{tab:3.8} перечислены некоторые возможные причины сигналов контрольной карты.

\begin{table}[htbp]
	\centering
	\caption{Возможные причины сигналов контрольной карты.}
	\label{tab:3.8}
	\renewcommand{\arraystretch}{1.2}
	\begin{tabular}{p{6.5cm}p{7cm}}
		\toprule
		Сигнал & Возможная причина \\
		\midrule
		Одна точка за пределами контрольных границ, за которой не следуют аномальные паттерны & Ошибка измерения или записи \\
		Пятнадцать точек подряд в зоне C (как выше, так и ниже \(CL\)); см. рисунки \ref{fig-3.8}–\ref{fig-3.11} для иллюстрации различных зон & Изолированное/временное событие, такое как замена персонала \\
		& Данные были сфальсифицированы, т.е. экстремальные значения не регистрируются или регистрируются неправильно \\
		Тренды, например, шесть точек подряд, монотонно возрастающих или убывающих & Постепенный дрейф данных может указывать на износ измерительного оборудования или механизмов. Это также может быть результатом прогрева и охлаждения машины или недостаточного технического обслуживания \\
		Циклический паттерн; см. рисунок \ref{fig-2.3} & Это может быть результатом изменений окружающей среды, таких как изменения температуры \\
		& Это также может быть результатом усталости оператора, колебаний напряжения или давления, или какой-либо другой переменной в оборудовании \\
		& Химические свойства сырья также могут играть роль \\
		& Стратификация может быть результатом неправильного расчёта контрольных границ или неправильного разделения на подгруппы \\
		& Это также может быть результатом не пересчёта контрольных границ после улучшения процесса \\
		Внезапный скачок & Это может быть результатом повреждения оборудования, появления нового персонала и т.д. \\
		\bottomrule
	\end{tabular}
\end{table}

\subsection{Правила сигнализации на основе серий}

Правила сигнализации на основе серий используются для повышения чувствительности карты Шухарта к обнаружению меньших сдвигов процесса при сохранении простоты карт Шухарта. Их также называют дополнительными правилами, поскольку они используются для дополнения исходной карты, чтобы сделать её более чувствительной к обнаружению изменений процесса. Исходная карта Шухарта использует статистику для нанесения на карту из текущей выборки, чтобы проверить, есть ли сигнал или нет, для определения состояния процесса; таким образом, эта карта также известна как карта 1-из-1 (основанная на правиле 1-из-1). Хотя это разумно, может быть полезно рассмотреть статистики для нанесения на карту из нескольких предыдущих выборок, чтобы увидеть, есть ли паттерн в сигналах. Например, если у нас есть сигналы из двух последовательных выборок (серия из двух), возникнет более сильное свидетельство изменения. Карты, использующие такие сигнальные правила, называются картами, улучшенными правилами серий. Двумя популярными сигнальными правилами являются правила серий 2-из-2 и 2-из-3, и эти карты, улучшенные правилами серий, обозначаются как карты 2-из-2 и 2-из-3.

Отметим, что правило серий \(k\)-из-\(k\) (\(k \geq 2\)) является общим случаем правила серий 2-из-2. Правило серий \(k\)-из-\(k\) подаёт сигнал, когда \(k\) последовательных статистик для нанесения на карту все отображаются на или за пределами контрольной границы(границ). Обобщением правила серий \(k\)-из-\(k\) является правило серий \(k\)-из-\(w\) (\(1 \leq k \leq w\)), которое подаёт сигнал, когда \(k\) из последних \(w\) статистик для нанесения на карту отображаются на или за пределами контрольной границы(границ).

Хотя карты, улучшенные правилами серий, могут сделать карту Шухарта более чувствительной к изменениям, могут быть некоторые потенциальные издержки. FAR может увеличиться, и результирующая карта может потерять способность немедленно обнаруживать большой сдвиг процесса. Правила сигнализации 2-из-3, например, требуют, чтобы сигнализировали как минимум последние две или три статистики для нанесения на карту, что может быть потенциально затратным для практика, поскольку способность карты обнаруживать большой сдвиг в процессе задерживается до тех пор, пока не будут собраны как минимум три выборки. Таким образом, издержки и выгоды должны быть тщательно взвешены, прежде чем принимать решение об использовании дополнительной карты, улучшенной правилами серий.

В литературе было предложено множество параметрических контрольных карт с правилами сигнализации на основе серий. Заинтересованный читатель может ознакомиться, например, с работами Roberts (1958), Champ and Woodall (1987), Champ (1992), Klein (2000), Shmueli and Cohen (2003), Acosta-Mejia (2007), Khoo and Ariffin (2006), Lim and Cho (2009), Cheng and Chen (2011), Santiago and Smith (2013) и других. Однако предмет нашего обсуждения — непараметрические контрольные карты, и соответственно мы сосредоточимся на непараметрических контрольных картах для мониторинга положения и/или масштаба с правилами сигнализации на основе серий. В целом, чтобы узнать о преимуществах и недостатках контрольных карт типа Шухарта, дополненных правилами сигнализации на основе серий, обратитесь к Nelson (1985). Koutras et al. (2007) показывают, что общая вероятность ложной тревоги при использовании \(h\) решающих правил равна

\begin{equation}
	\alpha = \prod_{i=1}^{h}(1 - \alpha_{i}),
	\label{eq:3.16}
\end{equation}

при условии, что все \(h\) решающих правил независимы, причём правило \(i\) имеет вероятность \(\alpha_{i}\) того, что статистика для нанесения на карту отобразится на или за пределами контрольных границ, когда процесс находится в состоянии IC. Отметим, что уравнение \eqref{eq:3.16} не справедливо для типичных правил серий; см. Nelson (1985) и Montgomery (2009).

\subsubsection{Сигнальные индикаторы}

Пусть статистика для нанесения на карту для \(i\)-й подгруппы обозначается как \(Q_{i}\) для \(i = 1, 2, 3, \ldots\). Пусть индикаторные функции для односторонних контрольных карт обозначаются как

\begin{equation}
	\xi_{i}^{+} = I(Q_{i} \geq UCL) = \begin{cases}
		1 & \text{если } Q_{i} \geq UCL \\
		0 & \text{если } Q_{i} < UCL
	\end{cases}
	\label{eq:3.17}
\end{equation}

и

\begin{equation}
	\xi_{i}^{-} = I(Q_{i} \leq LCL) = \begin{cases}
		1 & \text{если } Q_{i} \leq LCL \\
		0 & \text{если } Q_{i} > LCL
	\end{cases}
	\label{eq:3.18}
\end{equation}

соответственно. Таким образом, \(\xi_{i}^{+}\) равен 1 и указывает на сигнал для верхней односторонней карты, когда \(Q_{i}\) отображается на или выше UCL. Аналогично, \(\xi_{i}^{-}\) принимает значение 1 и указывает на сигнал для нижней односторонней карты, когда \(Q_{i}\) отображается на или ниже \(LCL\). Таким образом, эти бинарные переменные называются сигнальными индикаторами. Для двусторонней контрольной карты сигнальный индикатор может быть обозначен как

\begin{equation}
	\xi_{i} = \begin{cases}
		1 & \text{если } Q_{i} \geq UCL \\
		0 & \text{если } LCL < Q_{i} < UCL \\
		2 & \text{если } Q_{i} \leq LCL
	\end{cases}
	\label{eq:3.19}
\end{equation}

так что \(\xi_{i}\) указывает, отображается ли статистика для нанесения на карту \(Q_{i}\) на или выше \(UCL\) (\(\xi_{i} = 1\)), между двумя границами (\(\xi_{i} = 0\)) или на или ниже \(LCL\) (\(\xi_{i} = 2\)). Значения 1 и 2 указывают на наличие сигнала, соответственно на верхней или нижней стороне, тогда как значение 0 указывает на отсутствие сигнала и на то, что процесс находится в состоянии IC.

\paragraph{Правило сигнализации на основе серий 1-из-1}
Наиболее часто используемыми и наименее сложными картами являются обычные карты Шухарта, называемые картами 1-из-1, которые не дополнены никакими правилами серий. Однако известно, что они нечувствительны к обнаружению малых сдвигов процесса. Эти карты подают сигнал, когда происходит событие \(A_{1}\), или \(A_{2}\), или \((A_{1} \cup A_{2})\) соответственно, где

\begin{enumerate}
	\item Правило сигнализации на основе серий 1-из-1 для верхней односторонней карты:
	\[
	A_{1}: \{Q_{i} \geq UCL\} \leftrightarrow \{\xi_{i}^{+} = 1\}.
	\]
	\item Правило сигнализации на основе серий 1-из-1 для нижней односторонней карты:
	\[
	A_{2}: \{Q_{i} \leq LCL\} \leftrightarrow \{\xi_{i}^{-} = 1\}.
	\]
	\item Правило сигнализации на основе серий 1-из-1 для двусторонней карты:
	\[
	\begin{aligned}
		A &: \{Q_{i} \geq UCL \cup Q_{i} \leq LCL\} \leftrightarrow \{A_{1} \cup A_{2}\} \\
		&\leftrightarrow \{\xi_{i}^{+} = 1 \cup \xi_{i}^{-} = 1\} \leftrightarrow \{\xi_{i} = 1 \cup \xi_{i} = 2\}.
	\end{aligned}
	\]
\end{enumerate}

На рисунке \ref{fig-3.12} проиллюстрированы пункты (i) и (ii) правила сигнализации на основе серий 1-из-1 для верхней односторонней и нижней односторонней карты соответственно. Панель (a) обнаруживает восходящий сдвиг в момент времени \(i = 7\), когда \(Q_{7}\) отображается выше \(UCL\), тогда как панель (b) обнаруживает нисходящий сдвиг в момент времени \(i = 7\), когда \(Q_{7}\) отображается ниже \(LCL\). Для каждой карты процесс объявляется находящимся в состоянии OOC, и может быть начат поиск назначаемых причин.

\begin{figure}[!ht]
	\centering
	\includegraphics[width=0.7\linewidth]{"./images/fig-3.12.png"}
	\caption[]{Правило сигнализации на основе серий 1-из-1 для верхней односторонней и нижней односторонней карты}
	\label{fig-3.12}
\end{figure}

На рисунке \ref{fig-3.13} проиллюстрирован пункт (3) правила сигнализации на основе серий 1-из-1 для двусторонней карты. Обе карты подают сигнал в момент времени \(i = 7\); сначала обнаруживается восходящий сдвиг, когда \(Q_{7}\) отображается выше UCL, а затем обнаруживается нисходящий сдвиг, когда \(Q_{7}\) отображается ниже LCL.

\begin{figure}[!ht]
	\centering
	\includegraphics[width=0.7\linewidth]{"./images/fig-3.13.png"}
	\caption[]{Правило сигнализации на основе серий 1-из-1 для двусторонней карты}
	\label{fig-3.13}
\end{figure}

\paragraph{Правила сигнализации на основе серий \(k\)-из-\(k\) и \(k\)-из-\(w\)}
Несколько правил сигнализации на основе серий, таких как правила серий 2-из-2 и 2-из-3, были исследованы многими авторами, и было доказано, что их карты, «улучшенные правилами серий», превосходят карту 1-из-1.

\subparagraph{Односторонние правила сигнализации на основе серий \(k\)-из-\(k\) и \(k\)-из-\(w\)}
Для иллюстрации мы начнём с правила сигнализации на основе серий 2-из-2. Это правило использует последние две статистики для нанесения на карту, \(Q_{i-1}\) и \(Q_{i}\), чтобы определить, находится ли процесс в состоянии IC или OOC. Для правила сигнализации на основе серий 2-из-2 верхняя односторонняя карта подаёт сигнал, когда происходит событие \(B_{1}\), а нижняя односторонняя карта подаёт сигнал, когда происходит событие \(B_{2}\), где

\begin{enumerate}
	\item верхняя односторонняя карта, дополненная правилом сигнализации на основе серий 2-из-2:
	\[
	B_{1}: \{\xi_{i-1}^{+} = \xi_{i}^{+} = 1\}
	\]
	\item нижняя односторонняя карта, дополненная правилом сигнализации на основе серий 2-из-2:
	\[
	B_{2}: \{\xi_{i-1}^{-} = \xi_{i}^{-} = 1\}.
	\]
\end{enumerate}

Панели (a) и (b) на рисунке \ref{fig-3.14} иллюстрируют эти сигнальные события \(B_{1}\) и \(B_{2}\), и обе карты подают сигнал в момент времени \(i = 7\), когда две последовательные точки отображаются выше (ниже) \(UCL\) (\(LCL\)). Для каждой карты процесс объявляется находящимся в состоянии OOC, и может быть начат поиск назначаемых причин.

\begin{figure}[!ht]
	\centering
	\includegraphics[width=0.7\linewidth]{"./images/fig-3.14.png"}
	\caption[]{Правило сигнализации на основе серий 2-из-2 для верхней односторонней и нижней односторонней карты}
	\label{fig-3.14}
\end{figure}

Далее мы проиллюстрируем правило сигнализации на основе серий 2-из-3 для верхней односторонней карты и нижней односторонней карты соответственно. Эти карты используют две из последних трёх статистик для нанесения на карту \(Q_{i-2}\), \(Q_{i-1}\) и \(Q_{i}\), чтобы определить, находится ли процесс в состоянии IC или OOC. Хотя существует три способа, которыми две из трёх статистик для нанесения на карту могут привести к сигналу, только в двух из этих возможностей последняя статистика для нанесения на карту отображается на или выше (ниже) верхней (нижней) контрольной границы. Односторонние карты, дополненные правилом сигнализации на основе серий 2-из-3, подают сигнал, когда происходит событие \(C(+) = \{C_1 \cup C_2\}\) или \(C(-) = \{C_3 \cup C_4\}\), где:

\begin{table}[htbp]
	\centering
	\renewcommand{\arraystretch}{1.2}
	\begin{tabular}{p{13.5cm}}
		\toprule
		Верхняя односторонняя карта, дополненная правилом сигнализации на основе серий 2-из-3 \\
		\midrule
		(1) \(C_1: \{Q_{i-2} < UCL \text{ и } Q_{i-1} \geq UCL \text{ и } Q_{i} \geq UCL\} \leftrightarrow \{\xi_{i-2}^{+} = 0, \xi_{i-1}^{+} = 1, \xi_{i}^{+} = 1\}\) \\
		(2) \(C_2: \{Q_{i-2} \geq UCL \text{ и } Q_{i-1} < UCL \text{ и } Q_{i} \geq UCL\} \leftrightarrow \{\xi_{i-2}^{+} = 1, \xi_{i-1}^{+} = 0, \xi_{i}^{+} = 1\}\) \\
		\midrule
		Нижняя односторонняя карта, дополненная правилом сигнализации на основе серий 2-из-3 \\
		\midrule
		(3) \(C_3: \{Q_{i-2} > LCL \text{ и } Q_{i-1} \leq LCL \text{ и } Q_{i} \leq LCL\} \leftrightarrow \{\xi_{i-2}^{-} = 0, \xi_{i-1}^{-} = 1, \xi_{i}^{-} = 1\}\) \\
		(4) \(C_4: \{Q_{i-2} \leq LCL \text{ и } Q_{i-1} > LCL \text{ и } Q_{i} \leq LCL\} \leftrightarrow \{\xi_{i-2}^{-} = 1, \xi_{i-1}^{-} = 0, \xi_{i}^{-} = 1\}\) \\
		\bottomrule
	\end{tabular}
\end{table}

На рисунке \ref{fig-3.15} проиллюстрированы эти сигнальные события \(C_1\), \(C_2\), \(C_3\) и \(C_4\) для верхней односторонней карты и нижней односторонней карты, дополненных правилом сигнализации на основе серий 2-из-3, соответственно. Верхние односторонние карты на панели (a) подают сигнал в момент времени \(i = 7\), когда две из последних трёх точек отображаются выше UCL. Для каждой карты процесс объявляется находящимся в состоянии OOC, и может быть начат поиск назначаемых причин. Аналогично, нижние односторонние карты на панели (b) подают сигнал в момент времени \(i = 7\), когда две из последних трёх точек отображаются ниже LCL. Для каждой карты процесс объявляется находящимся в состоянии OOC, и может быть начат поиск назначаемых причин.

\begin{figure}[!ht]
	\centering
	\includegraphics[width=0.7\linewidth]{"./images/fig-3.15.png"}
	\caption[]{Правило сигнализации на основе серий 2-из-3 для верхней односторонней карты и нижней односторонней карты}
	\label{fig-3.15}
\end{figure}

Некоторые из паттернов правил сигнализации на основе серий 2-из-3 исключаются. Если первая и вторая статистики для нанесения на карту выходят за пределы контрольных границ, но третья (или последняя) статистика для нанесения на карту отображается между контрольными границами, процесс не может быть объявлен находящимся в состоянии OOC, поскольку процесс вернулся в состояние IC, и соответственно эти паттерны исключаются из правил сигнализации на основе серий 2-из-3. См. рисунок \ref{fig-3.16}, где верхняя односторонняя карта, дополненная правилом сигнализации на основе серий 2-из-3, и нижняя односторонняя карта, дополненная правилом сигнализации на основе серий 2-из-3, исключаются как сигнальные события.

\begin{figure}[!ht]
	\centering
	\includegraphics[width=0.7\linewidth]{"./images/fig-3.16.png"}
	\caption[]{Правила сигнализации на основе серий 2-из-3, которые исключаются для верхней односторонней и нижней односторонней карты}
	\label{fig-3.16}
\end{figure}

\subparagraph{Двусторонние карты, дополненные правилами сигнализации на основе серий \(k\)-из-\(k\) и \(k\)-из-\(w\)}
Двусторонняя карта имеет как верхнюю, так и нижнюю контрольную границу, таким образом, она может обнаруживать либо восходящий, либо нисходящий сдвиг. Для иллюстрации мы начнём с правила сигнализации на основе серий 2-из-2. Двусторонняя карта, дополненная правилом сигнализации на основе серий 2-из-2, во многом схожа с односторонними картами в том смысле, что она использует последние две статистики для нанесения на карту, \(Q_{i-1}\) и \(Q_{i}\), чтобы определить, находится ли процесс в состоянии IC или OOC, но она способна обнаруживать как восходящий, так и нисходящий сдвиг. Двусторонняя карта, дополненная правилом сигнализации на основе серий 2-из-2, подаёт сигнал, когда происходит событие \(D_{i}\), \(i = 1, 2, 3, 4\), \(\{D_{1}, D_{2}, D_{3} \text{ или } D_{4}\}\) где

\begin{enumerate}
	\item обе статистики для нанесения на карту отображаются на или выше UCL:
	\[
	D_{1}: \{Q_{i-1} > UCL \text{ и } Q_{i} > UCL\} \leftrightarrow \{\xi_{i-1} = 1, \xi_{i} = 1\}
	\]
	\item обе статистики для нанесения на карту отображаются на или ниже \(LCL\):
	\[
	D_{2}: \{Q_{i-1} \leq LCL \text{ и } Q_{i} \leq LCL\} \leftrightarrow \{\xi_{i-1} = 2, \xi_{i} = 2\}
	\]
	\item первая статистика для нанесения на карту отображается на или выше \(UCL\), а вторая отображается на или ниже \(LCL\):
	\[
	D_{3}: \{Q_{i-1} \geq UCL \text{ и } Q_{i} \leq LCL\} \leftrightarrow \{\xi_{i-1} = 1, \xi_{i} = 2\}
	\]
	\item первая статистика для нанесения на карту отображается на или ниже \(LCL\), а вторая отображается на или выше \(UCL\):
	\[
	D_{4}: \{Q_{i-1} \leq LCL \text{ и } Q_{i} \geq UCL\} \leftrightarrow \{\xi_{i-1} = 2, \xi_{i} = 1\}.
	\]
\end{enumerate}

На рисунке \ref{fig-3.17} проиллюстрированы эти сигнальные события \(D_{1}\), \(D_{2}\), \(D_{3}\) и \(D_{4}\) для двусторонней карты, дополненной правилами сигнализации на основе серий 2-из-2, соответственно. Панели (a) и (b) иллюстрируют события \(D_{1}\) и \(D_{2}\); обе карты подают сигнал в момент времени \(i = 7\), когда две последовательные точки отображаются выше (ниже) UCL (LCL). Панели (c) и (d) иллюстрируют события \(D_{3}\) и \(D_{4}\), когда восходящий сдвиг немедленно сменяется нисходящим сдвигом, или наоборот.

\begin{figure}[!ht]
	\centering
	\includegraphics[width=0.7\linewidth]{"./images/fig-3.17.png"}
	\caption[]{Правило сигнализации на основе серий 2-из-2 для двусторонней карты}
	\label{fig-3.17}
\end{figure}

Далее мы проиллюстрируем двустороннюю карту с правилами сигнализации на основе серий 2-из-3. Подобно односторонней карте, двусторонняя карта использует две из последних трёх статистик для нанесения на карту, \(Q_{i-2}\), \(Q_{i-1}\) и \(Q_{i}\), чтобы определить, находится ли процесс в состоянии IC или OOC. Существует до 12 сценариев, в которых ровно две из трёх статистик для нанесения на карту могут выходить за пределы контрольных границ; однако мы сосредоточимся только на четырёх сигнальных событиях \(E_{i}\), \(i = 1, 2, 3, 4\), \(\{E_{1}, E_{2}, E_{3} \text{ или } E_{4}\}\), которые могут произойти, где:

\begin{table}[htbp]
	\centering
	\renewcommand{\arraystretch}{1.2}
	\begin{tabular}{p{13.5cm}}
		\toprule
		Две из трёх статистик для нанесения на карту отображаются на или выше UCL\\
		\midrule
		(1) \(E_1: \{LCL < Q_{i-2} < UCL \text{ и } Q_{i-1} \geq UCL \text{ и } Q_{i} \geq UCL\} \leftrightarrow \{\xi_{i-2} = 0, \xi_{i-1} = 1, \xi_{i} = 1\}\)\\
		(2) \(E_2: \{Q_{i-2} \geq UCL \text{ и } LCL < Q_{i-1} < UCL \text{ и } Q_{i} \geq UCL\} \leftrightarrow \{\xi_{i-2} = 1, \xi_{i-1} = 0, \xi_{i} = 1\}\)\\
		\midrule
		Две из трёх статистик для нанесения на карту отображаются на или ниже LCL\\
		\midrule
		(3) \(E_3: \{LCL < Q_{i-2} < UCL \text{ и } Q_{i-1} \leq LCL \text{ и } Q_{i} \leq LCL\} \leftrightarrow \{\xi_{i-2} = 0, \xi_{i-1} = 2, \xi_{i} = 2\}\)\\
		(4) \(E_4: \{Q_{i-2} \leq LCL \text{ и } LCL < Q_{i-1} < UCL \text{ и } Q_{i} \leq LCL\} \leftrightarrow \{\xi_{i-2} = 2, \xi_{i-1} = 0, \xi_{i} = 2\}\)\\
		\bottomrule
		\end{tabular}
	\end{table}

На рисунке \ref{fig-3.18} проиллюстрированы эти сигнальные события \(E_{1}\), \(E_{2}\), \(E_{3}\) и \(E_{4}\) для двусторонней карты, дополненной правилами сигнализации на основе серий 2-из-3. Панель (a) иллюстрирует события \(E_{1}\) и \(E_{2}\), где ровно две из трёх статистик для нанесения на карту отображаются на или выше UCL, тогда как панель (b) иллюстрирует события \(E_{3}\) и \(E_{4}\), где ровно две из трёх статистик для нанесения на карту отображаются на или ниже LCL.

\begin{figure}[!ht]
\centering
\includegraphics[width=0.7\linewidth]{"./images/fig-3.18.png"}
\caption[]{Правило сигнализации на основе серий 2-из-3 для двусторонней карты}
\label{fig-3.18}
\end{figure}

На рисунке \ref{fig-3.19} проиллюстрированы восемь из 12 сценариев, которые исключаются как сигнальные события. Панель (a) показывает возможный обнаруженный размах, тогда как панель (b) иллюстрирует нисходящий или восходящий тренд в процессе. Как упоминалось на рисунке \ref{fig-3.16}, четыре события, проиллюстрированные на панели (c), будут исключены как сигнальные события из-за того, что последняя точка отображается между контрольными границами; следовательно, процесс не может быть объявлен находящимся в состоянии OOC, поскольку процесс вернулся в состояние IC.

\begin{figure}[!ht]
\centering
\includegraphics[width=0.7\linewidth]{"./images/fig-3.19.png"}
\caption[]{Правила сигнализации на основе серий 2-из-3, которые исключаются для двусторонней карты}
\label{fig-3.19}
\end{figure}

\section{Распределение длины серии в случае заданных параметров (случай K)}

Эффективность контрольной карты анализируется через распределение длины серии и связанные с ним характеристики, такие как различные моменты и процентили. В этом разделе мы обсуждаем методы расчёта распределения длины серии контрольных карт. Обсуждение носит общий характер и применимо ко всем типам контрольных карт, включая параметрические и непараметрические контрольные карты, а также случаи известных (случай K) и неизвестных (случай U) параметров. Сначала мы рассмотрим случай K.

\subsection{Методы расчёта распределения длины серии}

Существует четыре метода расчёта (или, по крайней мере, аппроксимации) распределения длины серии контрольной карты. В частности, в случае K это:

\begin{enumerate}
	\item точный подход (для карт Шухарта и некоторых карт типа Шухарта);
	\item подход на основе цепей Маркова (MC);
	\item подход на основе интегральных уравнений;
	\item подход на основе компьютерного моделирования (Монте-Карло).
\end{enumerate}

Далее следует обсуждение каждого метода. Сначала мы рассмотрим случай карт Шухарта.

\subsubsection{Точный подход (для карт Шухарта и некоторых карт типа Шухарта)}

Для карт Шухарта и некоторых контрольных карт типа Шухарта можно легко рассчитать характеристики распределения длины серии точно в случае K. Для двусторонней 3-сигма карты Шухарта \(\overline{X}\) длина серии, которая представляет собой количество выборок до первого сигнала, следует \textbf{геометрическому распределению} (geometric 
distribution) с вероятностью \(P\left(\frac{\sqrt{n}|\overline{X} - \mu|}{\sigma} > z_{\alpha}\right)\), которая равна FAR, обозначаемой \(\alpha\), когда процесс находится в состоянии IC. Мы приводим более подробную информацию об этом позже. Таким образом, в случае K все свойства IC контрольной карты Шухарта \(\overline{X}\) могут быть охарактеризованы через FAR, \(\alpha\), и получены из свойств геометрического распределения. Так, например, \(ARL_{\text{IC}} = \frac{1}{\alpha}\), \(SDRL_{\text{IC}} = \frac{1 - \alpha}{\alpha^2}\) и так далее.

\begin{figure}[!ht]
	\centering
	\includegraphics[width=0.7\linewidth]{"./images/fig-3.20.png"}
	\caption[]{Распределение длины серии контрольной карты Шухарта для среднего в случае K при \(\alpha = 0.05\)}
	\label{fig-3.20}
\end{figure}

На рисунке \ref{fig-3.20} показано распределение длины серии для \(\alpha = 0.05\). Правосторонняя асимметрия распределения очевидна. Для меньших значений \(\alpha\), таких как 0.0027, что довольно типично для промышленности, распределение длины серии имеет ещё более выраженную правостороннюю асимметрию с ещё более длинным правым хвостом.

Для некоторых хорошо известных контрольных карт, таких как CUSUM, EWMA и карты Шухарта с правилами серий, более распространённым подходом к определению распределения длины серии было умелое применение идеи \textbf{цепи Маркова} (Markov chain -- MC). Теория MC чаще используется в области случайных процессов, и литература по контрольным картам заимствовала некоторые важные результаты из этой литературы. Ниже мы описываем некоторые концепции и ключевые результаты.

\subsubsection{Подход на основе цепей Маркова}

Некоторые теория и концепции в области так называемой конечной однородной MC были применены для вывода распределения длины серии и различных характеристик распределения длины серии некоторых карт с взвешиванием по времени, таких как CUSUM и EWMA. Читатель может обратиться, например, к Fu and Lou (2003) и Balakrishnan and Koutras (2002) для более подробного рассмотрения этих тем; однако большинство существующих книг по SPC, по-видимому, не рассматривают эту важную тему. В этом разделе приводятся некоторые основные термины, определения, результаты и теоремы, которые дают необходимую основу для расчёта распределения длины серии некоторых контрольных карт с помощью подхода MC. Эти теоремы и результаты имеют решающее значение для последующих глав; приводятся иллюстрации.

Подход MC предполагает, что статистика для нанесения на карту рассматривается как следующая MC (как случайный процесс), с пространством состояний \(S\) и матрицей переходных вероятностей \(M\). Пространство состояний состоит из двух типов состояний:

\begin{enumerate}
	\item одно поглощающее состояние (т.е. это состояние достигается, когда карта подаёт сигнал, т.е. когда статистика для нанесения на карту больше или равна \(UCL\), или меньше или равна \(LCL\) для двусторонней карты); и
	\item \(\nu\) переходных или невозвратных состояний, так что всего имеется \(\nu + 1\) состояний.
\end{enumerate}

Матрица переходных вероятностей размера \((\nu + 1) \times (\nu + 1)\), \(M\), часто записывается в блочной форме

\begin{equation}
	M_{(\nu+1)\times(\nu+1)} = 
	\left(
	\begin{array}{c|c}
		Q_{\nu\times\nu} & \underline{p}_{\nu\times1} \\
		\hline
		\underline{0}'_{1\times\nu} & 1_{1\times1}
	\end{array}
	\right)
	\label{eq:3.20}
\end{equation}

где подматрица \(Q_{\nu \times \nu}\) содержит все вероятности перехода из одного переходного состояния в другое и называется существенной матрицей переходных вероятностей. Вектор-столбец \(\underline{p}_{\nu \times 1}\) содержит все вероятности перехода из каждого переходного состояния в поглощающее состояние; \(\underline{0}_{1 \times \nu}^{\prime}\) — это вектор-строка нулей, состоящая из вероятностей перехода из поглощающего состояния в каждое переходное состояние (все они равны нулю), а скалярное значение 1 — это вероятность остаться в поглощающем состоянии после того, как оно было достигнуто. Отметим, что ключевым компонентом при использовании подхода MC является получение существенной матрицы переходных вероятностей \(Q_{\nu \times \nu}\).

В контексте MC случайная величина длины серии \(N\) контрольной карты рассматривается как время ожидания до первого входа MC в поглощающее состояние. Используя эту аналогию и предполагая, что процесс начинается в состоянии IC, когда карта внедряется, функция вероятности (pmf), ожидаемое значение (ARL), стандартное отклонение (SDRL) и функция распределения (cdf) \(N\) задаются как (см. Fu and Lou, 2003, глава 5)

\begin{equation}
	P(N = t) = \xi Q^{t-1}(I - Q)\underline{1} \quad \text{для } t = 1, 2, 3, \ldots
	\label{eq:3.21}
\end{equation}

\begin{equation}
	E(N) = \xi (I - Q)^{-1}\underline{1}
	\label{eq:3.22}
\end{equation}

\begin{equation}
	SD(N) = \sqrt{\xi(I + Q)(I - Q)^{-2}\underline{1} - (E(N))^{2}}
	\label{eq:3.23}
\end{equation}

\begin{equation}
	P(N \leq t) = 1 - \xi Q^{t}\underline{1} \quad \text{для } t = 1, 2, 3, \ldots
	\label{eq:3.24}
\end{equation}

соответственно, где \(I = I_{\nu \times \nu}\) — единичная матрица, \(Q = Q_{\nu \times \nu}\) — существенная матрица переходных вероятностей, \(\underline{1} = \underline{1}_{\nu \times 1}\) — вектор-столбец, все элементы которого равны единице, а \(\underline{\xi} = \underline{\xi}_{1 \times \nu}\) — вектор-строка, называемая начальным вектором вероятностей, который содержит вероятности того, что MC начнётся в заданном состоянии. Вектор \(\xi = (\xi_{-s}, \ldots, \xi_{s})\) обычно выбирается так, чтобы \(\sum_{i=-s}^{s} \xi_{i} = 1\).

Режимы анализа нулевого состояния и стационарного состояния используются для характеристики краткосрочных и долгосрочных свойств длины серии контрольной карты соответственно. То есть длина серии нулевого состояния — это количество статистик для нанесения на карту, нанесённых до подачи сигнала, когда процесс начинается в некотором начальном состоянии; однако длина серии стационарного состояния — это количество статистик для нанесения на карту, нанесённых до подачи сигнала, когда процесс начинается и остаётся в состоянии IC в течение очень долгого времени. Затем, в некоторый случайный момент времени, наблюдается сигнал OOC; см., например, Champ (1992) и Machado and Costa (2014a, b). Пусть \(q_{1 \times \nu}\) обозначает вектор-строку начальных вероятностей, связанных с режимом нулевого состояния, то есть

\begin{equation}
	\xi = q_{1 \times \nu} = (0 \ 0 \ \ldots \ 1 \ \ldots \ 0 \ 0)
	\label{eq:3.25}
\end{equation}

так что элемент начального состояния в матрице переходных вероятностей соответствует значению 1 в уравнении \eqref{eq:3.25}. Следует отметить, что 1 находится в середине вектора-строки, так что процесс начинается в состоянии IC. Это делается путём принудительной установки начального состояния как можно ближе к \(CL\), то есть выбора среднего значения, соответствующего среднему дискретизированному подинтервалу, который, в свою очередь, является наиболее близким к \(CL\) из возможных начальных значений. В стационарном состоянии вектор \(\xi = q_{1 \times \nu}\) заменяется вектором \(s_{1 \times \nu}\), то есть вектором начальных вероятностей стационарного состояния. Он задаётся как

\begin{equation}
	\xi = \xi_{1 \times \nu} = (s_{1} \ s_{2} \ \ldots \ s_{M})
	\label{eq:3.26}
\end{equation}

где элементы в уравнении \eqref{eq:3.26} в сумме равны единице. Например, для \(v = 5\) \(\xi\) может быть равно \(\xi_{1 \times 5} = (0.15 \ 0.2 \ 0.3 \ 0.2 \ 0.15)\). Это означало бы, что можно начать ближе к \(CL\) с наибольшей вероятностью 0.3, или начать дальше от \(CL\) с меньшей вероятностью 0.2, и, наконец, можно начать близко к контрольным границам с наименьшей вероятностью 0.15. Наименьшая вероятность назначается для начала близко к \(UCL\) и близко к \(LCL\) соответственно, поскольку это не идеально. Другой пример: для \(v = 7\) имеем \(\xi_{1 \times 7} = (0 \ 0 \ 0.1 \ 0.8 \ 0.1 \ 0 \ 0)\). В этом примере некоторые элементы \(\xi\) равны нулю, а другие ненулевые. Здесь это означает, что можно начать ближе к \(CL\) с наибольшей вероятностью 0.8, или начать дальше от \(CL\) с меньшей вероятностью 0.1, и, наконец, можно начать близко к контрольным границам с нулевой вероятностью.

Из уравнения \eqref{eq:3.24} ясно, что для применения подхода MC для нахождения распределения длины серии данной контрольной карты необходимо задать существенную матрицу переходных вероятностей \(Q\). Как только матрица \(Q\) найдена, распределение длины серии и связанные с ним характеристики могут быть легко получены с использованием приведённых выше формул для заданного начального вектора вероятностей.

Теперь, чтобы рассчитать элементы \(Q\), сначала необходимо определить пространство состояний. Для контрольной карты по альтернативным признакам с дискретной статистикой для нанесения на карту это легко, поскольку количество состояний будет фиксированным. Однако для контрольной карты по количественным признакам, которая использует статистику для нанесения на карту с непрерывным распределением, состояния процесса должны быть определены тщательно. Ниже приведены шаги, как это обычно делается.

Следует отметить, что при использовании правил сигнализации на основе серий количество состояний определяется типом используемого правила, например, используется ли правило серий 2-из-2 или 2-из-3. Нет необходимости дискретизировать интервал между \(LCL\) и \(UCL\) на \(\nu = 2s + 1\) подинтервалов (как объясняется на шаге 1 ниже), поскольку, как указывалось ранее, для правил сигнализации на основе серий количество состояний зависит от типа правила. Также следует отметить, что для карт CUSUM количество состояний является фиксированным числом (см. пример 3.6). Нет необходимости дискретизировать интервал между \(LCL\) и \(UCL\) на \(\nu = 2s + 1\) подинтервалов (как объясняется на шаге 1 ниже), поскольку, как указывалось ранее, для карт CUSUM количество состояний будет фиксированным числом. Приведённые ниже шаги применимы к картам EWMA.

\noindent\textbf{Шаг 1:} Разделите интервал между \(LCL\) и \(UCL\) на \(\nu = 2s + 1\) подинтервалов шириной \(2\tau = (UCL - LCL)/\nu\) (см. рисунок \ref{fig-3.21}), где каждый подинтервал соответствует переходному (или достижимому) состоянию MC. Говорят, что состояние \(j\) достижимо из состояния \(i\) тогда и только тогда, когда, начиная из состояния \(i\), возможно, что процесс на некотором этапе войдёт в состояние \(j\). Отметим, что \(2\tau = (UCL - LCL)/\nu\) будет равно \(2\tau = 2UCL/\nu\), так что \(\tau = UCL/\nu\), если контрольные границы расположены симметрично относительно нуля (т.е. \(-LCL = UCL\)). Следует отметить, что взятие больших значений \(\nu\) приведёт к более точным результатам. Подробное обсуждение выбора \(\nu\) приведено в замечании 3.2.

\noindent\textbf{Шаг 2:} Выберите количество дискретизированных подинтервалов \(\nu\) в виде нечётного положительного целого числа, так чтобы существовал единственный средний элемент.

\noindent\textbf{Шаг 3:} Объявите, что статистика для нанесения на карту \(Z_{i}\) находится в переходном состоянии \(j\) в момент времени \(i\), если \(S_{j} - \tau < Z_{i} \leq S_{j} + \tau\) для \(j = -s, -s+1, \ldots, s-1\) и \(S_{j} - \tau < Z_{i} < S_{j} + \tau\) для \(j = s\), где \(S_{j}\) обозначает середину \(j\)-го интервала.

\noindent\textbf{Шаг 4:} Рассчитайте одношаговые переходные вероятности (\(p_{ij}\)), где \(p_{ij}\) обозначает вероятность перехода из состояния \(i\) в состояние \(j\) за один шаг в любой момент времени.

\noindent\textbf{Шаг 5:} Постройте матрицу переходных вероятностей, состоящую из одношаговых переходных вероятностей, чтобы найти распределение длины серии.

Отметим, что середины (определённые на шаге 3) могут быть найдены с использованием следующей общей расчётной формулы

\[
S_{j} = LCL + (2(s + j) + 1)\tau \quad \text{и} \quad S_{0} = 0
\]

поскольку мы предполагаем (для простоты), что \(-LCL = UCL\). Говорят, что статистика для нанесения на карту находится в поглощающем состоянии или состоянии OOC (т.е. \(j = s + 1\)), если \(Z_{i}\) попадает на или за пределы контрольных границ. Эта область считается поглощающей, поскольку процесс останавливается, когда карта подаёт сигнал. Следовательно, процесс объявляется находящимся в состоянии OOC всякий раз, когда \(Z_{i}\) находится в поглощающем состоянии, тогда как процесс считается находящимся в состоянии IC всякий раз, когда \(Z_{i}\) находится в переходном состоянии (также называемом состоянием IC). Таким образом, область IC состоит из \(v\) невозвратных состояний, тогда как область OOC рассматривается как одно поглощающее состояние.

Как объяснялось ранее, для применения подхода MC нам нужны выражения для вероятностей сигнала, чтобы мы могли построить матрицу переходных вероятностей. Элементы существенной матрицы переходных вероятностей \(Q\) называются одношаговыми переходными вероятностями; \(Q = [p_{ij}]\) для \(i, j = -s, -s+1, \ldots, s-1, s\). Чтобы рассчитать эти вероятности после того, как состояния определены или идентифицированы, мы предполагаем, что статистика для нанесения на карту равна середине \(S_{i}\) всякий раз, когда она находится в состоянии \(i\). Это приближение, которое становится более точным при увеличении числа подинтервалов, и комментируется позже. Для каждой рассматриваемой контрольной карты одношаговые переходные вероятности рассчитываются по-разному. Далее мы приводим ещё один пример использования подхода MC для расчёта распределения длины серии карты CUSUM.

\begin{figure}[!ht]
	\centering
	\includegraphics[width=0.7\linewidth]{"./images/fig-3.21.png"}
	\caption[]{Разбиение интервала между \(LCL\) и \(UCL\) на \(\nu = 2s + 1\) подинтервалов}
	\label{fig-3.21}
\end{figure}

\noindent\textbf{Пример 3.6. Распределение длины серии карты CUSUM для нормального среднего с индивидуальными данными}

Рассмотрим только верхнюю одностороннюю статистику для нанесения на карту CUSUM, определённую в уравнении \eqref{eq:3.3}. Предположим, что наблюдения следуют стандартному нормальному распределению, так что \(\mu_0\) в уравнении \eqref{eq:3.3} равно нулю. Brook and Evans (1972) предложили метод оценки «не только средних длин серий для схемы, но также моментов и процентных точек распределения длины серии и точных вероятностей длины серии» для верхней односторонней карты CUSUM, рассматривая CUSUM как MC с пространством состояний, являющимся подмножеством \(\{0, 1, \ldots, H\}\). Методы MC имеют большое преимущество, поскольку они адаптируемы ко многим задачам, связанным с сериями, и часто упрощают решения конкретных задач, к которым они применяются. Fu, Spiring и Xie (2002) приводят три результата, которые должны быть выполнены перед реализацией подхода MC с конечным числом состояний. Пусть \(S_t^+\) — конечная однородная MC на пространстве состояний \(\Omega^+\) с матрицей переходных вероятностей, такой что (i) \(\Omega^+ = \{\zeta_0, \zeta_1, \ldots, \zeta_{r+s-1}\}\), где \(0 = \zeta_0 < \zeta_1 < \ldots < \zeta_{r+s-1} = H\) и \(\zeta_{r+s-1}\) является поглощающим состоянием; (ii) матрица переходных вероятностей задаётся как \(M = [p_{ij}]\) для \(i = 0, 1, \ldots, r+s-1\) и \(j = 0, 1, \ldots, r+s-1\), где \(r\) обозначает количество невозвратных состояний, а \(s\) обозначает количество поглощающих состояний соответственно, и (iii) начальное значение должно находиться в «фиктивном» состоянии с вероятностью единица, то есть \(P(S_0^+ = \zeta_0) = 1\), чтобы гарантировать, что процесс начинается в состоянии IC. Предположим, что MC \(S_t^+\) удовлетворяет условиям (i), (ii) и (iii). Тогда из Fu, Spiring и Xie (2002) и Fu and Lou (2003) следует, что распределение длины серии и его различные характеристики могут быть получены с использованием уравнений \eqref{eq:3.21}–\eqref{eq:3.24}.

На этом этапе следует отметить, что для карт CUSUM количество невозвратных состояний является фиксированным числом (см. пример ниже). Нет необходимости дискретизировать интервал между \(LCL\) и \(UCL\) на \(\nu = 2s + 1\) подинтервалов (как объяснялось на шаге 1 ранее), поскольку, как указывалось ранее, для карт CUSUM количество невозвратных состояний будет фиксированным числом.

Время, когда процедура подаёт сигнал, — это первый момент, когда конечная MC \(S_t^+\) входит в одно из поглощающих состояний, где пространство состояний задаётся как

\begin{equation}
	\begin{aligned}
		&\Omega^{+} = \{\zeta_{0}, \zeta_{1}, \ldots, \zeta_{r+s-1}\}, \quad S_{0}^{+} = 0 \\
		&\text{и} \\
		&S_{t}^{+} = \min\{H, \max\{0, S_{t-1}^{+} + X_{t} - K\}\}.
	\end{aligned}
	\label{eq:3.27}
\end{equation}

Состояние, соответствующее сигналу, поданному картой CUSUM, называется поглощающим состоянием. Очевидно, что существует только одно поглощающее состояние, поскольку карта подаёт сигнал, когда \(S_t^+\) попадает на или выше \(H\), то есть \(s = 1\).

При рассмотрении подхода MC для карты CUSUM Brook and Evans (1972) начинают с рассмотрения дискретного случая. Пусть \(X\) — положительная целочисленная случайная величина, а опорное значение \(K\) и решающий интервал \(H\) оба принимают положительные целочисленные значения. Дробные значения могут быть охвачены соответствующим масштабированием. Определяя \(X\), \(K\) и \(H\) таким образом, CUSUM \(\sum(X - K)\) может принимать только одно из целочисленных значений \(\{0, 1, \ldots, H\}\), например, для \(H = 4\) и \(K = 1\) пространство состояний рассчитывается в таблице \ref{tab:3.9}.

Следует отметить, что показаны только значения \(X_t = -5(1)5\). Любое значение \(X_t\) меньше \(-5\) приведёт к \(S_t^+ = 0\), а любое значение \(X_t\) больше \(+5\) приведёт к \(S_t^+ = 4\). Таким образом, пространство состояний задаётся таблицей \ref{tab:3.10}.

Из таблицы \ref{tab:3.10} видим, что существует три невозвратных состояния, то есть \(r = 3\), и одно поглощающее состояние, то есть \(s = 1\). Следовательно, соответствующая матрица переходных вероятностей будет матрицей \(4 \times 4\). Используя три состояния IC, матрица переходных вероятностей задаётся как

\[
M_{4 \times 4} = \left( \begin{array}{ccc|c}
	p_{00} & p_{01} & p_{02} & p_{03} \\
	p_{10} & p_{11} & p_{12} & p_{13} \\
	p_{20} & p_{21} & p_{22} & p_{23} \\
	\hline
	0 & 0 & 0 & 1
\end{array} \right)
\]

\begin{table}[htbp]
	\centering
	\caption{Расчёт пространства состояний при \(H = 4\) и \(K = 1\).}
	\label{tab:3.9}
	\renewcommand{\arraystretch}{1.1}
	\begin{tabular}{cccc}
		\toprule
		\(X_t\) & \(S_{t-1}^{+} + X_t - K\) & \(\max\{0, S_{t-1}^{+} + X_t - K\}\) & \(S_t^+ = \min\{H, \max\{0, S_{t-1}^{+} + X_t - K\}\}\) \\
		\midrule
		\(-5\) & \(0 + (-5) - 1 = -6\) & \(\max\{0, -6\} = 0\) & \(\min\{4, 0\} = 0\) \\
		\(-4\) & \(0 + (-4) - 1 = -5\) & \(\max\{0, -5\} = 0\) & \(\min\{4, 0\} = 0\) \\
		\(-3\) & \(0 + (-3) - 1 = -4\) & \(\max\{0, -4\} = 0\) & \(\min\{4, 0\} = 0\) \\
		\(-2\) & \(0 + (-2) - 1 = -3\) & \(\max\{0, -3\} = 0\) & \(\min\{4, 0\} = 0\) \\
		\(-1\) & \(0 + (-1) - 1 = -2\) & \(\max\{0, -2\} = 0\) & \(\min\{4, 0\} = 0\) \\
		\(0\) & \(0 + 0 - 1 = -1\) & \(\max\{0, -1\} = 0\) & \(\min\{4, 0\} = 0\) \\
		\(1\) & \(0 + 1 - 1 = 0\) & \(\max\{0, 0\} = 0\) & \(\min\{4, 0\} = 0\) \\
		\(2\) & \(0 + 2 - 1 = 1\) & \(\max\{0, 1\} = 1\) & \(\min\{4, 1\} = 1\) \\
		\(3\) & \(0 + 3 - 1 = 2\) & \(\max\{0, 2\} = 2\) & \(\min\{4, 2\} = 2\) \\
		\(4\) & \(0 + 4 - 1 = 3\) & \(\max\{0, 3\} = 3\) & \(\min\{4, 3\} = 3\) \\
		\(5\) & \(0 + 5 - 1 = 4\) & \(\max\{0, 4\} = 4\) & \(\min\{4, 4\} = 4\) \\
		\bottomrule
	\end{tabular}
\end{table}

\begin{table}[htbp]
	\centering
	\caption{Классификация и описание состояний.}
	\label{tab:3.10}
	\renewcommand{\arraystretch}{1.1}
	\begin{tabular}{ccc}
		\toprule
		Номер состояния & Описание состояния & Поглощающее (A) / Невозвратное (NA) \\
		\midrule
		0 & \(S_t^+ = 0\) & NA \\
		1 & \(S_t^+ = 1\) & NA \\
		2 & \(S_t^+ = 3\) & NA \\
		3 & \(S_t^+ = 4\) & A \\
		\bottomrule
	\end{tabular}
\end{table}

Напомним, что для получения характеристик распределения длины серии с использованием уравнений \eqref{eq:3.22}–\eqref{eq:3.24} необходима только существенная матрица переходных вероятностей \(Q\). Здесь существенная матрица переходных вероятностей задаётся как \(Q_{3 \times 3} = \left( \begin{array}{ccc} p_{00} & p_{01} & p_{02} \\ p_{10} & p_{11} & p_{12} \\ p_{20} & p_{21} & p_{22} \end{array} \right)\). Расчёт одношаговых переходных вероятностей для существенной матрицы переходных вероятностей приведён в таблице \ref{tab:3.11}.

Здесь существенная матрица переходных вероятностей задаётся как \(Q_{3 \times 3} = \left( \begin{array}{ccc} 0.841 & 0 & 0 \\ 0.500 & 0 & 0 \\ 0.023 & 0 & 0 \end{array} \right)\), \(\underline{1}_{3 \times 1} = \left( \begin{array}{c} 1 \\ 1 \\ 1 \end{array} \right)\), \(\underline{\xi}_{1 \times \nu} = \left( \begin{array}{ccc} 1 & 0 & 0 \end{array} \right)\) и \(I_{3 \times 3} = \left( \begin{array}{ccc} 1 & 0 & 0 \\ 0 & 1 & 0 \\ 0 & 0 & 1 \end{array} \right)\). Подставляя эти значения в уравнения \eqref{eq:3.22} и \eqref{eq:3.23}, находим, что \(ARL_{\text{IC}}\) и \(SDRL_{\text{IC}}\) равны 6.29 и 5.77 соответственно. Таким образом, когда процесс находится в состоянии IC, первая ложная тревога будет наблюдаться в среднем на каждой шестой нанесённой точке. Это, конечно, нехорошо, поскольку желательна высокая \(ARL_{\text{IC}}\). Однако для иллюстрации значения параметров проектирования \(K\) и \(H\) были выбраны так, чтобы получить малую существенную матрицу переходных вероятностей. Малая существенная матрица переходных вероятностей, здесь порядка \(3 \times 3\), желательна для иллюстрации, поскольку это приводит к показу только девяти расчётов одношаговых переходных вероятностей в таблице \ref{tab:3.11}. Продолжая расчёты характеристик длины серии в состоянии IC, некоторые репрезентативные процентили (скажем, 5-й, 25-й, 50-й, 75-й и 95-й) могут быть рассчитаны с использованием уравнения \eqref{eq:3.24}, в результате чего 5-й процентиль в состоянии IC равен 1, 25-й процентиль/первый квартиль в состоянии IC равен 2, 50-й процентиль/медиана в состоянии IC равен 5, 75-й процентиль/третий квартиль в состоянии IC равен 9, а 95-й процентиль в состоянии IC равен 18. Для лучшего понимания процентилей рассмотрим первый квартиль, равный 2, и медиану, равную 5. Поскольку первый квартиль равен 2, мы знаем, что ложная тревога не произойдёт в пределах первых двух нанесённых точек с вероятностью не более 75\%. Для медианы, поскольку медиана равна 5, мы знаем, что первая ложная тревога будет наблюдаться в пределах первых пяти нанесённых точек, по крайней мере, в 50\% случаев. В этом примере мы проиллюстрировали случай IC; однако расчёты для случая, когда процесс находится в состоянии OOC, могут быть легко выполнены. Это оставлено читателю в качестве упражнения.

\begin{table}[htbp]
	\centering
	\caption{Расчёт одношаговых переходных вероятностей существенной матрицы переходных вероятностей.}
	\label{tab:3.11}
	\renewcommand{\arraystretch}{1.3} % Немного увеличиваем отступ между строками для читаемости
	\small % Уменьшаем шрифт, чтобы длинные формулы не вылезали за поля
	\begin{tabular}{lll}
		\toprule
		\(p_{00}\) & \(p_{01}\) & \(p_{02}\) \\
		\midrule
		\(= P(S_t^+ = 0 \mid S_{t-1}^+ = 0)\) & \(= P(S_t^+ = 1 \mid S_{t-1}^+ = 0)\) & \(= P(S_t^+ = 3 \mid S_{t-1}^+ = 0)\) \\
		\(= P(\min\{4, \max\{0, X_t - 1\}\} = 0)\) & \(= P(\min\{4, \max\{0, X_t - 1\}\} = 1)\) & \(= P(\min\{4, \max\{0, X_t - 1\}\} = 3)\) \\
		\(= P(\max\{0, X_t - 1\} = 0)\) & \(= P(\max\{0, X_t - 1\} = 1)\) & \(= P(\max\{0, X_t - 1\} = 3)\) \\
		\(= P(X_t - 1 \leq 0)\) & \(= P(X_t - 1 = 1)\) & \(= P(X_t - 1 = 3)\) \\
		\(= P(X_t \leq 1)\) & \(= P(X_t = 2)\) & \(= P(X_t = 4)\) \\
		\(= 0.841\) & \(= 0\) & \(= 0\) \\
		\midrule
		\(p_{10}\) & \(p_{11}\) & \(p_{12}\) \\
		\midrule
		\(= P(S_t^+ = 0 \mid S_{t-1}^+ = 1)\) & \(= P(S_t^+ = 1 \mid S_{t-1}^+ = 1)\) & \(= P(S_t^+ = 3 \mid S_{t-1}^+ = 1)\) \\
		\(= P(\min\{4, \max\{0, X_t\}\} = 0)\) & \(= P(\min\{4, \max\{0, X_t\}\} = 1)\) & \(= P(\min\{4, \max\{0, X_t\}\} = 3)\) \\
		\(= P(X_t \leq 0)\) & \(= P(X_t = 1)\) & \(= P(X_t = 3)\) \\
		\(= 0.500\) & \(= 0\) & \(= 0\) \\
		\midrule
		\(p_{20}\) & \(p_{21}\) & \(p_{22}\) \\
		\midrule
		\(= P(S_t^+ = 0 \mid S_{t-1}^+ = 3)\) & \(= P(S_t^+ = 1 \mid S_{t-1}^+ = 3)\) & \(= P(S_t^+ = 3 \mid S_{t-1}^+ = 3)\) \\
		\(= P(\min\{4, \max\{0, 2 + X_t\}\} = 0)\) & \(= P(\min\{4, \max\{0, 2 + X_t\}\} = 1)\) & \(= P(\min\{4, \max\{0, 2 + X_t\}\} = 3)\) \\
		\(= P(\max\{0, 2 + X_t\} = 0)\) & \(= P(\max\{0, 2 + X_t\} = 1)\) & \(= P(\max\{0, 2 + X_t\} = 3)\) \\
		\(= P(2 + X_t \leq 0)\) & \(= P(2 + X_t = 1)\) & \(= P(2 + X_t = 3)\) \\
		\(= P(X_t \leq -2)\) & \(= P(X_t = -1)\) & \(= P(X_t = 1)\) \\
		\(= 0.023\) & \(= 0\) & \(= 0\) \\
		\bottomrule
	\end{tabular}
\end{table}

Мы иллюстрируем подход MC для расчёта распределения длины серии карты EWMA. Напомним, что в этом случае мы сначала должны дискретизировать матрицу переходных вероятностей с бесконечным числом состояний.

\noindent\textbf{Пример 3.7. Распределение длины серии карты EWMA для нормального среднего с индивидуальными данными}

После дискретизации матрицы переходных вероятностей нам нужно найти одношаговые переходные вероятности; \(M = [p_{ij}]\) для \(i, j = -s, -s+1, \ldots, s-1, s\). Чтобы рассчитать эти вероятности, мы предполагаем, что статистика для нанесения на карту равна \(S_{i}\) всякий раз, когда она находится в состоянии \(i\). Для невозвратного \(j\) получаем

\begin{equation}
p_{ij} = P(S_j - \tau < Z_k \leq S_j + \tau \mid \text{в } S_i) \quad \forall j \ (j \neq s)
\label{eq:3.28a}
\end{equation}

и

\begin{equation}
p_{ij} = P(S_j - \tau < Z_k < S_j + \tau \mid \text{в } S_i) \quad \text{для } j = s.
\label{eq:3.28b}
\end{equation}

Отметим, что индекс \(k\) не указывает на то, что статистика для нанесения на карту зависит от времени. Он просто относится к рассматриваемой статистике для нанесения на карту. Для поглощающего \(j\) мы используем тот факт, что строки \(M\) должны в сумме давать единицу, и поэтому вероятности перехода из переходного состояния в поглощающее состояние находятся вычитанием. Нахождение существенной матрицы переходных вероятностей \(Q\) включает нахождение вероятностей перехода из одного переходного состояния в другое. Чтобы рассчитать переходные вероятности для параметрической контрольной карты EWMA, мы используем статистику для нанесения на карту и подставляем её в уравнение \eqref{eq:3.28a}, получая

\[
\begin{aligned}
p_{ij} &= P(S_j - \tau < \lambda X_k + (1 - \lambda)Z_{k-1} \leq S_j + \tau \mid Z_{k-1} = S_i) \quad \forall j \ (j \neq s)
\end{aligned}
\]
и
\begin{equation}
p_{ij} = P(S_j - \tau < \lambda X_k + (1 - \lambda)Z_{k-1} < S_j + \tau \mid Z_{k-1} = S_i) \quad \text{для } j = s.
\label{eq:3.29}
\end{equation}

Переписывая и подставляя \(Z_{k-1} = S_{i}\) в уравнение \eqref{eq:3.29}, получаем

\begin{equation}
p_{ij} = P\left(\frac{S_j - \tau - (1 - \lambda)S_i}{\lambda} < X_k \leq \frac{S_j + \tau - (1 - \lambda)S_i}{\lambda}\right) \quad \forall j \ (j \neq s)
\label{eq:3.30a}
\end{equation}
и
\[
p_{ij} = P\left(\frac{S_j - \tau - (1 - \lambda)S_i}{\lambda} < X_k < \frac{S_j + \tau - (1 - \lambda)S_i}{\lambda}\right) \quad (j = s)
\]
которые могут быть получены с использованием
\[
p_{ij} = \Phi\left(\frac{S_j + \tau - (1 - \lambda)S_i}{\lambda}\right) - \Phi\left(\frac{S_j - \tau - (1 - \lambda)S_i}{\lambda}\right) \quad \forall j \ (j \neq s)
\]

где \(\Phi\) обозначает функцию распределения стандартного нормального распределения. На этом этапе важно отметить, что для расчёта одношаговых переходных вероятностей, необходимых для расчёта распределения длины серии контрольной карты, требуется только лежащее в основе распределение процесса.

Поскольку значения \(\tau\), \(\lambda\), \(S_{i}\) и \(S_{j}\) являются известными константами (\(\lambda\) выбирается практиком, а \(\tau\), \(S_{i}\) и \(S_{j}\) рассчитываются), требуемые вероятности могут быть легко рассчитаны. Для получения более подробной информации о подходе MC читатель отсылается к Приложению 3.2.

Для численной иллюстрации предположим, что мы рассматриваем параметрическую карту EWMA для индивидуальных данных из нормального распределения с параметрами проектирования \(\lambda = 0.05\) и \(L = 2.703\) (это приводит к \(ARL_{\text{IC}} = 370\); см. таблицу \ref{tab:3.2}), и \(v = 2s + 1 = 9\), так что \(s = 4\), то есть мы делим интервал между контрольными границами на девять подинтервалов (см. рисунок \ref{fig-3.22}), что приводит к «малой» матрице переходных вероятностей, имеющей только 10 строк и 10 столбцов соответственно.

\begin{figure}[!ht]
\centering
\includegraphics[width=0.7\linewidth]{"./images/fig-3.22.png"}
\caption[]{Разбиение интервала между \(LCL\) и \(UCL\) на \(\nu = 9\) подинтервалов при \(\lambda = 0.10\) и \(L = 2.703\)}
\label{fig-3.22}
\end{figure}

Подставляя \(\mu_0 = 0\) и \(\sigma_0 = 1\) для стандартного нормального распределения и параметры проектирования \(\lambda = 0.1\) и \(L = 2.703\) в уравнение \eqref{eq:3.12}, получаем \(CL\) и контрольные границы стационарного состояния:

\[
\begin{aligned}
UCL &= \mu_0 + L\sigma_0\sqrt{\frac{\lambda}{2 - \lambda}} = 0 + 2.703(1)\sqrt{\frac{0.1}{2 - 0.1}} = 0.620 \\
CL &= \mu_0 = 0 \\
LCL &= \mu_0 - L\sigma_0\sqrt{\frac{\lambda}{2 - \lambda}} = 0 - 2.703(1)\sqrt{\frac{0.1}{2 - 0.1}} = -0.620.
\end{aligned}
\]

Разделив интервал между \(LCL\) и \(UCL\) на \(\nu = 9\) подинтервалов шириной \(2\tau = (UCL - LCL)/\nu = (0.620 - (-0.620))/9 = 0.138\), так что \(\tau = 0.069\). Разбиение проиллюстрировано на рисунке \ref{fig-3.22}.

Для девяти состояний IC матрица переходных вероятностей задаётся как

\[
M_{10 \times 10} = \left( \begin{array}{cccccccccc}
p_{(-4)(-4)} & \cdots & p_{(-4)0} & \cdots & p_{(-4)4} & p_{(-4)5} \\
\vdots & \ddots & \vdots & & \vdots & \vdots \\
p_{0(-4)} & \cdots & p_{00} & \cdots & p_{04} & p_{05} \\
\vdots & & \vdots & \ddots & \vdots & \vdots \\
p_{4(-4)} & \cdots & p_{40} & \cdots & p_{44} & p_{45} \\
\hline
0 & \cdots & 0 & \cdots & 0 & 1
\end{array} \right)
\]

где переходные вероятности рассчитываются с использованием уравнения \eqref{eq:3.30a}. Расчёт \(p_{01}\) показан ниже.

\[
\begin{aligned}
p_{01} &= \Phi\left(\frac{S_1 + \tau - (1 - \lambda)S_0}{\lambda}\right) - \Phi\left(\frac{S_1 - \tau - (1 - \lambda)S_0}{\lambda}\right) \\
&= \Phi\left(\frac{0.138 + 0.069 - (1 - 0.1)0}{0.1}\right) - \Phi\left(\frac{0.138 - 0.069 - (1 - 0.1)0}{0.1}\right) \\
&= \Phi(2.067) - \Phi(0.689) \\
&= 0.226.
\end{aligned}
\]

После того как одношаговые переходные вероятности рассчитаны и найдена существенная матрица переходных вероятностей \(Q\), характеристики длины серии в состоянии IC получаются с использованием уравнений \eqref{eq:3.22}–\eqref{eq:3.24}. Это даёт \(ARL_{\text{IC}}\) и \(SDRL_{\text{IC}}\), равные 270.26 и 263.07 соответственно, с 5-м, 25-м, 50-м, 75-м и 95-м процентилями в состоянии IC, равными 21, 83, 190, 372 и 795 соответственно. Таким образом, когда процесс находится в состоянии IC, первая ложная тревога будет наблюдаться в среднем на каждой 270-й нанесённой точке. Поскольку первый квартиль равен 83, мы знаем, что ложная тревога не произойдёт в пределах первых 83 нанесённых точек с вероятностью не более 75\%. Для медианы, поскольку медиана равна 190, мы знаем, что первая ложная тревога будет наблюдаться в пределах первых 190 нанесённых точек, по крайней мере, в 50\% случаев.

Следует отметить на этом этапе, что более точный ответ мог бы быть получен, если бы количество подинтервалов между контрольными границами было увеличено; например, для \(\nu = 1001\) находим \(ARL_{\text{IC}} = 371.88\), что ближе к желаемому номинальному значению 370. Однако для этого примера было взято небольшое количество подинтервалов, чтобы проиллюстрировать процедуру.

Далее мы иллюстрируем подход MC для расчёта распределения длины серии карты Шухарта с правилами сигнализации на основе серий. Напомним, что в этом случае мы сначала должны дискретизировать матрицу переходных вероятностей с бесконечным числом состояний.

\noindent\textbf{Пример 3.8. Распределение длины серии контрольной карты Шухарта для мониторинга среднего процесса с дополнительными правилами серий}

Продолжая использовать те же данные, что и для иллюстрации контрольной карты Шухарта для мониторинга среднего процесса, здесь мы иллюстрируем контрольную карту Шухарта \(\overline{X}\) с правилами сигнализации на основе серий. Данные и график опущены для экономии места; однако читатель отсылается к карте Шухарта \(\overline{X}\) 1-из-1 в разделе 3.3.1, пример 3.1, для этих данных, статистик для нанесения на карту и т.д. Напомним, что правила сигнализации на основе серий 2-из-2 и 2-из-3 требуют, чтобы сигнализировали как минимум последние две или три статистики для нанесения на карту соответственно. Однако из рисунка \ref{fig-3.3} видно, что все статистики для нанесения на карту находятся между контрольными границами и, следовательно, процесс функционирует в состоянии статистической управляемости относительно целевого среднего.

Здесь мы иллюстрируем подход MC для расчёта распределения длины серии двусторонней карты Шухарта 2-из-2 для среднего, которая подаёт сигнал, когда две из последних двух последовательных статистик для нанесения на карту выходят за пределы контрольных границ. Эта карта была предложена Klein (2000). Карта 2-из-2 является одним из нескольких типов правил сигнализации на основе серий, которые были предложены для повышения эффективности базовой карты Шухарта. Мы обсуждаем их более подробно в разделе 3.5.2.

Отметим, что карта 2-из-2 подаёт сигнал только тогда, когда обе из двух последовательных статистик для нанесения на карту отображаются на или выше UCL (см. панель (a) рисунка \ref{fig-3.17}), или когда обе отображаются на или ниже LCL (см. панель (b) рисунка \ref{fig-3.17}). В первом случае сигнализируется увеличение или восходящий сдвиг среднего, тогда как в последнем случае обнаруживается уменьшение или нисходящий сдвиг.

Подход MC для расчёта распределения длины серии состоит в первую очередь в получении матрицы переходных вероятностей. Для этого нам нужно определить количество поглощающих и невозвратных состояний. Следует подчеркнуть, что при использовании правил сигнализации на основе серий количество невозвратных состояний определяется типом используемого правила, например, 2-из-2 или 2-из-3. Нет необходимости дискретизировать интервал между \(LCL\) и \(UCL\) на \(\nu = 2s + 1\) подинтервалов (как объяснялось на шаге 1 ранее), поскольку, как указывалось ранее, для правил сигнализации на основе серий количество невозвратных состояний зависит от типа правила. Для этого сначала определим три невозвратных состояния MC, а именно:

\noindent\textbf{Состояние 1:} Ни одна точка не выходит за пределы любой из контрольных границ.

\noindent\textbf{Состояние 2:} Только одна точка находится на или выше UCL.

\noindent\textbf{Состояние 3:} Только одна точка находится на или ниже LCL.

Подход MC содержит одно поглощающее состояние, а именно:

\noindent\textbf{Состояние 4:} Как только две последовательные точки оказываются на или за пределами любой из контрольных границ, процесс достигает поглощающего состояния.

Далее мы определяем вероятности, соответствующие каждому из этих состояний. Пусть \(p\) обозначает вероятность того, что процесс находится в состоянии 1, то есть статистика для нанесения на карту отображается между \(LCL\) и \(UCL\). Аналогично, пусть \(pU\) и \(pL\) обозначают вероятности того, что случайный процесс находится в состоянии 2 или состоянии 3 соответственно, то есть статистика для нанесения на карту отображается на или выше UCL или на или ниже LCL. Состояния и связанные с ними вероятности проиллюстрированы на рисунке \ref{fig-3.23} для карты 2-из-2.

\begin{figure}[!ht]
\centering
\includegraphics[width=0.7\linewidth]{"./images/fig-3.23.png"}
\caption[]{Иллюстрация состояний и связанных с ними вероятностей для карты 2-из-2}
\label{fig-3.23}
\end{figure}

Одношаговые переходные вероятности для правила сигнализации на основе серий 2-из-2 показаны у Klein (2000), и читатель отсылается к его статье для получения подробностей. Однако для получения распределения длины серии и связанных с ним характеристик с использованием уравнений \eqref{eq:3.22}–\eqref{eq:3.24} необходима только существенная матрица переходных вероятностей \(Q\). Для нашего примера существенная матрица переходных вероятностей задаётся как \(Q_{3 \times 3} = \begin{pmatrix} p & pU & pL \\ p & 0 & pL \\ p & pU & 0 \end{pmatrix}\). Таким образом, для иллюстрации, предполагая, что когда процесс находится в состоянии IC, \(p = 0.93\), \(pL = pU = 0.035\), так что \(Q_{3 \times 3} = \begin{pmatrix} 0.93 & 0.035 & 0.035 \\ 0.93 & 0 & 0.035 \\ 0.93 & 0.035 & 0 \end{pmatrix}\), \(\underline{1}_{3 \times 1} = \begin{pmatrix} 1 \\ 1 \\ 1 \end{pmatrix}\), \(\underline{\xi}_{1 \times 3} = \begin{pmatrix} 1 & 0 & 0 \end{pmatrix}\) и \(I_{3 \times 3} = \begin{pmatrix} 1 & 0 & 0 \\ 0 & 1 & 0 \\ 0 & 0 & 1 \end{pmatrix}\). Подставляя эти значения в уравнения \eqref{eq:3.22} и \eqref{eq:3.23}, находим, что \(ARL_{\text{IC}}\) и \(SDRL_{\text{IC}}\) равны 422.45 и 420.98 соответственно. Таким образом, когда процесс находится в состоянии IC, первая ложная тревога будет наблюдаться в среднем на каждой 422-й точке времени или нанесённой точке. Некоторые репрезентативные процентили (скажем, 5-й, 25-й, 50-й, 75-й и 95-й) могут быть рассчитаны с использованием функции распределения длины серии в уравнении \eqref{eq:3.24}, и мы находим, что 5-й процентиль в состоянии IC равен 23, 25-й процентиль/первый квартиль в состоянии IC равен 123, 50-й процентиль/медиана в состоянии IC равен 293, 75-й процентиль/третий квартиль в состоянии IC равен 585, а 95-й процентиль в состоянии IC равен 1263. Для лучшего понимания процентилей рассмотрим первый квартиль, равный 123, и медиану, равную 293. Поскольку первый квартиль равен 123, мы знаем, что ложная тревога не произойдёт в пределах первых 123 нанесённых точек с вероятностью не более 75\%. Для медианы, поскольку медиана равна 293, мы знаем, что первая ложная тревога будет наблюдаться в пределах первых 293 нанесённых точек, по крайней мере, в 50\% случаев. Аналогичные расчёты могут быть выполнены, когда процесс находится в состоянии OOC. Всё, что нам нужно сделать, это найти \(p\), \(pL\) и \(pU\), построить существенную матрицу переходных вероятностей и применить формулы, приведённые в уравнениях \eqref{eq:3.22}–\eqref{eq:3.24}. Отметим, однако, что, вообще говоря, для расчёта переходных вероятностей требуется распределение статистики для нанесения на карту. Например, отметим, что \(pU = P(\overline{X} > UCL)\), так что нам нужно распределение \(\overline{X}\) или распределение генеральной совокупности. Для задачи о нормальном среднем предполагается, что генеральная совокупность является нормальной, и вероятность \(pU\) может быть легко рассчитана. Это важный момент, который снова будет обсуждаться в Главе 4 для непараметрических карт.

Выше мы рассмотрели правило сигнализации на основе серий 2-из-2. Важно отметить, что, вообще говоря, можно иметь правило сигнализации на основе серий \(w\)-из-\((w + v)\), где, например, возьмём \(w = 2\) и \(v = 1\), тогда мы имеем правило сигнализации на основе серий 2-из-3, где сигнал подаётся, когда две из последних трёх последовательных точек отображаются в состоянии OOC. Этот подход аналогичен правилу сигнализации на основе серий 2-из-2; однако отметим, что количество состояний увеличивается до восьми. На этом этапе читателю напоминается, что при использовании правил сигнализации на основе серий количество невозвратных состояний определяется типом используемого правила, например, 2-из-2 или 2-из-3. Нет необходимости дискретизировать интервал между \(LCL\) и \(UCL\) на \(\nu = 2s + 1\) подинтервалов (как объяснялось на шаге 1 ранее), поскольку, как указывалось ранее, для правил сигнализации на основе серий количество невозвратных состояний зависит от типа правила. Здесь мы рассматриваем правило сигнализации на основе серий 2-из-3, где восемь состояний следующие:

\noindent\textbf{Состояние (OO):} Две последовательные точки между обеими контрольными границами.

\noindent\textbf{Состояние (OU):} Первая точка между обеими контрольными границами, а вторая точка выше \(UCL\).

\noindent\textbf{Состояние (OL):} Первая точка между обеими контрольными границами, а вторая точка ниже \(LCL\).

\noindent\textbf{Состояние (UL):} Первая точка выше \(UCL\), а вторая точка ниже \(LCL\).

\noindent\textbf{Состояние (LO):} Первая точка ниже \(LCL\), а вторая точка между контрольными границами.

\noindent\textbf{Состояние (LU):} Первая точка ниже \(LCL\), а вторая точка выше \(UCL\).

Подход MC содержит одно поглощающее состояние, а именно:

\noindent\textbf{Состояние (OOC):} Две из трёх последовательных точек либо ниже \(LCL\), либо выше \(UCL\).

Одношаговые переходные вероятности для правила сигнализации на основе серий 2-из-3 показаны у Klein (2000), и мы отсылаем читателя к его статье для получения более подробной информации. Однако, как и в случае правила 2-из-2, для получения распределения длины серии и связанных с ним характеристик с использованием уравнений \eqref{eq:3.22}–\eqref{eq:3.24} необходима только существенная матрица переходных вероятностей \(Q\). Для нашего примера существенная матрица переходных вероятностей задаётся как

\[
Q_{6 \times 6} = \begin{pmatrix}
p & pU & pL & 0 & 0 & 0 \\
0 & 0 & pL & pU & pL & 0 \\
0 & pU & 0 & 0 & pL & pU \\
0 & pU & pL & 0 & 0 & 0 \\
p & pU & 0 & 0 & 0 & 0 \\
p & 0 & pL & 0 & 0 & 0
\end{pmatrix}
\]

предполагая, что когда процесс находится в состоянии IC, \(p = 0.93\), \(pL = pU = 0.035\), так что

\[
Q_{6 \times 6} = \begin{pmatrix}
0.93 & 0.035 & 0.035 & 0 & 0 & 0 \\
0 & 0 & 0.035 & 0.035 & 0.035 & 0 \\
0 & 0.035 & 0 & 0 & 0.035 & 0.035 \\
0 & 0.035 & 0.035 & 0 & 0 & 0 \\
0.93 & 0.035 & 0 & 0 & 0 & 0 \\
0.93 & 0 & 0.035 & 0 & 0 & 0
\end{pmatrix}
\]

Подставляя эти значения в уравнения \eqref{eq:3.22} и \eqref{eq:3.23}, находим, что \(ARL_{\text{IC}}\) и \(SDRL_{\text{IC}}\) равны 222.13 и 220.23 соответственно. Таким образом, когда процесс находится в состоянии IC, первая ложная тревога будет наблюдаться в среднем на каждой 222-й точке времени или нанесённой точке. Некоторые репрезентативные процентили (скажем, 5-й, 25-й, 50-й, 75-й и 95-й) могут быть рассчитаны с использованием функции распределения длины серии в уравнении \eqref{eq:3.24}, и мы находим, что 5-й процентиль в состоянии IC равен 13, 25-й процентиль/первый квартиль в состоянии IC равен 65, 50-й процентиль/медиана в состоянии IC равен 155, 75-й процентиль/третий квартиль в состоянии IC равен 307, а 95-й процентиль в состоянии IC равен 662. Для лучшего понимания процентилей рассмотрим первый квартиль, равный 65, и медиану, равную 155. Поскольку первый квартиль равен 65, мы знаем, что ложная тревога не произойдёт в пределах первых 65 нанесённых точек с вероятностью не более 75\%. Для медианы, поскольку медиана равна 155, мы знаем, что первая ложная тревога будет наблюдаться в пределах первых 155 нанесённых точек, по крайней мере, в 50\% случаев.

Аналогичные расчёты могут быть выполнены, когда процесс находится в состоянии OOC. Всё, что нам нужно сделать, это найти \(p\), \(pL\) и \(pU\), построить существенную матрицу переходных вероятностей и применить формулы, приведённые в уравнениях \eqref{eq:3.22}–\eqref{eq:3.24}.

\subsubsection{Подход на основе интегральных уравнений}

Подход на основе интегральных уравнений использует математику и комбинаторику для нахождения выражения распределения длины серии в замкнутой форме. Этот подход иногда сложен тем, что получаемое выражение обычно является комплексным или трудным для численной оценки. Часто точное выражение распределения длины серии может быть найдено после значительной работы, но здесь вместо этого проводится моделирование, поскольку это намного проще.

\subsubsection{Подход на основе компьютерного моделирования (Монте-Карло)}

Имитации Монте-Карло могут быть использованы для расчёта характеристик распределения длины серии. Популярность этого метода обусловлена тем фактом, что, независимо от того, насколько сложно распределение длины серии в теории, компьютерные имитации почти всегда могут быть использованы с относительной лёгкостью для расчёта распределения длины серии и связанных с ним характеристик достаточно точно, при условии, что размер имитации достаточно велик. В данной работе мы используем 100 000 имитаций, поскольку хорошо известно, что ошибка характеристики длины серии может быть ограничена путём достаточного увеличения размера имитации.

Общие шаги процедуры компьютерного моделирования для расчёта распределения длины серии для двусторонней контрольной карты, где статистика для нанесения на карту рассчитывается по случайной выборке, приведены ниже. Это может быть реализовано в любом программном пакете, таком как R, SAS или Matlab.

\noindent\textbf{Шаг 1:} Задайте необходимые параметры, такие как объём подгруппы, номинальная средняя длина серии в состоянии IC (или FAR); рассчитайте контрольные границы.

\noindent\textbf{Шаг 2:} Предположите некоторое распределение процесса, такое как нормальное распределение, и задайте параметры. Сгенерируйте (смоделируйте) подгруппы случайных наблюдений из заданного распределения.

\noindent\textbf{Шаг 3:} Рассчитайте статистику для нанесения на карту для каждой подгруппы и сравните её с контрольными границами, рассчитанными на шаге 1.

\noindent\textbf{Шаг 4:} Количество подгрупп, необходимое до тех пор, пока статистика для нанесения на карту впервые не попадёт на или за пределы контрольных границ, записывается как наблюдение из распределения длины серии.

\noindent\textbf{Шаг 5:} Повторите шаги 1–4 в общей сложности 100 000 раз.

\noindent\textbf{Шаг 6:} После получения «набора данных» из 100 000 наблюдений из распределения длины серии могут быть рассчитаны интересующие характеристики распределения длины серии, которые являются описательными статистиками, такими как среднее, стандартное отклонение, медиана, процентили и т.д. Например, можно использовать proc univariate в SAS для нахождения этих значений. Minitab или R можно использовать для того же самого и для построения ящичных диаграмм или гистограмм значений длины серии для лучшего понимания распределения длины серии.

Можно отметить, что моделирование распределения длины серии и связанных с ним характеристик может быть особенно трудоёмким в случае IC, поскольку распределение длины серии имеет сильную правостороннюю асимметрию, что означает, что могут быть некоторые очень большие значения длины серии, которые могут потребовать очень много времени для реализации. В некоторых случаях нужно быть осторожным, чтобы интересующая характеристика длины серии, скажем, \(ARL\), была конечной, что не обязательно может иметь место. Возможным сокращением для расчёта распределения длины серии может быть использование явных выражений для расчёта вероятности сигнала (и обратной ей величины) в рамках имитаций, которые могут быть доступны для некоторых карт при предполагаемом распределении процесса.

\noindent\textbf{Замечание 3.2. Количество подинтервалов, которое следует использовать для карты EWMA}

При применении подхода MC для расчёта распределения длины серии в картах EWMA необходимо принять решение относительно количества подинтервалов между контрольными границами. Здесь мы приводим иллюстрацию с использованием параметрической карты EWMA. Однако аналогичная картина разворачивается для всех карт EWMA. Мы также сопоставляем результаты подхода MC с результатами подхода на основе имитационного моделирования.

Для подхода MC, когда количество подинтервалов \(\nu\) достаточно велико, он обеспечивает эффективный метод, позволяющий точно аппроксимировать свойства длины серии схемы контроля EWMA. На практике значения \(\nu\) около 100 дают удовлетворительные приближения (см. Yu, 2007, с. 20). Однако результаты имитационного моделирования (см. таблицу \ref{tab:3.12}) показывают, что расхождения между смоделированными значениями и значениями, найденными с использованием подхода MC, могут быть несколько большими для \(\nu \leq 100\). Для больших значений \(\nu\), превышающих 100, расхождения кажутся малыми, и они особенно малы при использовании большего значения, такого как \(\nu = 1001\). В принципе, взятие больших значений \(\nu\) должно приводить к более точным ответам при использовании подхода MC, но при этом некоторые характеристики длины серии не могут быть вычислены за практический период времени. Комбинации \((\lambda, L)\), использованные в таблице \ref{tab:3.12}, взяты из таблицы \ref{tab:3.2}, и числа в каждой ячейке показывают достигнутое значение \(ARL_{\text{IC}}\) и процентное отличие от номинального значения \(ARL_{\text{IC}}\), равного 370. В целом, однако, для расчёта распределения длины серии карты EWMA рекомендуется метод имитационного моделирования.

\begin{table}[htbp]
\centering
\caption{Значения \(ARL_{\text{IC}}\) и абсолютное процентное отличие между смоделированными значениями и значениями, полученными с использованием подхода MC для параметрической (нормальной теории) карты EWMA для различных значений \(\nu\).}
\label{tab:3.12}
\renewcommand{\arraystretch}{1.1}
\begin{tabular}{c cc cc}
\toprule
\multirow{2}{*}{} & \multicolumn{2}{c}{\(\lambda = 0.05\), \(L = 2.492\)} & \multicolumn{2}{c}{\(\lambda = 0.10\), \(L = 2.703\)} \\
\cmidrule(lr){2-3} \cmidrule(lr){4-5}
& Достигнутая ARL & Отличие (\%) & Достигнутая ARL & Отличие (\%) \\
\midrule
\multicolumn{5}{c}{Подход на основе цепей Маркова} \\
\midrule
\(\nu = 5\)   & 439.07 & (18.62\%) & 384.92 & (4.06\%) \\
\(59\)       & 386.23 & (4.34\%)  & 420.36 & (13.64\%) \\
\(67\)       & 379.34 & (2.48\%)  & 337.55 & (8.75\%) \\
\(75\)       & 410.76 & (10.97\%) & 360.21 & (2.62\%) \\
\(83\)       & 410.48 & (10.90\%) & 423.25 & (14.42\%) \\
\(101\)      & 396.12 & (7.02\%)  & 357.42 & (3.38\%) \\
\(151\)      & 380.61 & (2.83\%)  & 358.61 & (3.05\%) \\
\(201\)      & 391.25 & (5.70\%)  & 330.66 & (10.61\%) \\
\(251\)      & 368.50 & (0.45\%)  & 354.17 & (4.26\%) \\
\(301\)      & 369.42 & (0.20\%)  & 379.63 & (2.63\%) \\
\(401\)      & 365.71 & (1.20\%)  & 370.68 & (0.21\%) \\
\(501\)      & 377.47 & (1.98\%)  & 377.33 & (2.01\%) \\
\(1001\)     & 369.49 & (0.18\%)  & 373.80 & (1.05\%) \\
\midrule
\multicolumn{5}{c}{100 000 имитаций} \\
\midrule
& 370.15 & & 369.91 & \\
\bottomrule
\end{tabular}
\end{table}

Из таблицы \ref{tab:3.12} видно, что чем больше значение \(\nu\), тем меньше процентное отличие между смоделированными значениями и значениями, полученными с использованием подхода MC. Когда параметры проектирования взяты равными \(\lambda = 0.05\) и \(L = 2.492\), например, мы имеем большое процентное отличие \(18.62\%\) при малом значении \(\nu\) (\(\nu = 51\)) по сравнению с малым процентным отличием \(0.18\%\) при большом значении \(\nu\) (\(\nu = 1001\)). Рекомендация состоит в том, чтобы взять \(\nu = 1001\), что может быть несколько трудоёмким, однако результаты будут более точными. С практической точки зрения значение \(\nu\), равное 251 или 301, также представляется разумным.

\section{Задача оценивания параметров и её влияние на эффективность контрольной карты}

В предыдущих разделах мы изучали ситуации, когда процесс предполагается нормально распределённым (или статистика для нанесения на карту распределена приблизительно нормально), а его параметры заданы или известны. Однако во многих практических ситуациях, даже если нормальность можно считать приемлемой, параметры процесса могут быть неизвестны, и их приходится оценивать до начала мониторинга процесса. Эта ситуация называется «случаем неизвестных стандартов» (standard(s) unknown) или случаем неизвестных параметров и обозначается как \textbf{случай U}. Хорошо известно, что игнорировать эффекты оценивания параметров может быть неразумно и дорого: свойства длины серии карты сильно меняются, что может приводить, например, к гораздо большему числу ложных тревог, чем ожидается номинально (см., например, Jensen et al., 2006 и Psarakis et al., 2014), а это снижает эффективность контрольных карт и увеличивает затраты.

Обычно неизвестные параметры оценивают по эталонной выборке фиксированного объёма, полученной в ходе анализа Фазы I. Ниже мы обсуждаем вывод распределения длины серии в случае U для трёх основных типов контрольных карт. Как выясняется, традиционный взгляд на распределение длины серии нужно пересмотреть с учётом корректного обращения с оценками параметров, используемыми в контрольных картах (границах).

Когда используются оценки параметров, под распределением длины серии контрольной карты обычно понимают то, что сегодня называют \textbf{безусловным} (unconditional) или \textbf{маргинальным} распределением длины серии; оно получается усреднением распределения длины серии для заданного набора оценок по распределению этих оценок. Распределение длины серии, полученное для заданного набора оценок, называется \textbf{условным} (conditional) распределением длины серии. В ранние годы развития SPC оценки параметров обычно вычисляли и подставляли (plug-in) вместо самих параметров, а контрольные границы (и контрольные карты) строили для последующего мониторинга процесса в Фазе II. Влияние оценивания параметров, по существу, поглощалось математико-статистическим шагом при выводе (безусловного) распределения длины серии. Для характеристики эффективности карты обычно использовали распределение длины серии и связанные с ним характеристики (ARL и т.д.). В последние годы было признано, что контрольные карты с подстановкой оценок ведут себя совсем не так, как их аналоги с известными параметрами, если только для вычисления оценок не использовалось огромное количество данных. Необходимый объём данных может быть значительным и на практике может оказаться недоступным или непозволительно большим.

Подход, основанный на безусловном распределении длины серии и связанных с ним свойствах, таких как безусловное ARL, возможно, и имел достоинства, поскольку давал информацию о «средней» эффективности карты, но гораздо больше практических сведений может дать условное распределение длины серии и его характеристики: они помогают объяснить эффективность контрольной карты для пользователя с теми конкретными оценками параметров, которые получены по его эталонным выборкам. В этом контексте представляет интерес среднее, или математическое ожидание, условного распределения длины серии — фактически условное среднее ARL, поскольку само условное ARL является случайной величиной. Однако нужно исследовать и другие свойства, например дисперсию распределения условного ARL, чтобы лучше понять эффективность карты при оценивании параметров. Отметим, что различие между безусловным и условным распределениями длины серии исчезает в случае известных параметров, когда оба они равны единственному распределению длины серии контрольной карты.

Общая тема всего этого — конечно же, объём данных Фазы I, необходимый для оценивания параметров, чтобы можно было гарантировать предсказуемую эффективность карты. Возможно, придётся также по-новому взглянуть на меры эффективности карты и на то, как проектировать контрольные карты (то есть получать правильные константы для нанесения на карту) для имеющегося объёма данных, чтобы гарантировать такую эффективность.

Теперь мы вернёмся к контрольным картам, которые обсуждались ранее в этой главе для случая известных параметров (случай K), и к случаю, когда параметры неизвестны, то есть к случаю U. Отметим, что мы по-прежнему предполагаем, что лежащее в основе распределение процесса нормально. Когда параметры неизвестны, их обычно оценивают по ретроспективному анализу исторических данных, который часто называют анализом Фазы I. Контрольные карты играют важную роль в анализе Фазы I; некоторые карты Фазы I мы обсуждаем в Главе 5. После завершения анализа Фазы I получают набор эталонных данных, который, как предполагается, происходит из процесса в состоянии IC. Эти данные используются для оценивания параметров и контрольных границ, которые затем применяются проспективно для мониторинга поступающих новых, или тестовых, данных. Эта часть статистического контроля процессов называется анализом Фазы II. Таким образом, в этом разделе мы сначала проиллюстрируем контрольные карты Фазы II, в которых неизвестные параметры процесса оцениваются по Фазе I. Мы также обсудим свойства контрольных карт для случая U и проиллюстрируем так называемый метод обусловливания/деобусловливания (conditioning/unconditioning, CUC) для вывода распределения длины серии и его различных характеристик, таких как средняя длина серии. Отметим, что в последнее время в этой области проделано много работы, но ещё больше работы продолжается, и следует ожидать новых результатов и более ясного понимания.

\section{Параметрические контрольные карты по количественным признакам в случае U}

\subsection{Контрольные карты Шухарта в случае U}

\subsubsection{Контрольные карты Шухарта для среднего в случае U}

В нашем обсуждении параметрических контрольных карт по количественным признакам в случае U мы предполагаем, что процесс следует нормальному распределению, но хотя бы один (или оба) параметра процесса — среднее и стандартное отклонение — неизвестен(ы). Хотя на практике можно предполагать и другие распределения процесса, мы рассматриваем случай нормального распределения, чтобы быть последовательными с предшествующим развитием в этой главе, относящимся к случаю известных параметров (см. Раздел 3.2). Прежде чем продолжить, внесём более тонкую классификацию внутри случая U. Когда задействованы два параметра, как в данном случае, возможны четыре ситуации, а именно: случай KK (среднее известно и стандартное отклонение известно), который часто называют просто случаем K, случай KU (среднее известно, стандартное отклонение неизвестно), случай UK (среднее неизвестно, стандартное отклонение известно) и случай UU (и среднее, и стандартное отклонение неизвестны). В этом разделе мы рассматриваем, пожалуй, наиболее практичную ситуацию — случай UU, когда оба параметра неизвестны. Таким образом, случай U означает одну из трёх ситуаций: KU, UK или UU. Другие ситуации можно рассматривать аналогичным образом (см., например, Chakraborti, 2000). Кроме того, мы предполагаем, что после успешного анализа Фазы I доступна эталонная выборка Фазы I, состоящая из \(m\) выборок объёма \(n\) каждая, для оценивания этих неизвестных параметров.

Предположим, что \(\overline{X}_1, \overline{X}_2, \ldots, \overline{X}_m\) и \(S_1, S_2, \ldots, S_m\) обозначают соответственно \(m\) средних и стандартных отклонений выборок Фазы I. Пусть также \(R_1, R_2, \ldots, R_m\) обозначают \(m\) размахов выборок Фазы I. Среднее процесса \(\mu\) обычно оценивают средним \(m\) выборочных средних, или общим средним \(\overline{\overline{X}} = (\overline{X}_1 + \overline{X}_2 + \cdots + \overline{X}_m)/m\). С другой стороны, для оценивания стандартного отклонения процесса \(\sigma\) было предложено несколько оценок. Мы используем некоторые из них, чтобы проиллюстрировать различные моменты. Пусть \(\hat{\sigma}\) обозначает выбранную по Фазе I оценку \(\sigma\); тогда контрольные границы для двусторонней \(k\)-сигма контрольной карты Шухарта \(\overline{X}\) для среднего в Фазе II в случае UU задаются как
\begin{equation}
	\begin{aligned}
		UCL &= \overline{\overline{X}} + k\frac{\hat{\sigma}}{\sqrt{n}} \\
		CL &= \overline{\overline{X}} \\
		LCL &= \overline{\overline{X}} - k\frac{\hat{\sigma}}{\sqrt{n}}.
	\end{aligned}
	\label{eq:3.31}
\end{equation}
Контрольные границы уравнения 3.31 могут называться границами с подстановкой оценок (plug-in limits), поскольку они имеют ту же форму, что и в случае K, но вместо известных значений параметров в состоянии IC используются оценки среднего и стандартного отклонения. Выбор константы для нанесения на карту \(k\) — интересный вопрос. Многие годы значение 3 использовалось по аналогии с 3-сигма картой в случае K. Сейчас понятно, что при границах с подстановкой оценок \(k = 3\) — не лучший выбор, так как он увеличивает число ложных тревог, если только не имеется очень большого количества данных Фазы I, которое может составлять сотни и даже тысячи наблюдений. Фактически в случае U выбор константы для нанесения на карту зависит от \(m\), \(n\) и заданной номинальной средней длины серии в состоянии IC. В случае K, напротив, константа для нанесения на карту зависит только от \(n\) и номинальной средней длины серии в состоянии IC. Подробнее об этом говорится ниже.

\subsubsection{Контрольные карты Шухарта для стандартного отклонения в случае U}

Как и в случае K, изменчивость процесса необходимо контролировать, чтобы убедиться, что он остаётся в состоянии IC. Контрольные границы Шухарта для среднего процесса в случае U зависят от стандартного отклонения процесса. Поэтому, если стандартное отклонение не остаётся в состоянии IC, контрольная карта для среднего не будет информативной и полезной. Значит, нужно контролировать дисперсию или стандартное отклонение с помощью контрольной карты.

Как отмечалось ранее, существует несколько возможных оценок \(\hat{\sigma}\), которые можно использовать для контроля стандартного отклонения \(\sigma\). Популярные варианты включают функции выборочного размаха (\(R\)), выборочного стандартного отклонения (\(S\)) или выборочной дисперсии (\(S^2\)). Обычно мы используем контрольную карту для контроля среднего процесса вместе с контрольной картой для контроля разброса процесса. Если разброс находится в состоянии IC, мы переходим к рассмотрению контрольной карты для среднего.

Например, карта Шухарта \(\overline{X}\) для среднего часто используется вместе с картой Шухарта \(R\) для разброса. Отметим, что для иллюстрации мы используем карту Шухарта \(R\) для \(\sigma\), хотя недавняя литература рекомендует использовать другую карту, например основанную на стандартном отклонении (см., к примеру, Mahmoud et al., 2010 и Epprecht, Loureiro and Chakraborti, 2016). Мы делаем так потому, что карта Шухарта \(R\) проста и продолжает использоваться в промышленности.

Пусть \(\overline{R} = (R_1 + R_2 + \cdots + R_m)/m\) обозначает средний размах. \(CL\) и контрольные границы карты Шухарта \(R\) для \(\sigma\) в случае неизвестных параметров задаются как (см. Montgomery, 2009, с. 229)
\begin{equation}
	\begin{aligned}
		UCL &= D_4\overline{R} \\
		CL &= \overline{R} \\
		LCL &= D_3\overline{R},
	\end{aligned}
	\label{eq:3.32}
\end{equation}
где
\[
	D_4 = 1 + \frac{3d_3(n)}{d_2(n)}, \qquad D_3 = 1 - \frac{3d_3(n)}{d_2(n)},
\]
а \(d_2(n)\) и \(d_3(n)\) — соответственно среднее и стандартное отклонение размаха в случайной выборке объёма \(n\) из стандартного нормального распределения. Мы приводим здесь эти формулы для полноты, поскольку во многих учебниках их нет:
\[
	d_2(n) = E(R) = \int_0^{\infty} r f_R(r)\,dr, \qquad
	d_3(n) = \sqrt{\operatorname{var}(R)} = \sqrt{E(R^2) - [E(R)]^2},
\]
где \(E(R^2) = \int_0^{\infty} w^2 f_R(w)\,dw\) и
\[
	f_R(w) = n(n - 1)\int_{-\infty}^{\infty} [\Phi(v) - \Phi(v - w)]^{n-2}\,\varphi(v - w)\varphi(v)\,dv,
	\quad w > 0,
\]
обозначает плотность распределения размаха выборки объёма \(n\) из стандартного нормального распределения с плотностью \(\varphi\) и функцией распределения \(\Phi\) (см., например, Gibbons and Chakraborti, 2010). Интегралы в этих выражениях можно вычислить численно или оценить в имитационном исследовании, чтобы найти константы \(d_2\) и \(d_3\). Обе они зависят от объёма выборки \(n\) и сведены в таблицу C приложения A. Отметим, что \(R/d_2(n)\) — несмещённая оценка \(\sigma\), так что \(d_2(n)\) называют константой снятия смещения (unbiasing constant).

Контрольные границы уравнения 3.32 имеют ту же форму, что и в случае K, но значения констант для нанесения на карту \(D_4\) и \(D_3\) другие. Однако контрольные границы уравнения 3.32 — это 3-сигма границы, и у них есть серьёзные недостатки, см., например, Diko et al. (2016), особенно при малых \(n\). Ниже в этом разделе мы рассмотрим некоторые альтернативы контрольным границам карты \(R\), основанные на вероятностных границах.

Используя аналогичную идею, но со стандартным отклонением выборки, а не с размахом, получаем \(CL\) и контрольные границы карты Шухарта \(S\) для \(\sigma\) в случае неизвестных параметров (см. Montgomery, 2009, с. 253):
\begin{equation}
	\begin{aligned}
		UCL &= B_4\overline{S} \\
		CL &= \overline{S} \\
		LCL &= B_3\overline{S},
	\end{aligned}
	\label{eq:3.33}
\end{equation}
где \(\overline{S} = \dfrac{S_1 + S_2 + \cdots + S_m}{m}\) — среднее стандартное отклонение, а константы задаются как
\[
	B_4 = 1 + \frac{3}{c_4(n)}\sqrt{1 - c_4^2(n)}, \qquad
	B_3 = 1 - \frac{3}{c_4(n)}\sqrt{1 - c_4^2(n)}
\]
и
\[
	c_4(n) = \frac{\sqrt{\dfrac{2}{n-1}}\;\Gamma\!\left(\dfrac{n}{2}\right)}{\Gamma\!\left(\dfrac{n-1}{2}\right)},
\]
где \(\Gamma\) обозначает гамма-функцию, причём \(\Gamma(n) = (n - 1)!\), когда \(n\) — целое положительное число. Отметим, что \(S/c_4(n)\) — также несмещённая оценка \(\sigma\). Константы \(B_3\) и \(B_4\) легко вычисляются по формулам и сведены в таблицу H приложения A.

Опять же отметим, что эти контрольные границы — 3-сигма границы, и, как и в случае карт \(R\), у них тоже есть недостатки, особенно при малых \(n\). Мы также рассмотрим некоторые альтернативы контрольным границам карты \(S\).

Хорошо известно, что вероятностные границы более подходят для стандартного отклонения, поскольку распределение оценки \(\hat{\sigma}\) (используемой в статистике для нанесения на карту) имеет правостороннюю асимметрию. Асимметрия выражена сильнее при меньшем объёме выборки Фазы II \(n\), что обычно и имеет место в SPC. Интуитивно нижняя контрольная граница карты для \(\sigma\) должна быть положительной (а не нулевой), чтобы карта могла обнаруживать как увеличения, так и уменьшения \(\sigma\). Увеличение \(\sigma\) указывает на деградацию процесса и, возможно, важнее, но уменьшение \(\sigma\) тоже может быть значимым, так как указывает на возможные улучшения процесса. Однако стандартные 3-сигма границы, используемые в случае K и основанные на нормальной аппроксимации распределения \(\hat{\sigma}\), могут привести к отрицательной \(LCL\), которую затем обычно заменяют нулём. Так и происходит при \(n \le 6\), что включает наиболее часто используемую на практике подгруппу, а именно \(n = 5\). Хотя 3-сигма границы, возможно, более популярны благодаря своей простоте, отметим, что многие авторы выступали за их применение со статистиками для нанесения на карту, имеющими асимметричное распределение. Montgomery (2009, с. 242) упоминает вероятностные границы и ссылается на Grant and Leavenworth (1980) в отношении констант, необходимых для построения таких контрольных границ. Однако отметим, что эти границы корректны и уместны в случае K, то есть в случае известных параметров. Для случая неизвестных параметров эти константы нужно скорректировать или адаптировать; это недавно рассматривалось в работе Diko et al. (2016). Далее мы кратко обсуждаем эту важную практическую проблему.

Для этого отметим, что вероятностные границы Фазы II для карты \(R\) (с оценкой \(R/d_2\)) и карты \(S\) (с оценкой \(S/c_4\)) в случае U задаются как
\[
	\begin{aligned}
		UCL &= D_4^*(m,n)\overline{R} \\
		CL &= \overline{R} \\
		LCL &= D_3^*(m,n)\overline{R}
	\end{aligned}
\]
и
\[
	\begin{aligned}
		UCL &= B_4^*(m,n)\overline{S} \\
		CL &= \overline{S} \\
		LCL &= B_3^*(m,n)\overline{S}
	\end{aligned}
\]
соответственно, где статистиками для нанесения на карту являются размах \(R_i\) и стандартное отклонение \(S_i\) \(i\)-й выборки Фазы II, а константы \((D_3^*, D_4^*)\) и \((B_3^*, B_4^*)\) получаются для заданных значений \(m\), \(n\) и номинального значения средней длины серии в состоянии IC, например 370. Эти константы приведены в таблице H приложения A. Некоторые подробности их вывода даны в приложении 3.3. Величина \(p\), показанная в этих таблицах, представляет суммарную вероятность хвостов за нижней и верхней контрольными границами при распределении статистики для нанесения на карту в состоянии IC. Она нужна при вычислении вероятностных границ.

Наконец, рассмотрим \(S_p = \sqrt{\sum_{i=1}^{m} S_i^2/m}\) как «объединённую» (pooled) оценку \(\sigma\), основанную на усреднении выборочных дисперсий Фазы I. Это оценка, рекомендованная в литературе (см., например, Mahmoud et al., 2010) благодаря её лучшим статистическим свойствам. \(CL\) и контрольные границы карты \(S\) Фазы II с \(S_p\), используемой в качестве оценки \(\sigma\) в случае неизвестных параметров, можно записать как
\begin{equation}
	\begin{aligned}
		UCL &= H_4^*(m,n)S_p \\
		CL &= S_p \\
		LCL &= H_3^*(m,n)S_p,
	\end{aligned}
	\label{eq:3.34}
\end{equation}
где статистикой для нанесения на карту является \(S_i\), а константы \(H_3^*(m,n)\) и \(H_4^*(m,n)\) приведены в таблице H приложения A для различных значений \(m\) и \(n = 5\) и 10, а также для номинальных средних длин серии в состоянии IC, равных 370 и 500 соответственно. Отметим, что в контрольных картах можно использовать и несмещённую версию \(S_p\), но видно, что константа снятия смещения приблизительно равна единице для типичных значений \(m\) и \(n\). Опять же, некоторые подробности вывода этих констант даны в приложении 3.3.

\subsection{Карта CUSUM для среднего в случае U}

В случае U константы для нанесения на карту \(k\) и \(H\) зависят от факторов \(m\), \(n\) и получаются для заданной номинальной средней длины серии в состоянии IC; см. таблицу 3.20. Выбор констант для нанесения на карту параметрической карты CUSUM не так прост, как в случае K. Это показано на примере в Разделе 3.9.

\subsection{Карта EWMA для среднего в случае U}

Для карты EWMA в случае U константы для нанесения на карту \(\lambda\) и \(L\) зависят от факторов \(m\), \(n\) и получаются для заданной номинальной средней длины серии в состоянии IC; см. таблицу 3.23. Как и в случае карты CUSUM, выбор констант для нанесения на карту параметрической карты EWMA не так прост, как в случае K. Это иллюстрируется примером в Разделе 3.9.

\section{Типы параметрических контрольных карт в случае U: иллюстративные примеры}

В этом разделе мы приводим некоторые примеры классических параметрических, то есть карт нормальной теории, контрольных карт в случае U.

\subsection{Контрольные карты для среднего}

\noindent\textbf{Пример 3.9. Контрольная карта Шухарта \(\overline{X}\) для среднего в Фазе II, когда и \(\mu\), и \(\sigma\) неизвестны}

Столбец (a) таблицы \ref{tab:3.13} содержит некоторые смоделированные данные из нормального распределения, которые представляют собой измерения, взятые из \(m = 25\) независимых случайных выборок, каждая объёма \(n = 5\), по некоторому типу пластины. Предположим, что это эталонные данные из процесса в состоянии IC, полученные после тщательного анализа Фазы I. Среднее и стандартное отклонение каждой эталонной выборки показаны в столбцах (b) и (c) таблицы \ref{tab:3.13} соответственно.

Оценкой среднего \(\mu\) является среднее выборочных средних Фазы I, или общее среднее

\[
\overline{X} = \frac{1.5119 + 1.4951 + \cdots + 1.5264}{25} = 1.5056.
\]

Как отмечалось выше, оценки \(\hat{\sigma}\), используемые в контрольных границах, основаны либо на (i) среднем выборочных размахов, \(\overline{R} = 0.3256\), либо на (ii) среднем выборочных стандартных отклонений, \(\overline{S} = 0.1316\), либо на объединённой оценке

\[
S_p = \sqrt{\frac{\sum_{i=1}^{m} S_i^2}{m}} = 0.1391.
\]

Отметим, что все три оценки \(\sigma\) близки друг к другу.

Контрольные карты Фазы II для среднего приведены ниже.

\begin{equation}
	\begin{aligned}
		UCL &= \overline{X} + k\frac{\hat{\sigma}}{\sqrt{n}} \\
		CL &= \overline{X} \\
		LCL &= \overline{X} - k\frac{\hat{\sigma}}{\sqrt{n}}.
	\end{aligned}
	\label{eq:3.35}
\end{equation}

Константы для нанесения на карту для карты Шухарта в случае U сведены в таблицу \ref{tab:3.14}. Для удобства та же таблица приведена как таблица G среди таблиц в Приложении A. Они основаны на работах Chakraborti (2000, 2006) с использованием безусловного метода. Из таблицы G видно, что константа для нанесения на карту для \(m = 25\), \(n = 5\) равна 2.9725 и 3.0574 для номинальных \(ARL_{\text{IC}} = 370\) и 500 соответственно. Так, например, контрольные границы Фазы II для среднего с использованием объединённой оценки \(S_p\) задаются как \(UCL = 1.5056 + 2.9725 \times \frac{0.1391}{\sqrt{5}} = 1.6905\), \(CL = 1.5056\) и \(LCL = 1.5056 - 2.9725 \times \frac{0.1391}{\sqrt{5}} = 1.3207\) соответственно.

Интересно рассмотреть элементы таблицы \ref{tab:3.14}. Очевидно, что требуемая константа для нанесения на карту не равна своему аналогу в случае K, если не используется большое число выборок Фазы I. Например, для \(ARL_{\text{IC}} = 370\) только тогда, когда \(m\) составляет примерно 100 или более, при \(n = 5\) константа для нанесения на карту становится ближе к 3 — значению для 3-сигма границы, которое подходит для случая известных параметров. Отметим также, что контрольные границы, определяемые безусловным \(ARL_{\text{IC}}\), по-видимому, немного уже, чем границы в случае K.

\begin{table}[htbp]
	\centering
	\caption{Измерения, использованные для построения контрольной карты Шухарта \(\overline{X}\).}
	\label{tab:3.13}
	\renewcommand{\arraystretch}{1.05} % Уменьшаем высоту строк, так как их много
	\begin{tabular}{c *{5}{r} r r r}
		\toprule
		Выборка & \multicolumn{5}{c}{(a)} & (b) & (c) & (d) \\
		\cmidrule(lr){2-6}
		\(i\) & \(X_1\) & \(X_2\) & \(X_3\) & \(X_4\) & \(X_5\) & \(\overline{X}_i\) & \(S_i\) & \(R_i\) \\
		\midrule
		1  & 1.3235 & 1.4128 & 1.6744 & 1.4573 & 1.6914 & 1.5119 & 0.1635 & 0.3679 \\
		2  & 1.4314 & 1.3592 & 1.6075 & 1.4666 & 1.6109 & 1.4951 & 0.1111 & 0.2517 \\
		3  & 1.4284 & 1.4871 & 1.4932 & 1.4324 & 1.5674 & 1.4817 & 0.0565 & 0.1390 \\
		4  & 1.5028 & 1.6352 & 1.3841 & 1.2831 & 1.5507 & 1.4712 & 0.1389 & 0.3521 \\
		5  & 1.5604 & 1.2735 & 1.5265 & 1.4363 & 1.6441 & 1.4882 & 0.1412 & 0.3706 \\
		6  & 1.5955 & 1.5451 & 1.3574 & 1.3281 & 1.4198 & 1.4492 & 0.1168 & 0.2674 \\
		7  & 1.6274 & 1.5064 & 1.8366 & 1.4177 & 1.5144 & 1.5805 & 0.1614 & 0.4189 \\
		8  & 1.4190 & 1.4303 & 1.6637 & 1.6067 & 1.5519 & 1.5343 & 0.1077 & 0.2447 \\
		9  & 1.3884 & 1.7277 & 1.5355 & 1.5176 & 1.3688 & 1.5076 & 0.1439 & 0.3589 \\
		10 & 1.4039 & 1.6697 & 1.5089 & 1.4627 & 1.5220 & 1.5134 & 0.0988 & 0.2658 \\
		11 & 1.4158 & 1.7667 & 1.4278 & 1.5928 & 1.4181 & 1.5242 & 0.1548 & 0.3509 \\
		12 & 1.5821 & 1.3355 & 1.5777 & 1.3908 & 1.7559 & 1.5284 & 0.1682 & 0.4204 \\
		13 & 1.2856 & 1.4106 & 1.4447 & 1.6398 & 1.1928 & 1.3947 & 0.1699 & 0.4470 \\
		14 & 1.4951 & 1.4036 & 1.5893 & 1.6458 & 1.4969 & 1.5261 & 0.0937 & 0.2422 \\
		15 & 1.3589 & 1.2863 & 1.5996 & 1.2497 & 1.5471 & 1.4083 & 0.1568 & 0.3499 \\
		16 & 1.5747 & 1.5301 & 1.5171 & 1.1839 & 1.8662 & 1.5344 & 0.2423 & 0.6823 \\
		17 & 1.3680 & 1.7269 & 1.3957 & 1.5014 & 1.4449 & 1.4874 & 0.1432 & 0.3589 \\
		18 & 1.4163 & 1.3864 & 1.3057 & 1.6210 & 1.5573 & 1.4573 & 0.1289 & 0.3153 \\
		19 & 1.5796 & 1.4185 & 1.6541 & 1.5116 & 1.7247 & 1.5777 & 0.1195 & 0.3062 \\
		20 & 1.7106 & 1.4412 & 1.2361 & 1.3820 & 1.7601 & 1.5060 & 0.2230 & 0.5240 \\
		21 & 1.4371 & 1.5051 & 1.3485 & 1.5670 & 1.4880 & 1.4691 & 0.0819 & 0.2185 \\
		22 & 1.4738 & 1.5936 & 1.6583 & 1.4973 & 1.4720 & 1.5390 & 0.0832 & 0.1863 \\
		23 & 1.5917 & 1.4333 & 1.5551 & 1.5295 & 1.6866 & 1.5592 & 0.0922 & 0.2533 \\
		24 & 1.6399 & 1.5243 & 1.5705 & 1.5563 & 1.5530 & 1.5688 & 0.0431 & 0.1156 \\
		25 & 1.5797 & 1.3663 & 1.6240 & 1.3732 & 1.6887 & 1.5264 & 0.1482 & 0.3224 \\
		\bottomrule
	\end{tabular}
\end{table}

\begin{table}[htbp]
	\centering
	\caption{Константы для нанесения на карту для контрольной карты Шухарта \(\overline{X}\) для среднего в случае UU для \(n = 5\), при изменении \(m\) и \(ARL_{\text{IC}} = 370\) и 500 соответственно, с использованием безусловного метода.} % CHECK: в оригинале «Case UU» — по-видимому, опечатка вместо «Case U»
	\label{tab:3.14}
	\begin{tabular}{ccc}
		\toprule
		\(m\) & \(ARL_{\text{IC}} = 370\) & \(ARL_{\text{IC}} = 500\) \\
		\midrule
		\bottomrule
	\end{tabular}
\end{table}

Предположим теперь, что по завершении анализа Фазы I становятся доступны 10 будущих (Фазы II) независимых случайных выборок. Выборочные средние показаны в столбце (b) таблицы \ref{tab:3.15}.

На рисунке \ref{fig-3.24} показаны средние этих 10 будущих выборочных средних, \(\overline{X}_i\), нанесённые на контрольную карту Шухарта \(\overline{X}\) вместе с обоими наборами контрольных границ: 3-сигма границами и скорректированными (исправленными) границами. Это отображение обычно называется контрольной картой Шухарта \(\overline{X}\) или картой «X-bar».

\begin{table}[htbp]
	\centering
	\caption{Данные Фазы II для контрольной карты Шухарта \(\overline{X}\) и статистики для нанесения на карту.}
	\label{tab:3.15}
	\renewcommand{\arraystretch}{1.05}
	\begin{tabular}{c *{5}{r} r r r}
		\toprule
		Выборка & \multicolumn{5}{c}{(a)} & (b) & (c) & (d) \\
		\cmidrule(lr){2-6}
		\(i\) & \(X_1\) & \(X_2\) & \(X_3\) & \(X_4\) & \(X_5\) & \(\overline{X}_i\) & \(R_i\) & \(S_i\) \\
		\midrule
		1  & 1.4231 & 1.3457 & 1.3211 & 1.4654 & 1.3845 & 1.3880 & 0.5099 & 0.0581 \\
		2  & 1.4213 & 1.4620 & 1.3394 & 1.6338 & 1.1239 & 1.3961 & 0.2970 & 0.1863 \\
		3  & 1.4885 & 1.5635 & 1.5446 & 1.5252 & 1.2665 & 1.4777 & 0.2576 & 0.1213 \\
		4  & 1.4158 & 1.5190 & 1.6259 & 1.3683 & 1.4117 & 1.4681 & 0.2999 & 0.1041 \\
		5  & 1.5129 & 1.3825 & 1.5784 & 1.3347 & 1.2785 & 1.4174 & 0.5099 & 0.1249 \\
		6  & 1.4660 & 1.3371 & 1.5378 & 1.6740 & 1.6637 & 1.5357 & 0.3369 & 0.1413 \\
		7  & 1.7342 & 1.2176 & 1.4747 & 1.6500 & 1.3623 & 1.4878 & 0.5166 & 0.2097 \\
		8  & 1.4146 & 1.2826 & 1.2620 & 1.3767 & 1.4036 & 1.3479 & 0.1526 & 0.0707 \\
		9  & 1.2262 & 1.4307 & 1.4523 & 1.5234 & 1.4574 & 1.4180 & 0.2971 & 0.1127 \\
		10 & 1.4625 & 1.4568 & 1.6827 & 1.4944 & 1.4810 & 1.5155 & 0.2259 & 0.0947 \\
		\bottomrule
	\end{tabular}
\end{table}

% CHECK: столбец (c) \(R_i\) в таблице 3.15 воспроизведён точно по книге, но значения в оригинале,
% по-видимому, сдвинуты на строку: для строки 1 напечатано 0.5099, тогда как размах
% 1.4654 - 1.3211 = 0.1443; значение 0.5099 соответствует размаху строки 2, 0.2970 — строки 3 и т.д.

\begin{figure}[!ht]
	\centering
	\includegraphics[width=0.7\linewidth]{"./images/fig-3.24.png"}
	\caption[]{Контрольная карта Шухарта \(\overline{X}\) для мониторинга среднего процесса, когда и среднее процесса, и стандартное отклонение неизвестны.}
	\label{fig-3.24}
\end{figure}

Отметим, что оценкой \(\sigma\), использованной на карте, была объединённая оценка \(S_p\). Как всегда, после построения контрольной карты принимается решение о состоянии процесса. На рисунке \ref{fig-3.24} ни одно из выборочных средних не выходит за пределы контрольных границ, и точки, по-видимому, не демонстрируют какого-либо неслучайного паттерна. Следовательно, мы заключаем, что процесс функционирует в состоянии статистической управляемости относительно среднего. Видно также, что два набора границ довольно близки друг к другу, причём скорректированные границы немного уже.
\noindent\textbf{Пример 3.12. Параметрическая (нормальной теории) карта CUSUM для среднего в случае неизвестных параметров при индивидуальных данных}

Столбец (a) таблицы \ref{tab:3.16} содержит 20 наблюдений, сгенерированных из нормально распределённого процесса. Предположим, что это эталонные данные из процесса в состоянии IC, полученные после тщательного анализа Фазы I. Предположим далее, что нас интересует обнаружение увеличения (сдвига) среднего на величину \(1\sigma\), что является сдвигом среднего размера. Поскольку \(\sigma\) неизвестно и его нужно оценить, мы используем скользящие размахи, то есть \(MR_i = |x_i - x_{i-1}|\), так как работаем с индивидуальными измерениями, то есть \(n = 1\). Скользящие размахи показаны в столбце (b) таблицы \ref{tab:3.16}. Тогда оценка \(\sigma\) задаётся как \(\overline{MR}/d_2\), где \(d_2 = 1.128\) (см. таблицу C приложения A при \(n = 2\), поскольку для вычисления скользящего размаха используются 2 наблюдения). Таким образом, находим \(\overline{MR} = \dfrac{1.46 + 1.30 + \cdots + 2.32}{19} = 1.550\), так что \(\hat{\sigma} = \dfrac{\overline{MR}}{d_2} = \dfrac{1.550}{1.128} = 1.374\). Нам также нужна оценка среднего, которая задаётся как \(\overline{X} = \dfrac{x_1 + x_2 + \cdots + x_{20}}{20} = 9.996\).

Предположим теперь, что по завершении анализа Фазы I становятся доступны 10 будущих независимых измерений.

В своей таблице 3.1 Hawkins (1993) рекомендует для сдвига среднего процесса примерно на \(1\sigma\) брать \(k = 0.5\) и \(h = 4.77\), так что \(K = k\hat{\sigma} = 0.5\hat{\sigma}\) и \(H = h\hat{\sigma} = 4.77\hat{\sigma}\) дают \(ARL_{\text{IC}} = 370\). Действуя таким образом, находим параметры проектирования \(K = k\hat{\sigma} = 0.5 \times 1.374 = 0.687\) и \(H = h\hat{\sigma} = 4.77 \times 1.374 = 6.554\) соответственно.

Среднее процесса мы проспективно контролируем в Фазе II по 10 выборкам, приведённым в таблице \ref{tab:3.17}. Из таблицы \ref{tab:3.18} видно, что константа для нанесения на карту \(h\) для \(m = 20\), \(n = 1\), \(k = 0.5\) и номинального \(ARL_0 = 370\) равна 4.1, а не 4.77, как в случае K. Здесь же отметим, что значения в таблице \ref{tab:3.18} получены имитационным моделированием и показывают влияние оценивания параметров на контрольные границы Фазы II. Правильное значение \(H\), таким образом, есть \(H = h\hat{\sigma} = 4.1 \times 1.374 = 5.634\). Для больших значений \(m\), порядка 100, значения \(h\), по-видимому, сходятся к своим значениям в случае K.

\begin{table}[htbp]
	\centering
	\caption{Измерения, использованные для построения параметрической карты CUSUM для индивидуальных данных, и статистика для нанесения на карту.}
	\label{tab:3.16}
	\begin{tabular}{crr}
		\toprule
		Период, \(i\) & (a) \(x_i\) & (b) \(MR_i\) \\
		\midrule
		1  & 9.45  &       \\
		2  & 7.99  & 1.46  \\
		3  & 9.29  & 1.30  \\
		4  & 11.66 & 2.37  \\
		5  & 12.16 & 0.50  \\
		6  & 10.18 & 1.98  \\
		7  & 8.04  & 2.14  \\
		8  & 11.46 & 3.42  \\
		9  & 9.20  & 2.26  \\
		10 & 10.34 & 1.14  \\
		11 & 9.03  & 1.31  \\
		12 & 11.47 & 2.44  \\
		13 & 10.51 & 0.96  \\
		14 & 9.40  & 1.11  \\
		15 & 10.08 & 0.68  \\
		16 & 9.37  & 0.71  \\
		17 & 10.62 & 1.25  \\
		18 & 10.31 & 0.31  \\
		19 & 8.52  & 1.79  \\
		20 & 10.84 & 2.32  \\
		\bottomrule
	\end{tabular}
\end{table}

\begin{table}[htbp]
	\centering
	\caption{Данные Фазы II для параметрической карты CUSUM.}
	\label{tab:3.17}
	\begin{tabular}{cr}
		\toprule
		Период, \(i\) & \(x_i\) \\
		\midrule
		1  & 10.90 \\
		2  & 9.33  \\
		3  & 12.29 \\
		4  & 11.50 \\
		5  & 10.60 \\
		6  & 11.08 \\
		7  & 10.38 \\
		8  & 11.62 \\
		9  & 11.31 \\
		10 & 10.52 \\
		\bottomrule
	\end{tabular}
\end{table}

\begin{table}[htbp]
	\centering
	\caption{Константа для нанесения на карту \(h\) для параметрической (нормальной теории) карты CUSUM для среднего в случае UU при \(n = 1\), \(k = 0.5\), для различных \(m\) и \(ARL_{\text{IC}} = 370\) и 500.}
	\label{tab:3.18}
	\begin{tabular}{ccc}
		\toprule
		\(m\) & \(ARL_{\text{IC}} = 370\) & \(ARL_{\text{IC}} = 500\) \\
		\midrule
		20   & 4.100 & 4.300 \\
		25   & 4.390 & 4.603 \\
		50   & 4.470 & 4.980 \\
		75   & 4.476 & 5.132 \\
		100  & 4.476 & 5.132 \\
		125  & 4.770 & 5.132 \\
		150  & 4.770 & 5.132 \\
		200  & 4.770 & 5.132 \\
		250  & 4.770 & 5.132 \\
		300  & 4.770 & 5.132 \\
		500  & 4.770 & 5.132 \\
		1000 & 4.770 & 5.120 \\
		\bottomrule
	\end{tabular}
\end{table}

Статистики для нанесения на карту CUSUM \(C_i^+\) и \(C_i^-\) показаны в столбцах (b) и (d) таблицы \ref{tab:3.19} соответственно, а соответствующие величины \(N^+\) и \(N^-\) приведены в столбцах (c) и (e) соответственно. Чтобы проиллюстрировать вычисления, рассмотрим период 1. Статистики для нанесения на карту вычисляются как в случае известных параметров, за исключением того, что вместо \(\mu_0\) мы используем \(\overline{X}\). Таким образом,
\[
\begin{aligned}
	C_1^+ &= \max[0, X_1 - \overline{X} - K + C_0^+] = \max[0, 10.9 - (9.996 + 0.687) + 0] \\
	      &= \max[0, 0.217] = 0.217
\end{aligned}
\]
и
\[
\begin{aligned}
	C_1^- &= \min[0, X_1 - \overline{X} + K + C_0^-] = \min[0, 10.9 - 9.996 + 0.687 + 0] \\
	      &= \min[0, 1.591] = 0.
\end{aligned}
\]

На рисунке \ref{fig-3.28} показаны статистики CUSUM \(C_i^+\) и \(C_i^-\) из столбцов (a) и (c) соответственно, которые нанесены на параметрическую карту CUSUM вместе с обоими наборами контрольных границ: с 3-сигма границами и со скорректированными на оценивание параметров границами.

Как всегда, после построения контрольной карты принимается решение о состоянии процесса. На рисунке \ref{fig-3.28} ни одна из статистик для нанесения на карту не выходит за пределы контрольных границ, и точки, по-видимому, не демонстрируют какого-либо неслучайного паттерна. Следовательно, мы заключаем, что процесс функционирует в состоянии статистической управляемости относительно среднего.

\begin{table}[htbp]
	\centering
	\caption{Данные Фазы II для параметрической карты CUSUM.}
	\label{tab:3.19}
	\begin{tabular}{crrrrr}
		\toprule
		& (a) & (b) & (c) & (d) & (e) \\
		Период, \(i\) & \(x_i\) & \(C_i^+\) & \(N^+\) & \(C_i^-\) & \(N^-\) \\
		\midrule
		& & 0 & & 0 & \\
		1  & 10.90 & 0.217 & 1  & 0 & 0 \\
		2  & 9.33  & 0.000 & 2  & 0 & 0 \\
		3  & 12.29 & 1.607 & 3  & 0 & 0 \\
		4  & 11.50 & 2.424 & 4  & 0 & 0 \\
		5  & 10.60 & 2.341 & 5  & 0 & 0 \\
		6  & 11.08 & 2.738 & 6  & 0 & 0 \\
		7  & 10.38 & 2.435 & 7  & 0 & 0 \\
		8  & 11.62 & 3.372 & 8  & 0 & 0 \\
		9  & 11.31 & 3.999 & 9  & 0 & 0 \\
		10 & 10.52 & 3.836 & 10 & 0 & 0 \\
		\bottomrule
	\end{tabular}
\end{table}

\begin{figure}[!ht]
	\centering
	\includegraphics[width=0.7\linewidth]{"./images/fig-3.28.png"}
	\caption[]{Параметрическая карта CUSUM для мониторинга среднего процесса, когда и среднее процесса, и стандартное отклонение неизвестны.}
	\label{fig-3.28}
\end{figure}

В этом примере мы проиллюстрировали параметрическую карту CUSUM для индивидуальных измерений. Однако если используются рациональные подгруппы объёма \(n > 1\), то в предыдущих уравнениях достаточно заменить \(X_i\) на \(\overline{X}_i\), а \(\sigma_0\) на \(\sigma_{\overline{X}} = \sigma_0/\sqrt{n}\). В этом случае оптимальная пара \(k\) и \(h\) может отличаться от пары, использованной в этом примере для индивидуальных данных. В таблице \ref{tab:3.18} константа для нанесения на карту \(h\) была дана для \(n = 1\). В таблице \ref{tab:3.20} мы приводим соответствующие значения для \(n = 5\).

\begin{table}[htbp]
	\centering
	\caption{Константа для нанесения на карту \(h\) для параметрической карты CUSUM для среднего в случае U при \(n = 5\), \(k = 0.5\), для различных \(m\) и \(ARL_{\text{IC}} = 370\) и 500 соответственно.}
	\label{tab:3.20}
	\begin{tabular}{ccc}
		\toprule
		\(m\) & \(ARL_{\text{IC}} = 370\) & \(ARL_{\text{IC}} = 500\) \\
		\midrule
		20   & 4.290 & 4.310 \\
		25   & 4.300 & 4.455 \\
		50   & 4.750 & 5.050 \\
		75   & 4.900 & 5.220 \\
		100  & 4.950 & 5.247 \\
		125  & 4.950 & 5.250 \\
		150  & 4.950 & 5.250 \\
		200  & 4.950 & 5.250 \\
		250  & 4.950 & 5.250 \\
		300  & 4.900 & 5.250 \\
		500  & 4.900 & 5.250 \\
		1000 & 4.900 & 5.250 \\
		\bottomrule
	\end{tabular}
\end{table}

\noindent\textbf{Пример 3.13. Параметрическая карта EWMA для среднего в случае неизвестных параметров}

Те же данные, которые использовались для иллюстрации параметрической карты CUSUM, используются и для иллюстрации параметрической карты EWMA в случае U. Для иллюстрации мы берём \(\lambda = 0.1\). Montgomery (2009; таблица 9.13, с. 425) рекомендует для \(\lambda = 0.1\) брать \(L = 2.703\), что даёт \(ARL_{\text{IC}} = 370\).

Среднее процесса мы проспективно контролируем в Фазе II по 10 выборкам, приведённым в таблице \ref{tab:3.21}. Будет видно, что из таблицы \ref{tab:3.23} константа для нанесения на карту \(L\) для \(m = 20\), \(n = 1\), \(\lambda = 0.1\) и номинального \(ARL_{\text{IC}} = 370\) равна 2.650, а не 2.703, как в случае K. Следует отметить, что значения в таблице \ref{tab:3.21} получены имитационным моделированием и показывают влияние оценивания параметров на контрольные границы Фазы II.

Статистика для нанесения на карту EWMA приведена в столбце (b) таблицы \ref{tab:3.22}, а зависящие от времени нижняя и верхняя контрольные границы — в столбцах (c) и (d) соответственно. В качестве начального значения берём \(Z_0 = \hat{\mu} = 9.996\). Чтобы проиллюстрировать вычисления, рассмотрим наблюдение номер 1. Первая статистика для нанесения на карту вычисляется так:
\[
	Z_1 = \lambda X_1 + (1 - \lambda)Z_0 = (0.1)(10.9) + (1 - 0.1)(9.996) = 10.086.
\]

\begin{table}[htbp]
	\centering
	\caption{Константа для нанесения на карту \(L\) для параметрической карты EWMA для среднего в случае UU при \(n = 1\), \(\lambda = 0.05\), \(0.10\), \(0.15\) и \(0.20\), для различных \(m\) и \(ARL_{\text{IC}} = 370\) и 500 соответственно.}
	\label{tab:3.21}
	\begin{tabular}{ccccc}
		\toprule
		\multirow{2}{*}{\(m\)} & \multicolumn{4}{c}{\(ARL_{\text{IC}} = 370\)} \\
		\cmidrule(lr){2-5}
		& \(\lambda = 0.05\) & \(\lambda = 0.10\) & \(\lambda = 0.15\) & \(\lambda = 0.20\) \\
		\midrule
		20   & 2.600 & 2.650 & 2.690 & 2.710 \\
		25   & 2.629 & 2.700 & 2.727 & 2.737 \\
		50   & 2.659 & 2.776 & 2.815 & 2.840 \\
		75   & 2.650 & 2.777 & 2.830 & 2.855 \\
		100  & 2.640 & 2.762 & 2.825 & 2.863 \\
		125  & 2.632 & 2.762 & 2.820 & 2.858 \\
		150  & 2.618 & 2.758 & 2.818 & 2.854 \\
		200  & 2.595 & 2.751 & 2.815 & 2.852 \\
		250  & 2.575 & 2.744 & 2.812 & 2.850 \\
		300  & 2.570 & 2.732 & 2.810 & 2.849 \\
		500  & 2.535 & 2.720 & 2.800 & 2.846 \\
		1000 & 2.515 & 2.710 & 2.787 & 2.840 \\
		\midrule
		\multirow{2}{*}{\(m\)} & \multicolumn{4}{c}{\(ARL_{\text{IC}} = 500\)} \\
		\cmidrule(lr){2-5}
		& \(\lambda = 0.05\) & \(\lambda = 0.10\) & \(\lambda = 0.15\) & \(\lambda = 0.20\) \\
		\midrule
		20   & 2.700 & 2.720 & 2.740 & 2.760 \\
		25   & 2.718 & 2.780 & 2.799 & 2.820 \\
		50   & 2.777 & 2.871 & 2.910 & 2.920 \\
		75   & 2.777 & 2.877 & 2.927 & 2.947 \\
		100  & 2.765 & 2.875 & 2.930 & 2.962 \\
		125  & 2.760 & 2.874 & 2.929 & 2.961 \\
		150  & 2.752 & 2.872 & 9.925 & 2.960 \\ % CHECK: в книге напечатано 9.925 — по-видимому, опечатка вместо 2.925
		200  & 2.730 & 2.867 & 2.920 & 2.959 \\
		250  & 2.715 & 2.852 & 2.919 & 2.958 \\
		300  & 2.710 & 2.850 & 2.912 & 2.958 \\
		500  & 2.665 & 2.835 & 2.905 & 2.955 \\
		1000 & 2.640 & 2.810 & 2.895 & 2.945 \\
		\bottomrule
	\end{tabular}
\end{table}

Подставляя \(\hat{\mu} = 9.996\), \(\hat{\sigma} = 1.374\), \(\lambda = 0.1\), \(L = 2.65\) и \(i = 1\) в уравнение 3.11, получаем \(CL\) и точные контрольные границы для данных (при \(i = 1\)) таблицы \ref{tab:3.22}:
\[
\begin{aligned}
	UCL &= 9.996 + 2.65(1.374)\sqrt{\frac{0.1}{2 - 0.1}\left(1 - (1 - 0.1)^{2(1)}\right)} = 10.360, \\
	CL &= 9.996, \\
	LCL &= 9.996 - 2.65(1.374)\sqrt{\frac{0.1}{2 - 0.1}\left(1 - (1 - 0.1)^{2(1)}\right)} = 9.632.
\end{aligned}
\]

\begin{table}[htbp]
	\centering
	\caption{Измерения, использованные для построения параметрической (нормальной теории) карты EWMA, статистика для нанесения на карту и контрольные границы.}
	\label{tab:3.22}
	\begin{tabular}{crrrr}
		\toprule
		& (a) & (b) & (c) & (d) \\
		Номер наблюдения \(i\) & \(X_i\) & \(Z_i\) & \(LCL_i\) & \(UCL_i\) \\
		\midrule
		&      & 9.996 &       &       \\
		1  & 10.90 & 10.086 & 9.632 & 10.360 \\
		2  & 9.33  & 10.011 & 9.506 & 10.486 \\
		3  & 12.29 & 10.239 & 9.424 & 10.568 \\
		4  & 11.50 & 10.365 & 9.366 & 10.626 \\
		5  & 10.60 & 10.388 & 9.322 & 10.670 \\
		6  & 11.08 & 10.458 & 9.288 & 10.704 \\
		7  & 10.38 & 10.450 & 9.262 & 10.730 \\
		8  & 11.62 & 10.567 & 9.242 & 10.750 \\
		9  & 11.31 & 10.641 & 9.226 & 10.766 \\
		10 & 10.52 & 10.629 & 9.213 & 10.779 \\
		\bottomrule
	\end{tabular}
\end{table}

Стационарные контрольные границы можно рассчитать с помощью уравнения 3.12; они равны \(LCL = 9.161\) и \(UCL = 10.831\) соответственно.

На рисунке \ref{fig-3.29} показаны статистики EWMA \(Z_i\) из столбца (b) таблицы \ref{tab:3.22}, нанесённые на параметрическую карту EWMA вместе с двумя наборами зависящих от времени контрольных границ: 3-сигма границами (сплошная линия) и скорректированными на оценивание параметров границами (пунктирная линия).

\begin{figure}[!ht]
	\centering
	\includegraphics[width=0.7\linewidth]{"./images/fig-3.29.png"}
	\caption[]{Параметрическая карта EWMA для мониторинга среднего процесса с 3-сигма и зависящими от времени контрольными границами, когда и среднее процесса, и стандартное отклонение неизвестны.}
	\label{fig-3.29}
\end{figure}

\begin{table}[htbp]
	\centering
	\caption{Константа для нанесения на карту \(L\) для параметрической (нормальной теории) карты EWMA для среднего в случае U при \(n = 5\), \(\lambda = 0.05\), \(0.10\), \(0.15\) и \(0.20\), для различных \(m\) и \(ARL_{\text{IC}} = 370\) и 500 соответственно.}
	\label{tab:3.23}
	\begin{tabular}{ccccc}
		\toprule
		\multirow{2}{*}{\(m\)} & \multicolumn{4}{c}{\(ARL_{\text{IC}} = 370\)} \\
		\cmidrule(lr){2-5}
		& \(\lambda = 0.05\) & \(\lambda = 0.10\) & \(\lambda = 0.15\) & \(\lambda = 0.20\) \\
		\midrule
		20   & 2.600 & 2.660 & 2.675 & 2.690 \\
		25   & 2.660 & 2.720 & 2.721 & 2.740 \\
		50   & 2.740 & 2.820 & 2.845 & 2.850 \\
		75   & 2.737 & 2.840 & 2.879 & 2.880 \\
		100  & 2.735 & 2.835 & 2.877 & 2.900 \\
		125  & 2.720 & 2.835 & 2.875 & 2.900 \\
		150  & 2.700 & 2.830 & 2.873 & 2.900 \\
		200  & 2.680 & 2.810 & 2.870 & 2.895 \\
		250  & 2.660 & 2.805 & 2.862 & 2.890 \\
		300  & 2.650 & 2.795 & 2.855 & 2.887 \\
		500  & 2.600 & 2.760 & 2.845 & 2.882 \\
		1000 & 2.550 & 2.730 & 2.820 & 2.860 \\
		\midrule
		\multirow{2}{*}{\(m\)} & \multicolumn{4}{c}{\(ARL_{\text{IC}} = 500\)} \\
		\cmidrule(lr){2-5}
		& \(\lambda = 0.05\) & \(\lambda = 0.10\) & \(\lambda = 0.15\) & \(\lambda = 0.20\) \\
		\midrule
		20   & 2.700 & 2.740 & 2.750 & 2.760 \\
		25   & 2.761 & 2.942 & 2.810 & 2.810 \\
		50   & 2.840 & 2.942 & 2.960 & 2.950 \\
		75   & 2.850 & 2.941 & 2.985 & 2.980 \\
		100  & 2.860 & 2.940 & 2.985 & 2.990 \\
		125  & 2.850 & 2.937 & 2.985 & 3.000 \\
		150  & 2.835 & 2.935 & 2.985 & 3.000 \\
		200  & 2.810 & 2.930 & 2.977 & 3.000 \\
		250  & 2.795 & 2.915 & 2.960 & 3.000 \\
		300  & 2.790 & 2.900 & 2.955 & 3.000 \\
		500  & 2.725 & 2.880 & 2.932 & 2.980 \\
		1000 & 2.690 & 2.850 & 2.920 & 2.970 \\
		\bottomrule
	\end{tabular}
\end{table}

Как всегда, после построения контрольной карты принимается решение о состоянии процесса. На рисунке \ref{fig-3.29} ни одна из статистик для нанесения на карту не выходит за пределы контрольных границ, и точки, по-видимому, не демонстрируют какого-либо неслучайного паттерна. Следовательно, мы заключаем, что процесс функционирует в состоянии статистической управляемости относительно среднего.
В этом примере мы проиллюстрировали параметрическую карту EWMA для индивидуальных измерений. Однако если используются рациональные подгруппы объёма \(n > 1\), то в предыдущих уравнениях достаточно заменить \(X_i\) на \(\overline{X}_i\), а \(\sigma_0\) на \(\sigma_{\overline{X}} = \sigma_0/\sqrt{n}\). В таблице \ref{tab:3.21} константа для нанесения на карту \(L\) была дана для \(n = 1\). Здесь мы приводим таблицу для \(n = 5\).

Как обсуждалось для случая K, важно помнить, что при мониторинге среднего процесса контролируются и среднее, и стандартное отклонение, поскольку стандартное отклонение входит в контрольные границы для среднего и, значит, должно находиться в состоянии IC, когда его оценивают и используют при расчёте контрольных границ. Это хорошо известный на практике момент, но влияние такого «совместного» мониторинга не всегда явно обсуждается и понимается.

\subsection{Карты для стандартного отклонения}

Здесь мы иллюстрируем карты \(R\) и \(S\), которые можно использовать для мониторинга разброса и стандартного отклонения соответственно. Начнём с примера карты Шухарта \(R\).

\noindent\textbf{Пример 3.10. Контрольная карта Шухарта \(R\) фазы II для стандартного отклонения в случае неизвестных параметров}

Те же данные, которые использовались для иллюстрации карты Шухарта \(\overline{X}\), теперь используются для иллюстрации карты Шухарта \(R\) в случае U. Напомним, что столбец (a) таблицы \ref{tab:3.15} содержит измерения, взятые из 25 независимых выборок Фазы I по пластинам, каждая объёма \((n = 5)\), из нормального распределения. % CHECK: в таблице 3.15 приведены 10 выборок фазы II, а не 25 выборок фазы I
Размах каждой выборки показан в столбце (d) таблицы \ref{tab:3.15}. % CHECK: в таблице 3.15 размах R находится в столбце (c), а не (d)
Для этих данных \(\overline{R} = 0.325\). Сначала для \(n = 5\) находим из таблицы C в Приложении A, что \(D_3 = 0\) и \(D_4 = 2.114\). Таким образом, при \(\overline{R} = 0.325\) 3-сигма контрольные границы карты \(R\) фазы II в случае U для \(\sigma\) задаются как \(LCL = 0\) и \(UCL = 2.114 \times 0.325 = 0.688\) при \(CL = 0.325\). Однако эти контрольные границы не учитывают должным образом оценивание параметров, и не следует ожидать, что они будут точными, если только нет большого количества данных Фазы I. Напротив, из таблицы H в Приложении A для \(m = 25\), \(n = 5\) и \(ARL_0 = 370\) находим, что \(D_3^*(25, 5) = 0.16603\) и \(D_4^*(25, 5) = 2.32788\). Следовательно, вероятностные контрольные границы карты Шухарта \(R\) фазы II для \(\sigma\), использующие оценку \(\overline{R}/d_2\), задаются как \(LCL = D_3^*(m, n)\overline{R} = (0.16603)(0.325) = 0.054\) и \(UCL = D_4^*(m, n)\overline{R} = (2.32788)(0.325) = 0.757\) при \(CL = \overline{R} = 0.325\). Отметим, что если \(LCL\) для 3-сигма карты \(R\) должна быть сброшена до 0, поскольку она отрицательна, то \(LCL\) для карты с вероятностными границами положительна, и такой корректировки не требуется. На рисунке \ref{fig-3.25} показаны выборочные размахи \(R_i\) из столбца (b), которые нанесены на контрольную карту Шухарта \(R\) вместе с контрольными границами. % CHECK: в таблице 3.15 размах R находится в столбце (c), а не (b)

На рисунке \ref{fig-3.25} все точки отображаются между контрольными границами, и в данных нет аномальных паттернов. Следовательно, мы заключаем, что процесс функционирует в статистически управляемом состоянии в отношении разброса. Отметим, что в случае неизвестных параметров скорректированные контрольные границы шире, что и ожидается, поскольку дополнительная вариация вносится через оценки параметров.

\begin{figure}[!ht]
	\centering
	\includegraphics[width=0.7\linewidth]{"./images/fig-3.25.png"}
	\caption[]{Контрольная карта Шухарта \(R\) для мониторинга разброса процесса при неизвестных параметрах.}
	\label{fig-3.25}
\end{figure}

\noindent\textbf{Пример 3.11. Контрольные карты Шухарта \(S\) фазы II для стандартного отклонения в случае неизвестных параметров с использованием оценки \(\overline{S}\)}

Те же данные, которые использовались для иллюстрации карты Шухарта \(\overline{X}\) и карты \(R\), теперь используются для иллюстрации карты Шухарта \(S\) в случае U. Напомним, что столбец (a) таблицы \ref{tab:3.15} содержит измерения, взятые из 25 независимых выборок Фазы I по пластинам, каждая объёма \((n = 5)\), из нормального распределения. Стандартное отклонение каждой выборки фазы II показано в столбце (c) таблицы \ref{tab:3.15}. % CHECK: в таблице 3.15 стандартное отклонение S находится в столбце (d), а не (c)
Для стандартной 3-сигма контрольной карты Шухарта \(S\) в случае U, использующей оценку \(\overline{S}/c_4\), контрольные границы получаются (см. таблицу C в Приложении A) следующим образом. Для \(n = 5\) находим \(B_3 = 0\) и \(B_4 = 2.089\). Следовательно, при \(\overline{S} = 0.1316\) контрольные границы задаются как \(LCL = 0\) и \(UCL = 2.089 \times 0.1316 = 0.2749\) при \(CL = 0.1316\). Опять же, эти границы не рекомендуется использовать на практике, если только нет большого объёма данных Фазы I. Вместо этого из таблицы H в Приложении A находим для \(m = 25\), \(n = 5\) и номинального \(ARL_{\text{IC}} = 370\), что \(B_3^*(25, 5) = 0.1685\) и \(B_4^*(25, 5) = 2.2587\). Следовательно, контрольные границы карты Шухарта \(S\) фазы II с вероятностными границами, использующей оценку \(\overline{S}/c_4\), равны \(LCL = (0.1685)(0.1391) = 0.0234\) и \(UCL = (2.2587)(0.1391) = 0.3142\) при \(CL = 0.1391\). % CHECK: в книге непоследовательно указаны \(\overline{S} = 0.1316\) (UCL = 0.2749) и 0.1391 (UCL = 0.3142); на рисунке 3.26 подписаны UCL(3-sigma) = 0.291 и CL = 0.139, что соответствует 0.1391
Отметим снова, что \(LCL\) карты с вероятностными границами положительна, и поэтому эта карта \(S\) способна обнаруживать как увеличения, так и уменьшения стандартного отклонения.

На рисунке \ref{fig-3.26} все точки отображаются между контрольными границами, и в данных нет аномальных паттернов. Следовательно, мы заключаем, что процесс функционирует в статистически управляемом состоянии в отношении вариации. Отметим снова, что в случае неизвестных параметров контрольные границы шире, что и ожидается, поскольку дополнительная вариация вносится через оценки параметров.

\begin{figure}[!ht]
	\centering
	\includegraphics[width=0.7\linewidth]{"./images/fig-3.26.png"}
	\caption[]{Контрольная карта Шухарта \(S\) с оценкой \(\overline{S}\) для мониторинга стандартного отклонения процесса при неизвестных параметрах.}
	\label{fig-3.26}
\end{figure}

\subsubsection{Использование оценки \(S_p\)}
Далее мы иллюстрируем карту Шухарта \(S\) с вероятностными границами, в которой в качестве оценки \(\sigma\) используется \(S_p\). Отметим, что для этого набора данных оценка по Фазе I равна \(S_p = \sqrt{\sum_{i=1}^{m} S_i^2/m} = 0.1391\). Из таблицы H в Приложении A находим необходимые константы для нанесения на карту: \(H_3^*(25, 5) = 0.1581\) и \(H_4^*(25, 5) = 2.1240\) для номинального \(ARL_{\text{IC}} = 370\). Таким образом, карта Шухарта \(S\) с вероятностными границами задаётся как \(LCL = H_3^*(m, n)S_p = (0.1581)(0.1391) = 0.0220\) и \(UCL = H_4^*(m, n)S_p = (2.1240)(0.1391) = 0.2955\) при \(CL = 0.1391\). Отметим, что эта карта немного уже, чем карта Шухарта \(S\) с вероятностными границами, использующая оценку \(\overline{S}/c_4\), показанная на рисунке \ref{fig-3.26}. Это ожидаемо, поскольку \(S_p\) является более эффективной оценкой, чем \(\overline{S}/c_4\).

На рисунке \ref{fig-3.27} все точки отображаются между контрольными границами, и в данных нет аномальных паттернов. Следовательно, мы заключаем, что процесс функционирует в статистически управляемом состоянии в отношении вариации.

\begin{figure}[!ht]
	\centering
	\includegraphics[width=0.7\linewidth]{"./images/fig-3.27.png"}
	\caption[]{Контрольная карта Шухарта \(S\) с оценкой \(S_p\) для мониторинга стандартного отклонения процесса при неизвестных параметрах.}
	\label{fig-3.27}
\end{figure}

\section{Распределение длины серии в случае неизвестных параметров (случай U)}

\subsection{Методы расчёта распределения длины серии и его свойств: метод обусловливания/деобусловливания}

\subsubsection{Карта Шухарта для среднего в случае U}

Метод обусловливания-деобусловливания (CUC), впервые явно введённый и применённый в Chakraborti (2000), объясняется здесь на примере параметрической карты Шухарта \(\overline{X}\) для среднего при предположении о нормальном распределении. Это развитие важно, поскольку случай «неизвестных стандартов», то есть случай U, часто встречается на практике, а карта Шухарта \(\overline{X}\) для среднего — одна из самых популярных карт, применяемых на практике. Перед подробным изложением дадим краткие предварительные сведения. Напомним, что в случае «известных стандартов», то есть когда параметры процесса известны (случай K), сигнальные события взаимно независимы, так что распределение длины серии является геометрическим с вероятностью успеха (а в SPC это сигнал), равной, скажем, некоторому \(\theta\). Этот результат полностью характеризует эффективность карты Шухарта в случае K, так что все характеристики эффективности карты можно получить из свойств распределения GEO(\(\theta\)). Таким образом, математическое ожидание распределения длины серии равно обратной величине \(\theta\). Для распределения длины серии в состоянии IC \(ARL_{\text{IC}}\) просто равна обратной величине FAR, то есть вероятности сигнала \(\theta\), когда процесс находится в состоянии IC. Вероятность сигнала в состоянии IC обозначается через \(\alpha\), и поэтому \(ARL_{\text{IC}} = 1/\alpha\), так что задание FAR задаёт \(ARL_{\text{IC}}\), и наоборот. Эта простая связь упрощает понимание эффективности карты Шухарта \(\overline{X}\) в случае K. Напротив, когда параметры процесса неизвестны и их нужно оценить, чтобы построить контрольные границы до начала мониторинга в Фазе II, сигнальные события уже не независимы, так что распределение длины серии уже не является геометрическим. Как следствие, например, ARL уже не является обратной величиной вероятности сигнала. В этом состоит фундаментальное концептуальное различие между случаем K и случаем U, которое имеет важные последствия, поскольку у некоторых может возникнуть соблазн использовать результаты случая K для проектирования контрольных карт даже тогда, когда лежащие в основе параметры процесса неизвестны. Практиков предостерегают от такой практики (см., например, Quesenberry, 1993), и поэтому мы обсуждаем, как правильно справляться с этой ситуацией. Для этого мы используем важный инструмент под названием «метод обусловливания-деобусловливания», предложенный в Chakraborti (2000) и широко применяемый в литературе многими исследователями. Отметим, что Chen (1997) и другие также использовали сходные идеи.

При мониторинге нормального среднего возможны четыре случая. Во-первых, простейший случай — когда известны или заданы и среднее, и стандартное отклонение. Это называется случаем известных стандартов и ранее обозначалось как случай K. В дальнейшем он называется случаем KK, причём каждая буква обозначает статус/предположение о соответствующем параметре: первая — о среднем, вторая — о стандартном отклонении. Далее, возможно наиболее распространённый случай — когда среднее процесса задано или известно (K), а стандартное отклонение процесса неизвестно (U). Он будет называться случаем KU. Два других случая — соответственно случай UK, когда среднее неизвестно, а стандартное отклонение известно, и, наконец, важный случай UU, когда неизвестны и среднее процесса, и стандартное отклонение процесса. Признаем, что случай UK более академичен и, возможно, менее практичен, поэтому мы оставляем его анализ читателю в качестве упражнения.

В каждом случае ключевые величины — это сигнальное событие, вероятность сигнального события и распределение длины серии. Как только распределение длины серии получено, связанные с ним характеристики эффективности карты, такие как среднее (математическое ожидание), стандартное отклонение и т.д., находятся легко.

Начнём для простоты со случая KK, чтобы задать контекст. В этом случае двусторонние контрольные границы Шухарта задаются как
\[
	LCL = \mu_0 - z_{\alpha/2}\frac{\sigma_0}{\sqrt{n}} \quad\text{и}\quad
	UCL = \mu_0 + z_{\alpha/2}\frac{\sigma_0}{\sqrt{n}},
\]
где \(\mu_0\) и \(\sigma_0\) обозначают известные среднее и стандартное отклонение процесса в состоянии IC соответственно. Процесс объявляется вышедшим из-под контроля (OOC) или, что эквивалентно, контрольная карта подаёт сигнал, когда статистика для нанесения на карту оказывается вне контрольных границ. Таким образом, в случае KK сигнальным событием является
\[
	\overline{X}_i \ge UCL \quad\text{или}\quad \overline{X}_i \le LCL, \quad i = 1, 2, \ldots
\]
или, что эквивалентно, если
\[
	\frac{\sqrt{n}(\overline{X}_i - \mu_0)}{\sigma_0} \ge z_{\alpha/2}
	\quad\text{или}\quad
	\frac{\sqrt{n}(\overline{X}_i - \mu_0)}{\sigma_0} \le -z_{\alpha/2}, \quad i = 1, 2, \ldots
\]
Следовательно, вероятность сигнала равна
\[
\begin{aligned}
	\theta &= P\left(\frac{\sqrt{n}(\overline{X}_i - \mu_0)}{\sigma_0} \ge z_{\alpha/2}\right)
	       + P\left(\frac{\sqrt{n}(\overline{X}_i - \mu_0)}{\sigma_0} \le -z_{\alpha/2}\right) \\
	       &= 1 - \Phi(z_{\alpha/2}) + \Phi(-z_{\alpha/2}).
\end{aligned}
\]
Таким образом, в случае KK распределение длины серии контрольной карты Шухарта \(\overline{X}\) является геометрическим с вероятностью успеха \(\theta\), то есть GEO(\(\theta\)). Это следует из того, что для случая KK выполняются четыре условия геометрического распределения, а именно: эксперименты/испытания должны быть независимыми испытаниями Бернулли с постоянной вероятностью успеха, где под испытанием/экспериментом здесь понимается сравнение выборочного среднего с контрольными границами. Кроме того, на каждом шаге мы наблюдаем испытание Бернулли, которое является либо «сигналом», когда \(\overline{X} \ge UCL\) или \(\overline{X}_i \le LCL\), что можно классифицировать как успех (\(S\)), либо «отсутствием сигнала», когда \(LCL < \overline{X}_i < UCL\), что можно классифицировать как неудачу (\(F\)). Более того, вероятность успеха равна \(P(S) = \theta\), а испытания в разные моменты времени независимы, поскольку выборки берутся независимо в каждый момент времени. Наконец, длина серии — это случайная величина, число испытаний до первого успеха в этой последовательности независимых испытаний Бернулли. Таким образом, распределение длины серии следует геометрическому распределению с вероятностью успеха \(P(S) = \theta\).

Поскольку распределение длины серии геометрическое, оно полностью характеризуется вероятностью успеха \(\theta\), и потому все свойства распределения, такие как среднее (ARL), стандартное отклонение (SDRL) и медианная длина серии (MRL) и т.д., можно найти через \(\theta\); например, ARL и SDRL задаются как
\[
	ARL = \frac{1}{\theta} \quad\text{и}\quad SDRL = \frac{\theta}{(1 - \theta)^2}
\]
соответственно. Как отмечалось ранее в Главе 2, \(100p\)-й процентиль \(Q(p)\) можно вычислить по формуле
\[
	Q(p) = \inf\{j \in \mathbb{Z} : P(N \le j) \ge p\}
	     = \inf\left\{j \in \mathbb{Z} : j \ge \frac{\ln(1 - p)}{\ln(1 - \theta)}\right\}, \quad j \ge 1,
\]
где \(\inf\) обозначает инфимум, а \(0 < p < 1\) — порядок процентиля. Таким образом, \(Q(p)\) — наименьшее целое число, не меньшее \(\dfrac{\ln(1 - p)}{\ln(1 - \theta)}\). Например, для медианной длины серии MRL имеем \(p = 0.5\), и потому
\[
	MRL = Q(0.5) = \inf\left\{j \in \mathbb{Z} : j \ge \frac{\ln(0.5)}{\ln(1 - \theta)}\right\}
	    = \inf\left\{j \in \mathbb{Z} : j \ge \frac{-\ln(2)}{\ln(1 - \theta)}\right\}.
\]

Далее рассмотрим случай, когда среднее известно или задано, а дисперсия неизвестна, то есть случай KU. Такая ситуация может возникать в задачах, где среднее задано некоторым производственным или нормативным ограничением, а стандартное отклонение процесса неизвестно. В этой ситуации есть более одного способа начать мониторинг процесса. Обычно неизвестный параметр (здесь — дисперсию) сначала оценивают по ретроспективному (Фаза I) анализу данных, уже имеющихся по процессу, когда он считался находящимся в состоянии IC. Этот анализ называется анализом Фазы I, а такие данные — эталонными данными, которые состоят из выборки фиксированного объёма. Другие варианты могут включать оценивание параметров «на ходу» с помощью так называемых самозапускающихся и последовательных схем выборки, но здесь мы сосредоточимся на более традиционном подходе с использованием эталонной выборки фиксированного объёма. Контрольные карты используются в Фазе I, чтобы привести процесс под контроль, но построение контрольных границ там не такое, как в Фазе II. Контрольные карты для анализа Фазы I мы обсуждаем позже, в Главе 5.

В Фазе II полученные в Фазе I контрольные границы применяются для постоянного мониторинга состояния процесса. Предположим, что в Фазе I процесс следует нормальному распределению, так что \(X_{ij} \sim N(\mu_0, \sigma^2)\) для \(i = 1, 2, 3, \ldots, m\) и \(j = 1, 2, 3, \ldots, n\). Кроме того, предположим, что неизвестное стандартное отклонение процесса в состоянии IC \(\sigma_0\) мы оцениваем точечной оценкой \(S = \sqrt{S^2} = \sqrt{\sum_{i=1}^{m}(X_i - \mu_0)^2/m}\). Двусторонние контрольные границы в случае KU получаются подстановкой (plug-in) точечной оценки в контрольные границы случая KK:
\[
	\widehat{LCL} = \mu_0 - z_{\alpha/2}\frac{S}{\sqrt{n}} \quad\text{и}\quad
	\widehat{UCL} = \mu_0 + z_{\alpha/2}\frac{S}{\sqrt{n}}.
\]
Заметим, что в этом случае сигналы и отсутствие сигналов относятся к событиям в Фазе II, то есть к выборке \((m + 1)\) и последующим. Предположим, что выборки Фазы II следуют распределению \(N(\mu_1, \sigma_0^2)\). Тогда сигнальное событие в Фазе II можно записать как
\[
	\overline{X}_i \le \widehat{LCL} \quad\text{или}\quad \overline{X}_i \ge \widehat{UCL}
	\quad\text{для } i = m + 1, m + 2, m + 3, \ldots,
\]
что можно переписать как
\[
	\frac{\sqrt{n}|\overline{X}_i - \mu_0|}{S} \ge z_{\alpha/2}
	\quad\text{для } i = m + 1, m + 2, m + 3, \ldots
\]
Вероятность сигнального события в случае KU вычисляется в два шага. Сначала вероятность сигнала вычисляется условно при заданной оценке параметра \(S\) (или \(S^2\)) из Фазы I, которая является случайной величиной. Поэтому эта вероятность называется условной вероятностью сигнала и играет ключевую роль в последующих построениях. Отметим, что при подстановке оценки стандартного отклонения в контрольные границы случая KK для построения границ случая UU можно было бы использовать константу для нанесения на карту из \(t\)-распределения с \(n - 1\) степенями свободы, особенно поскольку \(n\) обычно мало. Такая идея рассматривалась в литературе, но мы здесь этих деталей не касаемся.

Пусть \(\mu_1 = \mu_0 + \delta\sigma_1\), где величина \(\delta\) характеризует, есть ли сдвиг среднего, и потому такая запись позволяет рассматривать как случай IC (\(\delta = 0\)), так и случай OOC (\(\delta \ne 0\)). Тогда условная вероятность сигнала при заданном \(S\) в Фазе II равна
\[
	P\left(\frac{\sqrt{n}|\overline{X}_i - \mu_0|}{S} \ge z_{\alpha/2} \;\middle|\; \mu_1 = \mu_0 + \delta\sigma_1\right)
	\quad\text{для } i = m + 1, m + 2, m + 3, \ldots
\]
Эту вероятность можно переписать как
\[
	= 1 - P\left(\mu_0 - z_{\alpha/2}\frac{S}{\sqrt{n}} < \overline{X}_i < \mu_0 + z_{\alpha/2}\frac{S}{\sqrt{n}}
	\;\middle|\; \mu_1 = \mu_0 + \delta\sigma_1\right).
\]
Теперь, предполагая, что наряду с лежащими в основе нормальными распределениями дисперсии в Фазе I и Фазе II не меняются (то есть сдвиг ожидается только по среднему при \(\sigma_1 = \sigma_0 = \sigma\)), последнюю вероятность можно переписать как
\[
\begin{aligned}
	&= 1 - P\left(
	\frac{\mu_0 - z_{\alpha/2}S/\sqrt{n} - (\mu_0 + \delta\sigma)}{\sigma/\sqrt{n}}
	< \frac{\overline{X}_i - \mu_1}{\sigma/\sqrt{n}} <
	\frac{\mu_0 + z_{\alpha/2}S/\sqrt{n} - (\mu_0 + \delta\sigma)}{\sigma/\sqrt{n}}
	\;\middle|\; S\right) \\
	&= 1 - P\left(-z_{\alpha/2}\frac{S}{\sigma} - \delta\sqrt{n} < Z
	< z_{\alpha/2}\frac{S}{\sigma} - \delta\sqrt{n} \;\middle|\; S\right) \\
	&= 1 - P\left(-z_{\alpha/2}\frac{1}{\sqrt{v}}\sqrt{\frac{vS^2}{\sigma^2}} - \delta\sqrt{n} < Z
	< z_{\alpha/2}\frac{1}{\sqrt{v}}\sqrt{\frac{vS^2}{\sigma^2}} - \delta\sqrt{n} \;\middle|\; S\right) \\
	&= \beta(\alpha, \delta, m, n, S),\ \text{скажем}.
\end{aligned}
\]
Наконец, полагая \(Y = \dfrac{vS^2}{\sigma^2}\), что следует распределению хи-квадрат с \(\nu = mn\) степенями свободы, можно переписать последнее выражение как
\[
	\beta(\alpha, \delta, m, n, y) = 1 - P\left(-z_{\alpha/2}\sqrt{\frac{Y}{v}} - \delta\sqrt{n} < Z
	< z_{\alpha/2}\sqrt{\frac{Y}{v}} - \delta\sqrt{n} \;\middle|\; Y\right).
\]
Следовательно, поскольку \(Z\) имеет распределение \(N(0,1)\), условную вероятность сигнала можно удобно записать как
\[
\begin{aligned}
	P_C(\text{сигнал}) &= \beta(\alpha, \delta, m, n, Y) \\
	&= 1 - \Phi\left(z_{\alpha/2}\sqrt{\frac{Y}{v}} - \delta\sqrt{n}\right)
	   - \Phi\left(-z_{\alpha/2}\sqrt{\frac{Y}{v}} - \delta\sqrt{n}\right).
\end{aligned}
\]
Условная вероятность явно показывает, что эффективность карты будет зависеть от значения случайной величины \(Y\), которая соответствует оценке \(\sigma\), использованной и полученной по данным Фазы I. Как отмечалось ранее, условная вероятность сигнала играет важную роль в анализе эффективности контрольной карты при оценённых параметрах. Подробнее об этом мы скажем позже. Однако, поскольку эта вероятность является функцией случайной величины \(Y\), она сама является случайной величиной со своим собственным распределением вероятностей. В этом и состоит главный эффект оценивания параметров на эффективность контрольных карт, и для концептуального понимания он принципиален. Таким образом, когда контрольные границы используют оценку(и) параметров, в Фазе II сигнальные события зависимы, потому что в контрольных границах используется одна и та же оценка стандартного отклонения по Фазе I и каждая выборка Фазы II сравнивается с одними и теми же контрольными границами, содержащими эту оценку. Поэтому распределение длины серии контрольной карты Фазы II уже не следует геометрическому распределению, так как участвующие в этом случае испытания Бернулли (сигнал или отсутствие сигнала на каждой выборке или в каждый момент времени) уже не независимы, и потому для описания эффективности карты в этом случае обычно не используют FAR, поскольку FAR зависит от типа и значения оценок (полученных, например, в ходе анализа Фазы I). В случае IC \(\delta = 0\), и величина \(\beta(\alpha, 0, m, n, Y)\) обозначает условную вероятность сигнала в состоянии IC, то есть условную FAR в Фазе II.

Итак, подводя итог, в случае KU, то есть когда среднее известно или задано, а дисперсия неизвестна и потому оценена, для карты Шухарта \(\overline{X}\):

\begin{enumerate}
	\item Условно распределение длины серии в Фазе II является геометрическим с вероятностью успеха \(\beta(\alpha, \delta, m, n, y)\), где \(\alpha\) — номинальная FAR, \(\delta = (\mu_1 - \mu_0)/\sigma\), а \(y\) — реализация случайной величины \(Y\) с распределением хи-квадрат и \(\nu = mn\) степенями свободы. В этом контексте отметим, что соответствующая случайная величина \(Y\) — это масштабированная версия оценки \(S^2\). Можно было бы рассмотреть и другие оценки \(\sigma^2\) (например, основанную на среднем выборочных размахов), но теория для таких оценок, хотя в принципе и похожая, несколько сложнее и потому здесь опущена.
	\item Условно характеристики карты Шухарта \(\overline{X}\) Фазы II в случае KU можно получить из свойств геометрического распределения.
	\item Так, например, условное распределение длины серии задаётся своей функцией вероятности (pmf) \(P_C(N = j \mid Y = y) = \beta(y)[1 - \beta(y)]^{j-1}\), \(j = 1, 2, \ldots\), где для краткости мы пишем \(\beta(y) = \beta(\alpha, \delta, m, n, y)\).
\end{enumerate}

Следовательно, используя свойства геометрического распределения, мы можем полностью охарактеризовать условное распределение длины серии в случае KU и вычислить его различные характеристики, например условная средняя длина серии (\(ARL_C\)) карты Фазы II равна \([\beta(y)]^{-1}\). Отметим, что как функция \(Y\) величина \(ARL_C\) является случайной величиной \([\beta(Y)]^{-1}\), принимающей значение \([\beta(y)]^{-1}\) при значении \(y\) величины \(Y\). Аналогично, дисперсия и стандартное отклонение условного распределения длины серии — обе случайные величины, задаваемые как
\[
	\operatorname{var}_C(N \mid Y) = \frac{1 - \beta(Y)}{[\beta(Y)]^2}
\]
и
\[
	SDRL_C(N \mid Y) = \frac{1 - \beta(Y)}{\beta(Y)}
\]
соответственно.

Условное распределение длины серии и связанные с ним характеристики, такие как среднее, стандартное отклонение и процентили и т.д., помогают лучше понять эффекты оценивания параметров на эффективность контрольной карты Фазы II. Это важный момент, общий для всех контрольных карт Фазы II при оценённых параметрах, включая рассматриваемые здесь карты Шухарта. Он состоит в том, что оценка(и) параметров, являющаяся (являющиеся) случайной(ыми) величиной(ами), при простой подстановке в контрольные границы случая KK существенно меняет эффективность получаемой карты относительно случая KK. Более того, поскольку оценка является случайной величиной, две независимые эталонные выборки, обе из процесса в состоянии IC, скорее всего, дадут две разные оценки параметров, которые, в свою очередь, дадут две разные контрольные карты (границы) и, следовательно, разную эффективность в Фазе II. Таким образом, будет разброс в эффективности контрольных карт, даже если они используют оценки по независимым эталонным выборкам из одного и того же процесса в состоянии IC и даже если все карты настроены на одну и ту же номинальную \(ARL_{\text{IC}}\). Этот разброс эффективности называют вариацией «от пользователя к пользователю» (user-to-user variation). Признаем, что это может настораживать, и на поверхность выходят различные вопросы, особенно с точки зрения практика. В этой связи ряд недавних работ в литературе предлагал изучать условное распределение длины серии подробнее и тем самым лучше понимать влияние вариации «от пользователя к пользователю» на контрольные карты Фазы II. Некоторые из этих аспектов мы рассматриваем в последующих главах.

Хотя условное распределение длины серии помогает лучше понять эффекты оценивания параметров, можно дополнительно усреднить условное распределение по распределениям оценки параметра — процесс, известный как «деобусловливание», — и изучать «среднюю» эффективность карты. Например, безусловная вероятность сигнала в случае KU равна
\[
	P_U(\text{сигнал}) = \int_0^{\infty} \beta(y) f_{\chi^2_\nu}(y)\,dy,
\]
где \(f_{\chi^2_\nu}(y)\) обозначает плотность распределения хи-квадрат с \(\nu\) степенями свободы. Аналогично, безусловное, или маргинальное, распределение длины серии получается взятием условного распределения длины серии и усреднением по распределениям случайной величины \(Y\), которые для нормального распределения независимы. Таким образом, безусловная pmf длины серии задаётся как
\[
	P_U(N = j) = \int_0^{\infty} \beta(y)[1 - \beta(y)]^{j-1} f_{\chi^2_\nu}(y)\,dy
	\quad\text{для } j = 1, 2, 3, \ldots,
\]
что не является pmf геометрического распределения.

Такой способ вывода маргинального, или безусловного, распределения длины серии и связанных с ним характеристик называется методом обусловливания-деобусловливания (CUC). Вывод идёт сначала через обусловливание на оценке параметра, а затем через взятие математического ожидания результата по распределению вероятностей этой оценки.

Таким образом, в случае KU, используя тот факт, что условное распределение длины серии геометрическое с вероятностью успеха \(\beta(\alpha, \delta, m, n, y)\), и применяя метод CUC, среднее безусловного распределения длины серии равно
\[
	E_U(N) = \int_0^{\infty} \frac{1}{1 - \beta(y)}\, f_{\chi^2_\nu}(y)\,dy.
\]
Это безусловная ARL карты Шухарта \(\overline{X}\), обозначаемая \(ARL_U\). Отметим, что тот же результат можно получить из приведённой выше безусловной pmf \(P_U(N = j)\).

Помимо ARL, важной характеристикой эффективности является дисперсия (или стандартное отклонение) безусловного распределения длины серии. Вывод безусловной дисперсии интересен и потому объясняется подробнее. Выводу помогает результат, упомянутый в Главе 1: для любой случайной величины \(N\) безусловную дисперсию можно получить с помощью условного среднего, условной дисперсии и формулы
\[
	\operatorname{var}_U(N) = E(\operatorname{var}(N \mid Y)) + \operatorname{var}(E(N \mid Y))
\]
(см. также, например, Hogg, McKean and Craig, 2005), где математическое ожидание берётся по распределению случайной величины \(Y\). Для членов \(\operatorname{var}(N \mid Y)\) и \(E(N \mid Y)\) вспомним, что при заданном \(Y\) распределение длины серии геометрическое с вероятностью \(\beta(Y)\), и потому из свойств геометрического распределения
\[
	\operatorname{var}(N \mid Y) = \frac{\beta(Y)}{(1 - \beta(Y))^2}
\]
и
\[
	E(N \mid Y) = \frac{1}{1 - \beta(Y)}
\]
соответственно. Таким образом,
\[
\begin{aligned}
	\operatorname{var}(E(N \mid Y))
	&= E\big((E(N \mid Y))^2\big) - (E(E(N \mid Y)))^2 \\
	&= E\left(\left[\frac{1}{1 - \beta(Y)}\right]^2\right)
	   - \left[E\left(\frac{1}{1 - \beta(Y)}\right)\right]^2,
\end{aligned}
\]
и подстановкой получаем
\[
\operatorname{var}_U(N) = E\left(\frac{\beta(Y)}{(1 - \beta(Y))^2}\right)
	+ E\left(\left[\frac{1}{1 - \beta(Y)}\right]^2\right)
	- \left[E\left(\frac{1}{1 - \beta(Y)}\right)\right]^2.
\]
Переписав это выражение через интегралы, получаем безусловную дисперсию распределения длины серии в случае KU:
\[
	\operatorname{var}_U(N) = \int_0^{\infty} \frac{1 + \beta(y)}{(1 - \beta(y))^2}\, f_{\chi^2_\nu}(y)\,dy
	+ \int_0^{\infty} \left[\frac{1}{1 - \beta(y)}\right]^2 f_{\chi^2_\nu}(y)\,dy
	- \left(\int_0^{\infty} \frac{1}{1 - \beta(y)}\, f_{\chi^2_\nu}(y)\,dy\right)^2,
\]
где читателю следует помнить, что \(\beta(y) = \beta(\alpha, \delta, m, n, y)\), как определено ранее. Отметим, что хотя распределение длины серии всегда дискретно, контролируемая переменная, а значит, и соответствующие оценки параметров, могут быть как непрерывными (как в настоящем случае), так и дискретными, как в случае карт по альтернативным признакам (см. Montgomery, 2009, глава 7). В последнем случае оценка \(Y\) будет дискретной, и потому безусловное распределение длины серии, его среднее и дисперсию можно получить аналогичным образом с помощью метода CUC, за исключением того, что интеграл теперь заменяется суммой по точкам вероятностной массы оценки \(Y\), а плотность хи-квадрат в интеграле — соответствующей pmf величины \(Y\). Примеры читатель найдёт в Human and Chakraborti (2006).

Другие моменты и характеристики распределения длины серии также можно получить аналогично с помощью метода CUC, который использует тот факт, что (i) условная длина серии при заданной оценке параметра является геометрическим распределением, а (ii) деобусловливание включает усреднение (взятие математического ожидания) по распределению оценки параметра, которое следует некоторому известному распределению вероятностей.

Обычно безусловное распределение длины серии называют распределением длины серии контрольной карты, но условное распределение длины серии приобретает всё большую значимость, поскольку оно отражает и объясняет влияние присущего оценкам параметров разброса на эффективность контрольной карты. Подробности читатель найдёт в недавней литературе.

Подводя итог: безусловные свойства длины серии, такие как среднее и стандартное отклонение, фиксированы и описывают среднюю эффективность контрольной карты по многим эталонным выборкам, тогда как условные свойства длины серии меняются в зависимости от значений оценок параметров для данной эталонной выборки и показывают влияние оценивания параметров исходя из значения оценки, имеющейся у пользователя в его конкретной ситуации. И условная, и безусловная эффективность могут играть роль в практическом применении контрольных карт с оценёнными параметрами. Распределение длины серии в случае KK и связанные с ним характеристики более 50 лет служили эталонной нормой при проектировании контрольной карты — вплоть до недавнего времени, когда в свете эффектов оценивания параметров на первый план вышли безусловное распределение длины серии и его характеристики. Но, как отмечалось ранее, признаётся, что хотя процесс усреднения и может дать полезные общие ориентиры, условная длина серии и связанные с ней характеристики могут быть очень полезны при практическом внедрении и понимании эффективности карты, когда в контрольных картах используются оценённые параметры. То же резюме применимо ко всем типам контрольных карт — одномерным или многомерным, по количественным или альтернативным признакам, параметрическим или непараметрическим, — где неизвестные параметры оцениваются по данным.

Далее рассмотрим случай, когда оба параметра неизвестны, то есть случай UU. Анализ аналогичен случаю KU, за исключением того, что теперь участвуют два параметра и их оценки, что повышает сложность. Мы постараемся указать на сходство со случаем KU и выделить главные различия. Как отмечалось ранее, в Фазе II оценки параметров, а значит, и контрольные границы, получаются в Фазе I и применяются для постоянного мониторинга состояния процесса в Фазе II. Снова, как и в случае KU, предположим, что в Фазе I процесс следует нормальному распределению, так что \(X_{ij} \sim N(\mu_0, \sigma_0^2)\) для \(i = 1, 2, 3, \ldots, m\) и \(j = 1, 2, 3, \ldots, n\). Кроме того, предположим, что неизвестное среднее процесса в состоянии IC \(\mu_0\) и неизвестное стандартное отклонение процесса в состоянии IC \(\sigma_0\) мы оцениваем общим средним, которое является общим средним средних эталонных выборок,
\[
	\overline{\overline{X}} = \frac{1}{m}\sum_{i=1}^{m}\overline{X}_i
	= \frac{1}{mn}\sum_{i=1}^{m}\sum_{j=1}^{n} X_{ij},
\]
и объединённой оценкой
\[
	S_p = \sqrt{\frac{\sum_{i=1}^{m} S_i^2}{m}}.
\]
Таким образом, в отличие от случая KU, в анализе для случая UU участвуют две оценки, а именно \(\overline{\overline{X}}\) и \(S_p\). Отметим также, что, в отличие от случая KU, оценка \(\sigma^2\), а именно \(S_p^2\), имеет масштабированное распределение хи-квадрат, но с \(m(n - 1)\) степенями свободы. Таким образом, при оценивании среднего \(\mu\) по \(m\) выборкам Фазы I теряется \(m\) степеней свободы (напомним, что в случае KU было \(mn\) степеней свободы).

Снова двусторонние контрольные границы получаются подстановкой точечных оценок в контрольные границы для случая KK:
\[
	\widehat{LCL} = \overline{\overline{X}} - z_{\alpha/2}\frac{S_p}{\sqrt{n}}
	\quad\text{и}\quad
	\widehat{UCL} = \overline{\overline{X}} + z_{\alpha/2}\frac{S_p}{\sqrt{n}}.
\]
Как и раньше, предположим, что выборки Фазы II следуют распределению \(N(\mu_1, \sigma_1^2)\). Тогда сигнальные события в Фазе II можно записать как
\[
	\overline{X}_i < \widehat{LCL} \quad\text{или}\quad \overline{X}_i > \widehat{UCL}
	\quad\text{для } i = m + 1, m + 2, m + 3, \ldots,
\]
что можно переписать как
\[
	\frac{\sqrt{n}|\overline{X}_i - \overline{\overline{X}}|}{S_p} \ge z_{\alpha/2}
	\quad\text{для } i = m + 1, m + 2, m + 3, \ldots
\]
Сравним это сигнальное событие с событием в случае KU. Как и в случае KU, при использовании метода CUC вероятность сигнального события в случае UU вычисляется в два шага. Сначала вероятность сигнала вычисляется при заданных (обусловленно на) двух оценках параметров \(\overline{\overline{X}}\) и \(S_p\), которые обе являются случайными величинами. Эта вероятность называется условной вероятностью сигнала и играет ключевую роль в последующих построениях.

Чтобы показать это, отметим, что вероятность сигнала в случае UU задаётся как \(P\left(\dfrac{\sqrt{n}|\overline{X}_i - \overline{\overline{X}}|}{S_p} \ge z_{\alpha/2}\right)\) для \(i = m + 1, m + 2, m + 3, \ldots\). Её можно записать как
\[
\begin{aligned}
	&= 1 - P\left(\frac{-z_{\alpha/2}S_p}{\sqrt{n}} + \overline{\overline{X}} < X_i
	< \frac{z_{\alpha/2}S_p}{\sqrt{n}} + \overline{\overline{X}}\right) \\
	&= 1 - P\left(-\frac{z_{\alpha/2}S_p}{\sigma_1} + \frac{\sqrt{n}(\overline{\overline{X}} - \mu_0)}{\sigma_1}
	- \frac{\sqrt{n}(\mu_1 - \mu_0)}{\sigma_1} < \frac{\sqrt{n}(X_i - \mu_1)}{\sigma_1} \right.\\
	&\qquad\qquad \left.< \frac{z_{\alpha/2}S_p}{\sigma_1} + \frac{\sqrt{n}(\overline{\overline{X}} - \mu_0)}{\sigma_1}
	- \frac{\sqrt{n}(\mu_1 - \mu_0)}{\sigma_1}\right).
\end{aligned}
\]
Теперь, полагая \(Z = \dfrac{\sqrt{mn}(\overline{\overline{X}} - \mu_0)}{\sigma_0}\), \(Z_1 = \dfrac{\sqrt{n}(X_i - \mu_1)}{\sigma_1}\) и \(\delta = \dfrac{(\mu_1 - \mu_0)}{\sigma_0}\), можно переписать приведённую выше вероятность как
\[
\begin{aligned}
	&= 1 - P\left(\left(\frac{\sigma_0}{\sigma_1}\right)\frac{-z_{\alpha/2}S_p}{\sigma_0} + \frac{Z}{\sqrt{m}} - \delta\sqrt{n} < Z_1
	< \left(\frac{\sigma_0}{\sigma_1}\right)\frac{z_{\alpha/2}S_p}{\sigma_1} + \frac{Z}{\sqrt{m}} - \delta\sqrt{n}\right).
\end{aligned}
\]
Наконец, полагая \(Y = \dfrac{vS_p^2}{\sigma_0^2}\), что следует распределению хи-квадрат с \(\nu = m(n - 1)\) степенями свободы, и обозначая \(\tau = \dfrac{\sigma_0}{\sigma_1}\), можно переписать последнее выражение как
\[
\begin{aligned}
	&= 1 - P\left(\left(-z_{\alpha/2}\sqrt{\frac{Y}{v}} + \frac{Z}{\sqrt{m}} - \delta\sqrt{n}\right)\tau < Z_1
	< \left(z_{\alpha/2}\sqrt{\frac{Y}{v}} + \frac{Z}{\sqrt{m}} - \delta\sqrt{n}\right)\tau\right).
\end{aligned}
\]
При \(\sigma_1 = \sigma_0\) величина \(\tau = 1\). Наконец, поскольку \(Z_1\) имеет распределение \(N(0,1)\), условная вероятность сигнала в Фазе II для карты среднего в случае UU равна
\[
\begin{aligned}
	&1 - \Phi\left(\left(z_{\alpha/2}\sqrt{\frac{Y}{v}} + \frac{Z}{\sqrt{m}} - \delta\sqrt{n}\right)\tau\right) \\
	&\quad + \Phi\left(\left(-z_{\alpha/2}\sqrt{\frac{Y}{v}} + \frac{Z}{\sqrt{m}} - \delta\sqrt{n}\right)\tau\right)
	= \beta(\alpha, \delta, \tau, m, n, z, y),
\end{aligned}
\]
где \(Z\) и \(Y\) — независимые случайные величины. Величина \(Z\) имеет стандартное нормальное распределение, а \(Y\) — распределение хи-квадрат с \(v = m(n - 1)\) степенями свободы. Опять же можно было бы рассмотреть другие оценки \(\sigma^2\), и в дальнейшем пришлось бы работать с их распределениями. Мы оставляем это читателю.
% CHECK: формула для β(α, δ, τ, m, n, z, y) в книге набрана с разнобоем в знаках слагаемых Z/√m внутри двух Φ; воспроизведено дословно, форма в целом совпадает с условной вероятностью сигнала в случае UU.

В любом случае и в случае UU условная вероятность сигнала играет важную роль в анализе эффективности контрольной карты при оценённых параметрах. Подробнее об этом мы скажем позже. В случае IC имеем \(\delta = 0\), и при типичном предположении \(\sigma_1 = \sigma_0\), то есть \(\tau = 1\), величина \(\beta(\alpha, 0, 1, m, n, z, y)\) обозначает условную вероятность сигнала (FAR) в состоянии IC в Фазе II. Однако, поскольку эта вероятность является функцией случайных величин \(Z\) и \(Y\), она является случайной величиной со своим собственным распределением вероятностей. Именно поэтому, как отмечалось ранее, когда параметры оцениваются, в Фазе II обычно не используют FAR для описания эффективности карты: FAR зависит от типа оценок, полученных при анализе Фазы I, и сама является случайной величиной с распределением вероятностей.

Более того, как и в случае KU, в случае UU сигнальные события Фазы II зависимы, поскольку в контрольных границах используются одни и те же оценки параметров Фазы I и каждая выборка Фазы II сравнивается с одними и теми же контрольными границами, содержащими эти оценки. Таким образом, в случае UU, как и в случае KU, распределение длины серии контрольной карты Фазы II уже не следует геометрическому распределению, поскольку участвующие в этом случае испытания Бернулли уже не независимы. Однако, опять же как и в случае KU, при заданных оценках параметров Фазы I (или при обусловливании на соответствующие случайные величины) сигнальные события независимы. Это наблюдение существенно облегчает многие статистические расчёты характеристик карты Фазы II в случае UU, как и в случае KU. Итак, подводя итог, в случае UU для карты Шухарта \(\overline{X}\):

\begin{enumerate}
	\item Условное распределение длины серии в Фазе II геометрическое с вероятностью успеха \(\beta(\alpha, \delta, \tau, m, n, z, y)\).
	\item Следовательно, условно свойства карты Шухарта \(\overline{X}\) Фазы II следуют из свойств геометрического распределения; например, условное распределение длины серии задаётся своей функцией вероятности (pmf)
	\[
		P_C(N = j \mid z, y) = \beta(z, y)[1 - \beta(z, y)]^{j-1}, \quad j = 1, 2, \ldots,
	\]
	где для краткости мы пишем \(\beta(z, y) = \beta(\alpha, \delta, \tau, m, n, z, y)\). Теперь, используя свойства геометрического распределения, можно легко получить различные характеристики условного распределения длины серии. Условная средняя длина серии (\(ARL_C\)) карты Фазы II, например, равна \([\beta(z, y)]^{-1}\). Аналогично, дисперсия и стандартное отклонение задаются как \(\operatorname{var}_C(N \mid Z, Y) = \dfrac{1 - \beta(Z, Y)}{[\beta(Z, Y)]^2}\) и \(SDRL_C(N \mid Z, Y) = \dfrac{1 - \beta(Z, Y)}{\beta(Z, Y)}\) соответственно.
\end{enumerate}

Как отмечалось ранее, оценки параметров, подставленные в контрольные границы, могут существенно изменить/повлиять на эффективность контрольной карты и условное распределение длины серии. Связанные с ним характеристики, такие как среднее, стандартное отклонение и т.д., помогают описать эффекты оценивания параметров на эффективность контрольной карты Фазы II. Оценки — это случайные величины, и две независимые эталонные выборки из одного и того же процесса в состоянии IC скорее всего дадут два разных набора оценок параметров, которые, в свою очередь, дадут разные контрольные карты (границы) и, следовательно, разную эффективность карты в Фазе II. Суть в том, что будет разброс эффективности контрольных карт между пользователями, когда они используют оценки по независимым эталонным выборкам из одного и того же процесса в состоянии IC, и даже если контрольные карты были настроены на одну и ту же номинальную \(ARL_{\text{IC}}\). Этот разброс можно назвать вариацией «от пользователя к пользователю», и её непросто контролировать, если нет большого объёма данных, что не всегда возможно на практике. В этой связи, как отмечалось ранее, ряд недавних работ предлагал изучать условное распределение длины серии и влияние вариации «от пользователя к пользователю» на контрольные карты Фазы II. Некоторые из этих аспектов мы рассматриваем в последующих главах.

Усредняя по распределениям оценок параметров — процесс, известный как «деобусловливание», — можно изучать «среднюю» эффективность карты; например, безусловная вероятность сигнала в случае UU равна
\[
	P_U(\text{сигнал}) = \int_{-\infty}^{\infty}\int_0^{\infty} \beta(z, y) g(z, y)\,dy\,dz
	= \int_{-\infty}^{\infty}\int_0^{\infty} \beta(z, y)\varphi(z) f_{\chi^2_\nu}(y)\,dy\,dz,
\]
где \(f_{\chi^2_\nu}(y)\) обозначает плотность распределения хи-квадрат с \(\nu\) степенями свободы. Отметим, что для нормального распределения оценки среднего и дисперсии статистически независимы. Поэтому совместное распределение распадается в произведение отдельных распределений — одного стандартного нормального и одного хи-квадрат. Однако это может быть не так во всех случаях. Аналогично, используя метод CUC, безусловную pmf длины серии можно получить интегрированием условной pmf длины серии по совместной плотности \(Z\) и \(Y\), которая является произведением маргинальных плотностей \(Z\) и \(Y\) соответственно, поскольку \(Z\) и \(Y\) независимы. Следовательно,
\[
	P_U(N = j) = \int_{-\infty}^{\infty}\int_0^{\infty} \beta(z, y)[1 - \beta(z, y)]^{j-1}
	\varphi(z) f_{\chi^2_\nu}(y)\,dy\,dz \quad\text{для } j = 1, 2, 3, \ldots
\]
Аналогично, используя тот факт, что условное распределение длины серии геометрическое с вероятностью успеха \(\beta(z, y)\), и применяя метод CUC, математическое ожидание безусловного распределения длины серии равно
\[
	E_U(N) = \int_{-\infty}^{\infty}\int_0^{\infty} \frac{1}{1 - \beta(z, y)}\, \varphi(z) f_{\chi^2_\nu}(y)\,dy\,dz.
\]
Это безусловная ARL карты Шухарта \(\overline{X}\) в случае UU. В общем случае для моментов более высокого порядка, скажем \(k\)-й нецентральный момент распределения длины серии, его можно получить из \(k\)-го условного нецентрального момента, то есть \(E(N^k) = E(E(N^k \mid Z, Y))\), используя тот факт, что \(N \mid Z, Y\) следует геометрическому распределению. В частности, дисперсия (или стандартное отклонение) безусловного распределения длины серии — важная характеристика эффективности. Можно показать, что она равна
\[
\begin{aligned}
	\operatorname{var}_U(N)
	&= \int_{-\infty}^{\infty}\int_0^{\infty} \frac{1 + \beta(z, y)}{(1 - \beta(z, y))^2}\, \varphi(z) f_{\chi^2_\nu}(y)\,dy\,dz \\
	&\quad - \left(\int_{-\infty}^{\infty}\int_0^{\infty} \frac{1}{1 - \beta(z, y)}\, \varphi(z) f_{\chi^2_\nu}(y)\,dy\,dz\right)^2.
\end{aligned}
\]
Вывод безусловной дисперсии (и стандартного отклонения) интересен и объясняется подробнее. Он следует из результата, что для любой случайной величины \(N\)
\[
	\operatorname{var}_U(N) = E(\operatorname{var}(N \mid Z, Y)) + \operatorname{var}(E(N \mid Z, Y))
\]
(см., например, Hogg, McKean and Craig, 2005), где математическое ожидание берётся по распределениям двух независимых случайных величин \(Z\) и \(Y\). Более того, при предположении о нормальном распределении \(Z\) и \(Y\) независимы. Для членов \(\operatorname{var}(N \mid Z, Y)\) и \(E(N \mid Z, Y)\) отметим, что при заданных \(Z\) и \(Y\) распределение длины серии геометрическое с вероятностью \(\beta(Z, Y)\), и потому из свойств геометрического распределения
\[
	\operatorname{var}(N \mid Z, Y) = \frac{\beta(Z, Y)}{(1 - \beta(Z, Y))^2}
	\quad\text{и}\quad
	E(N \mid Z, Y) = \frac{1}{1 - \beta(Z, Y)}.
\]
Таким образом,
\[
\begin{aligned}
	\operatorname{var}(E(N \mid Z, Y))
	&= E\big((E(N \mid Z, Y))^2\big) - (E(E(N \mid Z, Y)))^2 \\
	&= E\left(\left[\frac{1}{1 - \beta(Z, Y)}\right]^2\right)
	   - \left[E\left(\frac{1}{1 - \beta(Z, Y)}\right)\right]^2,
\end{aligned}
\]
и подстановкой получаем
\[
	\operatorname{var}_U(N) = E\left(\frac{\beta(Z, Y)}{(1 - \beta(Z, Y))^2}\right)
	+ E\left(\left[\frac{1}{1 - \beta(Z, Y)}\right]^2\right)
	- \left[E\left(\frac{1}{1 - \beta(Z, Y)}\right)\right]^2.
\]
Переписав это выражение через интегралы, получаем безусловную дисперсию
\[
\begin{aligned}
	\operatorname{var}_U(N)
	&= \int_{-\infty}^{\infty}\int_0^{\infty} \frac{\beta(Z, Y) + 1}{(1 - \beta(Z, Y))^2}\,
	\varphi(z) f_{\chi^2_\nu}(y)\,dy\,dz \\
	&\quad + \int_{-\infty}^{\infty}\int_0^{\infty}\left[\frac{1}{1 - \beta(Z, Y)}\right]^2
	\varphi(z) f_{\chi^2_\nu}(y)\,dy\,dz,
\end{aligned}
\]
показанную выше. Следовательно, безусловное стандартное отклонение длины серии равно \(SDRL_U(N) = \sqrt{\operatorname{var}_U(N)}\).
% CHECK: в книге опечатка «SDRLU (N ) = varU (N )» — пропущен корень; в переводе восстановлен корень.

Обычно безусловное распределение длины серии называют распределением длины серии контрольной карты, и именно его использовали при традиционной оценке эффективности контрольной карты. Однако в последнее время условное распределение длины серии приобрело большую значимость, поскольку оно отражает и объясняет влияние присущего оценкам параметров разброса на эффективность контрольной карты. Другие моменты и характеристики распределения длины серии можно получить аналогично с помощью метода CUC, который использует тот факт, что (i) условная длина серии при заданных оценках параметров является геометрическим распределением, а (ii) деобусловливание включает усреднение (взятие математического ожидания) по распределениям двух оценок параметров, которые независимы и следуют некоторым известным распределениям вероятностей. Подводя итог: безусловные свойства длины серии, такие как среднее и стандартное отклонение, фиксированы и описывают среднюю эффективность контрольной карты по многим выборкам, тогда как условные свойства длины серии таковыми не являются и показывают влияние оценивания параметров исходя из оценок, имеющихся у пользователя в его ситуации. И условная, и безусловная эффективность могут играть роль в практическом применении контрольных карт с оценёнными параметрами.

Для случая UK вероятность сигнала, условное распределение длины серии, его различные моменты, безусловное распределение длины серии и его различные моменты можно получить аналогично; это оставлено читателю в качестве упражнений. Более подробно можно узнать в Chakraborti (2000).

Рассмотрев подробно карту Шухарта для среднего, мы теперь лучше подготовлены к изучению некоторых более продвинутых или сложных карт. К ним относятся карты CUSUM и EWMA, а также карты Шухарта с дополнительными правилами серий, которые, как ожидается, более эффективны при обнаружении малых, устойчивых изменений лежащего в основе параметра. Расчёт распределения длины серии и связанных с ним свойств для этих карт в случае U сложнее, но и здесь можно применить метод CUC. Ниже он кратко изложен, сначала в общих чертах. Напомним, что для этих карт выражения для распределения длины серии и некоторых связанных с ним характеристик в случае K были даны в уравнениях 3.21–3.24 с использованием подхода на основе цепей Маркова (MC). Это отправная точка в случае U. Как отмечалось ранее, ключ к использованию этих выражений — построение существенной матрицы переходных вероятностей для каждой карты. В случае U мы сначала используем эти формулы условно при заданных оценках параметров, а затем применяем метод CUC, чтобы найти безусловное распределение длины серии и связанные с ним свойства. Предположим, например, что имеется единственный неизвестный параметр, который нужно оценить, и соответствующая оценка по Фазе I равна \(T\). Тогда элементы существенной матрицы вероятностей \(Q\), так называемые переходные вероятности, мы вычисляем, используя условное распределение статистики для нанесения на карту в Фазе II при заданной оценке \(T\). Таким образом, каждый элемент матрицы переходных вероятностей вычисляется условно на \(T\); пусть \(Q(T)\) обозначает получающуюся условную существенную матрицу вероятностей. Далее мы можем применить выражения из уравнений 3.22–3.24 условно на \(T\) и получить
\begin{equation}
	P(N = t \mid T) = \xi Q(T)^{t-1}(I - Q(T))\mathbf{1}
	\quad\text{для } t = 1, 2, 3, \ldots
	\label{eq:3.36}
\end{equation}
\begin{equation}
	E(N \mid T) = \xi (I - Q(T))^{-1}\mathbf{1}
	\label{eq:3.37}
\end{equation}
\begin{equation}
	SD(N \mid T) = \sqrt{\xi (I + Q(T))(I - Q(T))^{-2}\mathbf{1} - (E(N \mid T))^2}
	\label{eq:3.38}
\end{equation}
\begin{equation}
	P(N \le j \mid T) = 1 - \xi Q(T)^j \mathbf{1}
	\quad\text{для } j = 1, 2, 3, \ldots
	\label{eq:3.39}
\end{equation}
Наконец, используя метод CUC, безусловное распределение длины серии получается взятием математического ожидания (усреднением или интегрированием) по распределению \(T\). Так, например, безусловная ARL карты равна
\begin{equation}
	\int_{-\infty}^{\infty} \{\xi [I - Q(t)]^{-1}\mathbf{1}\}\, f(t)\,dt,
	\label{eq:3.40}
\end{equation}
где предполагается, что \(T\) непрерывна с плотностью \(f(t)\). Другие моменты и связанные характеристики распределения длины серии можно найти аналогичным образом. Отметим, что аналогичный анализ возможен, если \(T\) дискретна; в этом случае указанные выражения для безусловного распределения длины серии получаются суммированием по pmf \(f(t)\) величины \(T\). Заинтересованным читателям предлагается попробовать некоторые примеры. Некоторые примеры можно найти в Chakraborti, Eryilmaz and Human (2009).
\subsubsection{Карта Шухарта для дисперсии в случае U}

Распределение длины серии карт Шухарта для дисперсии можно получить аналогично случаю среднего с помощью метода CUC. Здесь мы кратко набросаем пример для карты \(S\) с \(S_p\) в качестве оценки \(\sigma\) и выведем распределение длины серии в состоянии IC.

Отметим, что в этом случае вероятностные границы карты Фазы II задаются как \(LCL = H_3^*(m,n)S_p\), \(CL = S_p\) и \(UCL = H_4^*(m,n)S_p\), где \(H_3^*(m,n) = \sqrt{\chi^2_{\alpha(m,n)/2,\,n-1}/(n - 1)}\), \(H_4^* = \sqrt{\chi^2_{1-\alpha(m,n)/2,\,n-1}/(n - 1)}\), а \(\alpha(m,n)\) — безусловная вероятность того, что статистика для нанесения на карту Фазы II, стандартное отклонение \(S_i\), окажется вне нижней или верхней контрольной границы, когда процесс находится в состоянии IC. Отметим, что в случае известных параметров это вероятность ложной тревоги, то есть FAR, и она зависит только от объёма выборки \(n\), а в случае неизвестных параметров вероятность \(\alpha(m,n)\) зависит и от \(m\), и от \(n\). Как только величина \(\alpha(m,n)\) определена, вероятностные границы легко получить по соответствующим процентилям распределения хи-квадрат с \(n - 1\) степенями свободы. Хотя для удобства мы можем называть \(\alpha(m,n)\) в случае U FAR (вероятностью сигнала, когда процесс находится в состоянии IC), отметим, что она не имеет того же ясного смысла, что в случае K, поскольку в случае U сигнальные события зависимы. Некоторые подробности вывода \(\alpha(m,n)\) даны в приложении 3.3. Значения \(\alpha(m,n)\) рассчитаны и приведены в таблице H приложения A вместе с соответствующими контрольными границами \(H_3^*(m,n)\) и \(H_4^*(m,n)\) для \(n = 5\) и 10 и некоторых значений \(m\) от 5 до 100 при номинальных значениях \(ARL_{\text{IC}}\), равных 370 и 500. Отметим, что значения \(\alpha(m,n)\) зависят и от \(m\), и от \(n\), и при заданном \(n\) они возрастают с ростом \(m\) (с увеличением числа выборок Фазы I).

Другие моменты условного и безусловного распределения длины серии в состоянии IC можно найти аналогичным способом. Отметим, что когда в контрольных границах Фазы I используются другие оценки \(\sigma\), например основанные на \(R\) или \(\overline{S}\), или в Фазе II используется другая карта, например карта \(R\), аналогичный подход можно применить для вывода распределения длины серии и его различных моментов с помощью метода CUC. По существу, нам нужны распределения вероятностей статистики для нанесения на карту и оценки \(\sigma\) по Фазе I. Отметим, что тот же подход можно использовать и для вывода распределения длины серии в состоянии OOC, что полезно при изучении эффективности этих карт в обнаружении сдвигов и при сравнениях. Детали мы оставляем читателю.

\subsubsection{Карта CUSUM для среднего в случае U}

Читателю следует вернуться к примеру 3.6 в Разделе 3.6.1.2. В этом примере предполагается, что лежащее в основе распределение процесса — \(N(0,1)\). Рассмотрим, например, одношаговую переходную вероятность \(p_{00}\). Было показано, что \(p_{00} = P(X_t \le 1) = 0.841\). Однако напомним, что \(P(X_t \le 1) = 0.841\) получено с использованием функции распределения \(N(0,1)\), что требует знания не только нормальности распределения процесса, но и того, что среднее и стандартное отклонение равны соответственно 0 и 1. Предположим теперь, что мы находимся в случае KU, то есть среднее процесса известно (скажем, равно нулю), а дисперсия процесса неизвестна. В этом случае у нас есть единственный неизвестный параметр \(\sigma\), который нужно оценить некоторой оценкой \(T\) (обычно получаемой по анализу Фазы I), и мы вычисляем элементы существенной матрицы вероятностей \(Q\), то есть переходные вероятности, условно при заданной оценке \(T\). Теперь \(p_{00} = P(X_t \le 1)\) вычисляется как \(\hat{p}_{00} = P(X_t \le 1 \mid T) = \Phi(1/T)\), где \(T\) — некоторая оценка \(\sigma\). То же относится и к другим одношаговым переходным вероятностям. Таким образом получают \(Q(T)\) — оценённую условную существенную матрицу вероятностей — и подставляют её в уравнения 3.36–3.39. Наконец, с помощью метода CUC безусловное распределение длины серии получают взятием математического ожидания по распределению \(T\), как объяснялось выше и проиллюстрировано в уравнении 3.40 для безусловной ARL. Случай UK рассматривается аналогично. Однако на этот раз нужно оценить и учесть среднее процесса, а для случая UU, где нужно оценить и среднее процесса, и дисперсию процесса, мы обусловливаем на, скажем, \(T_1\) и \(T_2\), где \(T_1\) обозначает оценку среднего процесса, а \(T_2\) — оценку дисперсии процесса. Читателю следует также обратиться к Jones, Champ and Rigdon (2004) за альтернативным выводом распределения длины серии карты CUSUM с оценёнными параметрами.

\subsubsection{Карта EWMA для среднего в случае U}

Читателю следует вернуться к примеру 3.7 в Разделе 3.6.1.2. В этом примере предполагается, что лежащее в основе распределение процесса — \(N(0,1)\). Рассмотрим одношаговую переходную вероятность \(p_{01}\) из этого примера. Было показано, что \(p_{01}\) равна 0.226, но \(0.226 = \Phi(2.067) - \Phi(0.689)\) получено с использованием функции распределения \(N(0,1)\). Однако даже если мы предполагаем нормальность, один или несколько параметров могут быть неизвестны и их нужно оценить. В таком случае мы действуем так же, как в случае карты CUSUM, рассмотренном выше, и вычисляем элементы матрицы переходных вероятностей при условии заданных оценок параметров. Затем получают получающуюся условную существенную матрицу вероятностей и подставляют её в уравнения 3.36–3.39. Наконец, с помощью метода CUC безусловное распределение длины серии получают взятием математического ожидания по распределению оценки(ок), как объяснялось выше и проиллюстрировано в уравнении 3.40 для безусловной ARL. Читателю следует обратиться к Jones, Champ and Rigdon (2001) за подробным обсуждением распределения длины серии карты EWMA с оценёнными параметрами.

В случае неизвестных параметров метод CUC успешно применялся не только к различным параметрическим картам, но и к непараметрическим картам, как будет видно в следующей главе.

\section{Усовершенствования контрольной карты}

\subsection{Расчёт длины серии для правил сигнализации на основе серий в случае U}

Zhang and Castagliola (2010) рассмотрели карту Шухарта \(\overline{X}\) с правилами серий, когда параметры процесса неизвестны. Напомним, что если среднее и стандартное отклонение процесса в состоянии IC, \(\mu_0\) и \(\sigma_0\), оба неизвестны, в случае UU их оценивают по эталонному набору данных Фазы I объёмом \(m \ge 1\) подгрупп, каждая объёма \(n \ge 1\). Можно использовать широко распространённые оценки \(\hat{\mu}_0 = \overline{\overline{X}}\) и \(\hat{\sigma}_0 = S_p\), где напомним, что \(\overline{\overline{X}}\) — среднее средних выборок Фазы I, а \(S_p\) — так называемая объединённая оценка стандартного отклонения. Пусть \(\{Y_{i1}, \ldots, Y_{in}\}\) обозначает выборку Фазы II, а \(\overline{Y}_i\) — среднее этой выборки. Теперь мы подставляем оценки в известные \(k\)-сигма контрольные границы Шухарта и получаем оценённые границы \(\widehat{LCL}/\widehat{UCL} = \hat{\mu}_0 \pm k\hat{\sigma}_0/\sqrt{n}\), как показано ранее. Поскольку \(\hat{\mu}_0\) и \(\hat{\sigma}_0\) — случайные величины, \(\widehat{LCL}\) и \(\widehat{UCL}\) также случайные величины. Условные вероятности сигнала можно найти как \(p_L = P(\overline{Y}_i < \widehat{LCL} \mid \hat{\mu}_0, \hat{\sigma}_0)\) и \(p_U = P(\overline{Y}_i > \widehat{UCL} \mid \hat{\mu}_0, \hat{\sigma}_0)\), так что \(p = P(\widehat{LCL} \le \overline{Y}_i \le \widehat{UCL} \mid \hat{\mu}_0, \hat{\sigma}_0) = 1 - p_L - p_U\). Подставляя \(\widehat{LCL}/\widehat{UCL} = \hat{\mu}_0 \pm k\hat{\sigma}_0\) в выражения для \(p_L\) и \(p_U\), находим \(p_L = P(\overline{Y}_i < \hat{\mu}_0 - k\hat{\sigma}_0 \mid \hat{\mu}_0, \hat{\sigma}_0)\) и \(p_U = P(\overline{Y}_i > \hat{\mu}_0 + k\hat{\sigma}_0 \mid \hat{\mu}_0, \hat{\sigma}_0)\) соответственно. Далее мы не обсуждаем \(p_{IC}\) подробно, поскольку её можно найти через \(p = 1 - p_L - p_U\). Вычитая среднее Фазы II \(\mu_1 = \mu_0 + \delta\sigma_0\) и деля обе части неравенств на \(\sigma_0/\sqrt{n}\), находим
\[
	p_L = P\left((\overline{Y}_i - \mu_0 - \delta\sigma_0)\frac{\sqrt{n}}{\sigma_0}
	< (\hat{\mu}_0 - k\hat{\sigma}_0 - \mu_0 - \delta\sigma_0)\frac{\sqrt{n}}{\sigma_0}
	\;\middle|\; \hat{\mu}_0, \hat{\sigma}_0\right)
\]
и
\[
	p_U = P\left((\overline{Y}_i - \mu_0 - \delta\sigma_0)\frac{\sqrt{n}}{\sigma_0}
	> (\hat{\mu}_0 + k\hat{\sigma}_0 - \mu_0 - \delta\sigma_0)\frac{\sqrt{n}}{\sigma_0}
	\;\middle|\; \hat{\mu}_0, \hat{\sigma}_0\right)
\]
соответственно.

По определению, при предположении, что распределение Фазы II имеет сдвиг среднего, \(\overline{Y}_i \sim N(\mu_1 = \mu_0 + \delta\sigma_0, \sigma_0/\sqrt{n})\), так что приведённые выше выражения можно записать через функцию распределения стандартного нормального распределения:
\[
	p_L = \Phi\left((\hat{\mu}_0 - \mu_0)\frac{\sqrt{n}}{\sigma_0} - k\frac{\hat{\sigma}_0\sqrt{n}}{\sigma_0} - \delta\sqrt{n}\right)
\]
и
\[
	p_U = \Phi\left(-(\hat{\mu}_0 - \mu_0)\frac{\sqrt{n}}{\sigma_0} - k\frac{\hat{\sigma}_0\sqrt{n}}{\sigma_0} + \delta\sqrt{n}\right)
\]
соответственно.

Определим случайные величины \(U = (\hat{\mu}_0 - \mu_0)\dfrac{\sqrt{n}}{\sigma_0}\) и \(V = \dfrac{\hat{\sigma}_0\sqrt{n}}{\sigma_0}\) соответственно. Теперь вероятности сигнала можно записать в более простом виде, а именно
\[
	p_L = \Phi(U - kV - \delta\sqrt{n}) \quad\text{и}\quad
	p_U = \Phi(-U - kV + \delta\sqrt{n})
\]
соответственно. Плотность случайной величины \(U\) равна \(f_U(u) = \varphi(u \mid 0, 1/\sqrt{m})\), где \(\varphi\) — плотность стандартного нормального распределения, поскольку
\[
	\hat{\mu}_0 \sim N(\mu_0, \sigma_0/\sqrt{mn}) \quad\text{и}\quad U \sim N(0, 1/\sqrt{m}).
\]
Для случайной величины \(V\) воспользуемся тем, что
\[
	\hat{\sigma}_0^2/\sigma_0^2 \sim \mathrm{GAM}(m(n - 1)/2,\ 2/(m(n - 1))),
\]
так что
\[
	V^2 \sim \mathrm{GAM}(m(n - 1)/2,\ 2n/(m(n - 1))),
\]
и потому плотность \(V\) задаётся как
\[
	f_V(v) = 2v\, f_\gamma(v^2 \mid m(n - 1)/2,\ 2n/(m(n - 1))),
\]
где \(f_\gamma(\cdot)\) — плотность гамма-распределения с параметрами \(m(n - 1)/2\) и \(2n/(m(n - 1))\) соответственно. Этот ход вывода основан на работе Zhang and Castagliola (2010), но можно также последовать Chakraborti (2000) и переписать уравнения через более привычное распределение хи-квадрат.

Теперь безусловные свойства распределения длины серии можно найти с помощью техники CUC, обсуждавшейся в Разделе 3.7. Идея, которую мы хотели здесь ввести, состоит в том, что техника CUC, обсуждавшаяся в Разделе 3.10.1, также удобно применима к картам с правилами сигнализации на основе серий. Отметим, что, как упоминалось ранее, поскольку случайную величину \(V\) можно масштабировать так, чтобы её распределение записывалось через распределение хи-квадрат, приведённые выше выражения можно переписать через распределение хи-квадрат с некоторым подходящим числом степеней свободы. Это может быть полезнее с практической точки зрения, поскольку плотность хи-квадрат более доступна. Проверку этих деталей мы оставляем читателю.

\section{Контрольные карты Фазы I}

\subsection{\(\overline{X}\)-карта Фазы I}

Как обсуждалось ранее, анализ Фазы I — очень важная часть общей схемы SPC, и контрольные карты, в особенности карты Шухарта, играют важную роль в мониторинге подгрупп Фазы I. В этом разделе мы сначала обсудим \(\overline{X}\)-карту Фазы I. Предположим, что имеется \(m\) подгрупп Фазы I объёма \(n\) каждая, и \(X_{ij}\) обозначает \(j\)-е наблюдение из \(i\)-й подгруппы. Пусть \(\overline{X}_i\) и \(S_i^2\) обозначают соответственно среднее и дисперсию \(i\)-й подгруппы. Статистика для нанесения на карту для \(\overline{X}\)-карты Фазы I — это среднее \(\overline{\overline{X}}\) каждой подгруппы, а контрольные границы задаются как
\[
	LCL/UCL = \overline{\overline{X}} \pm k\hat{\sigma}/\sqrt{n},
\]
где \(\overline{\overline{X}} = \dfrac{1}{mn}\sum_{i=1}^{m}\sum_{j=1}^{n} X_{ij}\) — общее (или генеральное) среднее, а
\[
	\hat{\sigma} = \frac{1}{c_4}\sqrt{\frac{1}{m}\sum_{i=1}^{m} S_i^2}
\]
— несмещённая оценка стандартного отклонения процесса, основанная на среднем дисперсий подгрупп с константой \(c_4\) (см. таблицу C приложения A). Константа \(k > 0\) — константа проектирования карты, которая определяется для заданного номинального значения FAP, \(m\) и \(n\). Champ and Jones (2004) рассмотрели эту задачу для ряда оценок стандартного отклонения, включая указанную выше \(\hat{\sigma}\), и показали, что \(k\) можно получить по функции распределения многомерного \(t\)-распределения со специальной корреляционной структурой. За подробностями заинтересованному читателю следует обратиться к их работе. Они затабулировали константу \(k\) для \(m = 4(1)10\), \(m = 15\) и \(n = 3(1)10\) при номинальной FAP, равной 0.01, 0.05 и 0.10. Однако в последнее время число подгрупп Фазы I \(m\) рекомендуется брать гораздо большим — скажем, 50 или даже 100 и более, чтобы надлежащим образом учесть оценивание параметров. Champ and Jones (2004) также обсудили некоторые приближения для \(k\) при больших \(m\); однако недавние достижения в компьютерном программном обеспечении для вычисления функции распределения многомерного нормального и многомерного \(t\)-распределения теперь позволяют вычислять \(k\) точно для больших значений \(m\), даже для \(m = 300\). Недавно Yao et al. (2017) расширили их таблицы на большие значения \(m\) и предоставили R-пакет для расчётов по запросу.

\noindent\textbf{Пример 3.14. Иллюстрация \(\overline{X}\)-карты Фазы I}

Для иллюстрации набор данных из работы Montgomery (2005, таблица 5.3, с. 223) о внутренних диаметрах поршневых колец, изготовленных в процессе ковки, используется во многих примерах этой книги. Мы используем его и здесь. Однако в наборе данных Montgomery (2005) имеется 25 подгрупп (\(m = 25\)) объёма 5 каждая (\(n = 5\)), а мы используем только первые 15 подгрупп, чтобы воспользоваться таблицами Champ and Jones (2004). Столбец (a) таблицы \ref{tab:3.24} содержит 15 ретроспективных подгрупп Фазы I объёма пять каждая, собранных, когда процесс считался находящимся в состоянии IC. Таким образом, \(N = mn = 15 \times 5 = 75\). Столбцы (b) и (c) таблицы \ref{tab:3.24} показывают средние и стандартные отклонения для \(i = 1, 2, \ldots, m\). Значение параметра проектирования \(k\) (в данном случае \(k = 2.6728\)) находится в таблице III работы Champ and Jones (2004) для \(m = 15\), \(n = 5\) и номинальной FAP = 0.10.

Графическое представление \(\overline{X}\)-карты Фазы I показано на рисунке \ref{fig-3.30}. Все статистики для нанесения на карту попадают между контрольными границами без заметного паттерна, что указывает на то, что процесс находится в состоянии IC. Теперь эти данные можно назвать данными IC, или эталонными данными, по которым можно оценить любые необходимые параметры и затем найти подходящие контрольные границы, используемые для проспективного мониторинга процесса в Фазе II.

{\small
\begin{longtable}{crrrrrrrr}
	\caption{Данные Фазы I для \(\overline{X}\)-карты и статистики для нанесения на карту.}
	\label{tab:3.24} \\
	\toprule
	& \multicolumn{5}{c}{(a)} & (b) & (c) \\
	\cmidrule(lr){2-6}
	Подгруппа, \(i\) & \(X_{i1}\) & \(X_{i2}\) & \(X_{i3}\) & \(X_{i4}\) & \(X_{i5}\) & \(\overline{X}_i\) & \(S_i\) \\
	\midrule
	\endfirsthead
	\toprule
	Подгруппа, \(i\) & \(X_{i1}\) & \(X_{i2}\) & \(X_{i3}\) & \(X_{i4}\) & \(X_{i5}\) & \(\overline{X}_i\) & \(S_i\) \\
	\midrule
	\endhead
	\bottomrule
	\endfoot
	1  & 74.030 & 74.002 & 74.019 & 73.992 & 74.008 & 74.010 & 0.015 \\
	2  & 73.995 & 73.992 & 74.001 & 74.011 & 74.004 & 74.001 & 0.008 \\
	3  & 73.998 & 74.024 & 74.021 & 74.005 & 74.002 & 74.010 & 0.012 \\
	4  & 74.002 & 73.996 & 73.993 & 74.015 & 74.009 & 74.003 & 0.009 \\
	5  & 73.992 & 74.007 & 74.015 & 73.989 & 74.014 & 74.003 & 0.012 \\
	6  & 74.009 & 73.994 & 73.997 & 73.985 & 73.993 & 73.996 & 0.009 \\
	7  & 73.995 & 74.006 & 73.994 & 74.000 & 74.005 & 74.000 & 0.006 \\
	8  & 73.985 & 74.003 & 73.993 & 74.015 & 73.988 & 73.997 & 0.012 \\
	9  & 74.008 & 73.995 & 74.009 & 74.005 & 74.004 & 74.004 & 0.006 \\
	10 & 73.998 & 74.000 & 73.990 & 74.007 & 73.995 & 73.998 & 0.006 \\
	11 & 73.994 & 73.998 & 73.994 & 73.995 & 73.990 & 73.994 & 0.003 \\
	12 & 74.004 & 74.000 & 74.007 & 74.000 & 73.996 & 74.001 & 0.004 \\
	13 & 73.983 & 74.002 & 73.998 & 73.997 & 74.012 & 73.998 & 0.010 \\
	14 & 74.006 & 73.967 & 73.994 & 74.000 & 73.984 & 73.990 & 0.015 \\
	15 & 74.012 & 74.014 & 73.998 & 73.999 & 74.007 & 74.006 & 0.007 \\
	16 & 74.000 & 73.984 & 74.005 & 73.998 & 73.996 & 73.997 & 0.008 \\
	17 & 73.994 & 74.012 & 73.986 & 74.005 & 74.007 & 74.001 & 0.011 \\
	18 & 74.006 & 74.010 & 74.018 & 74.003 & 74.000 & 74.007 & 0.007 \\
	19 & 73.984 & 74.002 & 74.003 & 74.005 & 73.997 & 73.998 & 0.008 \\
	20 & 74.000 & 74.010 & 74.013 & 74.020 & 74.003 & 74.009 & 0.008 \\
	21 & 73.982 & 74.001 & 74.015 & 74.005 & 73.996 & 74.000 & 0.012 \\
	22 & 74.004 & 73.999 & 73.990 & 74.006 & 74.009 & 74.002 & 0.007 \\
	23 & 74.010 & 73.989 & 73.990 & 74.009 & 74.014 & 74.002 & 0.012 \\
	24 & 74.015 & 74.008 & 73.993 & 74.000 & 74.010 & 74.005 & 0.009 \\
	25 & 73.982 & 73.984 & 73.995 & 74.017 & 74.013 & 73.998 & 0.016 \\
\end{longtable}
}
% CHECK: в книге написано «15 ретроспективных подгрупп» и k найдено для m = 15, но в таблице 3.24 напечатаны 25 подгрупп; воспроизведено как в книге.

\begin{figure}[!ht]
	\centering
	\includegraphics[width=0.7\linewidth]{"./images/fig-3.30.png"}
	\caption[]{\(\overline{X}\)-карта Фазы I.}
	\label{fig-3.30}
\end{figure}

О мониторинге дисперсий в Фазе I читателю следует обратиться к Human, Chakraborti and Smit (2010a). Напомним, что карты Фазы I применяются итеративно: исследуют и, возможно, отбрасывают подгруппы в состоянии OOC, пересчитывают границы и строят карту заново, пока все подгруппы не окажутся в состоянии IC.


\section{Объём данных Фазы I}

Эталонный набор данных Фазы I обычно состоит из \(m \ge 1\) подгрупп объёма \(n \ge 1\) каждая. Таким образом, для оценивания параметров и построения контрольных границ, используемых в Фазе II, имеется всего \(mn\) наблюдений Фазы I (эталонных данных). Существуют, как правило, два подхода к получению эталонного набора данных Фазы I. Он может уже иметься (как ретроспективный, или исторический, набор данных), либо процесс нужно сначала привести в состояние IC, а затем собрать эталонные данные. Приведение процесса в состояние IC и получение тем самым эталонных данных включает тщательный анализ Фазы I (где контрольные границы рассматриваются как пробные и используются итеративно, пока все статистики для нанесения на карту не окажутся в состоянии IC), в том числе применение контрольных карт и других статистических инструментов. Читателям следует обратиться к Jones-Farmer et al. (2014) и Chakraborti, Human and Graham (2009) за более подробной информацией о построении карт Фазы I. Суть в том, что эффекты оценивания параметров реальны, они различаются от практика к практику и часто приводят к большему числу ложных тревог, чем номинально ожидается. Если доступно большое количество эталонных данных, границы Фазы II могут работать так же, как их аналоги при известных параметрах. И наоборот, контрольные границы можно рассчитать под имеющиеся данные. К лежащим в основе вопросам относятся знание формы/вида лежащего в основе распределения, интересующий параметр,
тип оценки, тип карты и т.д. При оценивании изменчивости процесса, например, обычно требуется гораздо больший объём выборки, чем при оценивании среднего процесса. Напомним (из обсуждения анализа Фаз I и II в Разделе 2.1.12), что одна из главных целей анализа Фазы I — оценить изменчивость процесса.

Подводя итог: точное и аккуратное оценивание параметров в Фазе I критически важно для удовлетворительной работы контрольной карты в фазе мониторинга, то есть в Фазе II, поскольку свойства длины серии сильно зависят от оценивания, что может приводить, например, к гораздо большему числу ложных тревог, чем номинально ожидается исходя из заданной средней длины серии в состоянии IC; см., например, Jensen et al. (2006) и Psarakis, Vyniou and Castagliola (2014). Эти и другие авторы пришли к выводу, что число эталонных выборок Фазы I \(m\) часто должно быть довольно большим (во многих сотнях наблюдений), чтобы с разумной уверенностью можно было утверждать, что эффективность контрольной карты будет близка к известной эффективности в случае K или KK. Отметим, что «магического числа» для \(m\) не существует, поскольку выбор \(m\) в значительной степени зависит от типа применяемой контрольной карты и используемых оценок параметров. В литературе (например, Chakraborti, 2006; Saleh et al., 2015; Epprecht, Loureiro and Chakraborti, 2016) исследуются различные аспекты и показано, что для получения устойчивых свойств карты, близких к случаю известных параметров, требуется гораздо больший объём данных Фазы I — часто в несколько раз больше, чем обычно рекомендуется в учебниках. Случай известных параметров принят за эталон, поскольку эти значения не меняются. Поэтому нужно быть очень осторожным, начиная процедуру построения карты Фазы II при ограниченном объёме исторических данных. Другой подход — скорректировать контрольные границы под имеющийся объём данных. Исследования в этом направлении продолжаются.

\section{Робастность параметрических контрольных карт}

Робастная статистическая процедура — это процедура, которая хорошо работает не только в идеальных условиях (в которых она проектируется и математически выводится), но и при некотором отклонении от идеала. Важный вопрос о статистическом критерии проверки гипотез — робастность его ошибки I рода, поскольку без неё свойства мощности критерия сомнительны. В том же духе робастность в состоянии IC является важным свойством контрольной карты, которое нужно оценивать и учитывать, поскольку без неё способность карты обнаруживать сдвиги, по сути, теряет смысл. Для параметрических карт эта проблема серьезна, так как они обычно строятся при предположении о распределении (модели), и потому их свойства становятся сомнительными, когда это предположение не выполняется или не может быть адекватно обосновано в данной ситуации. Многие параметрические карты, включая карты Шухарта, CUSUM и EWMA, как показано, в той или иной степени страдают отсутствием робастности в состоянии IC при отклонении от нормальности. Следствие этого — сокращение средней длины серии в состоянии IC, причём настолько, что наблюдается гораздо больше ложных тревог, чем номинально ожидается. Это может снизить эффективность любой контрольной карты и навредить всей схеме мониторинга. Хотя некоторые авторы предполагали, что отдельные параметрические контрольные карты можно спроектировать (настроить) так, чтобы они были робастны к ряду распределений (см. Borror, Montgomery and Runger, 1999, о традиционной параметрической карте EWMA), это может оказаться проблематичным, поскольку знания о форме распределения и ожидаемом сдвиге, необходимые для настройки, не всегда доступны. В действительности Human, Kritzinger and Chakraborti (2011) показали, что традиционная параметрическая карта EWMA может терять робастность в состоянии IC для некоторых ненормальных распределений, таких как симметричное бимодальное и загрязнённое нормальное распределение. Их наблюдения ставят под вопрос рутинное применение традиционной параметрической карты EWMA на практике. Проблема робастности в состоянии IC может быть серьёзной заботой практиков, и в таких случаях можно применять непараметрические контрольные карты. Они робастны в состоянии IC по определению, поскольку их распределение длины серии в состоянии IC остаётся неизменным для всех непрерывных распределений или непрерывных и симметричных распределений. Это одно из главных практических преимуществ непараметрических карт. Мы рекомендуем: если предположение о распределении вызывает сомнения или не может быть обосновано из-за отсутствия доступной информации или данных, следует использовать непараметрическую, или свободную от распределения, контрольную карту. Это следует делать даже при возможности некоторой потери эффективности, когда истинная модель задана правильно, поскольку истинная модель обычно неизвестна.

\section*{Приложение 3.1. Некоторые выводы для карты EWMA}

\noindent\textbf{Результат 3.1. Дисперсия статистики для нанесения на карту EWMA}

Экспоненциально взвешенное скользящее среднее определяется как
\begin{equation}
	Z_i = \lambda X_i + (1 - \lambda)Z_{i-1} \quad \text{при } Z_0 = \mu_0.
	\tag{A3.1}
	\label{eq:A3.1}
\end{equation}
Рекурсивно раскрывая уравнение \eqref{eq:A3.1}, находим
\[
\begin{aligned}
	Z_i &= \lambda X_i + (1 - \lambda)Z_{i-1} \\
	    &= \lambda X_i + (1 - \lambda)(\lambda X_{i-1} + (1 - \lambda)Z_{i-2}) \\
	    &\;\;\vdots \\
	    &= \lambda \sum_{j=0}^{i-1}(1 - \lambda)^j X_{i-j} + (1 - \lambda)^i Z_0.
\end{aligned}
\]
Продолжая рекурсивную подстановку для \(Z_{i-j}\), \(j = 2, 3, \ldots, t\), получаем
\begin{equation}
	Z_i = \lambda \sum_{j=0}^{i-1}(1 - \lambda)^j X_{i-j} + (1 - \lambda)^i Z_0,
	\tag{A3.2}
	\label{eq:A3.2}
\end{equation}
так что
\[
\begin{aligned}
	\operatorname{var}(Z_i)
	&= \operatorname{var}\left(\lambda \sum_{j=0}^{i-1}(1 - \lambda)^j X_{i-j} + (1 - \lambda)^i Z_0\right) \\
	&= \lambda^2 \operatorname{var}\left(\sum_{j=0}^{i-1}(1 - \lambda)^j X_{i-j}\right)
	   \quad \text{поскольку } (1 - \lambda)^i Z_0 \text{ — константа} \\
	&= \lambda^2 \big(\operatorname{var}(X_i + (1 - \lambda)X_{i-1} + (1 - \lambda)^2 X_{i-2} \\
	&\qquad + \cdots + (1 - \lambda)^{i-2} X_2 + (1 - \lambda)^{i-1} X_1)\big) \\
	&= \lambda^2 \big(\sigma^2 + (1 - \lambda)^2 \sigma^2 + (1 - \lambda)^4 \sigma^2 \\
	&\qquad + \cdots + (1 - \lambda)^{2i-4}\sigma^2 + (1 - \lambda)^{2i-2}\sigma^2\big).
\end{aligned}
\]
Таким образом,
\begin{equation}
	\operatorname{var}(Z_i) = \lambda^2 \sigma^2 \left(1 + (1 - \lambda)^2 + (1 - \lambda)^4 + \cdots
	+ (1 - \lambda)^{2i-4} + (1 - \lambda)^{2i-2}\right).
	\tag{A3.3}
	\label{eq:A3.3}
\end{equation}
В общем случае для конечной геометрической прогрессии сумма первых \(n\) членов равна
\begin{equation}
	\sum_{j=0}^{n-1} a r^j = a\frac{1 - r^n}{1 - r} \quad \text{при } |r| < 1.
	\tag{A3.4}
	\label{eq:A3.4}
\end{equation}
Применяя этот результат к уравнению \eqref{eq:A3.3}, получаем
\[
\begin{aligned}
	\operatorname{var}(Z_i)
	&= \lambda^2 \sigma^2 \left(\frac{1 - ((1 - \lambda)^2)^i}{1 - (1 - \lambda)^2}\right) \\
	&= \lambda^2 \sigma^2 \left(\frac{1 - (1 - \lambda)^{2i}}{2\lambda - \lambda^2}\right) \\
	&= \lambda \sigma^2 \left(\frac{1 - (1 - \lambda)^{2i}}{2 - \lambda}\right).
\end{aligned}
\]
Таким образом,
\begin{equation}
	\operatorname{var}(Z_i) = \sigma^2 \frac{\lambda}{2 - \lambda}\left(1 - (1 - \lambda)^{2i}\right).
	\tag{A3.5}
	\label{eq:A3.5}
\end{equation}

\noindent\textbf{Результат 3.2. EWMA — несмещённая оценка среднего процесса}

Предположим, что процесс находится в состоянии IC со средним, равным \(\mu_0\). Используя рекурсивную подстановку, мы показали, что \(Z_i = \lambda \sum_{j=0}^{i-1}(1 - \lambda)^j X_{i-j} + (1 - \lambda)^i Z_0\). Предположим, что либо начальное значение \(Z_0\) равно нулю, либо \(i\) настолько велико, что начальное значение \(Z_0\) не оказывает влияния (равно нулю). Тогда имеем \(Z_i = \lambda \sum_{j=0}^{\infty}(1 - \lambda)^j X_{i-j}\). Следовательно,
\[
\begin{aligned}
	E(Z_i) &= E\left(\lambda \sum_{j=0}^{\infty}(1 - \lambda)^j X_{i-j}\right)
	       = \lambda \sum_{j=0}^{\infty}(1 - \lambda)^j E(X_{i-j})
\end{aligned}
\]
в силу свойства математического ожидания. Поскольку \(E(X_{i-j}) = \mu_0\),
\[
\begin{aligned}
	E(Z_i) &= \lambda \sum_{j=0}^{\infty}(1 - \lambda)^j \mu_0 \\
	       &= 1 \times \mu_0 = \mu_0
\end{aligned}
\]
с учётом того, что \(\sum_{j=0}^{\infty}(1 - \lambda)^j = 1/\lambda\), что следует из свойств бесконечной геометрической прогрессии. Этот результат для суммы бесконечной геометрической прогрессии можно получить из \eqref{eq:A3.4}, переходя к пределу суммы при \(n \to \infty\).

\section*{Приложение 3.2. Цепи Маркова}

\noindent\textbf{Определение 3.1.} Стохастический процесс — это набор случайных величин \(\{Y_t : t \in T\}\). Все возможные значения, которые может принимать \(\{Y_t, t \in T\}\), называются состояниями стохастического процесса, а множество всех возможных состояний называется пространством состояний и обозначается \(\Omega\). Множество всех индексирующих параметров \(T\) называется пространством параметров. Если \(t\) обозначает время, то значение \(\{Y_t\}\) в момент времени \(t\) называется состоянием, в котором процесс находится в момент \(t\).

\noindent\textbf{Определение 3.2.} Последовательность \(\{Y_t\}\) называется цепью Маркова (MC), если для любой последовательности
\[
	\{Y_0 = \omega_0, Y_1 = \omega_1, \ldots, Y_{t-1} = \omega_{t-1}, Y_t = \omega_t\}, \quad t = 1, 2, 3, \ldots,
\]
верно, что
\[
	P(Y_t = \omega_t \mid Y_{t-1} = \omega_{t-1}, \ldots, Y_0 = \omega_0) = P(Y_t = \omega_t \mid Y_{t-1} = \omega_{t-1}),
\]
где \(\{\omega_0, \omega_1, \omega_2, \ldots, \omega_t\} \in \Omega\).

Таким образом, последовательность \(\{Y_t\}\) является MC, если вероятность того, что система перейдёт в состояние \(\omega_t\) в момент \(t\), зависит только от непосредственно предшествующего состояния \(\omega_{t-1}\), в котором система находилась в момент \(t - 1\). В контексте SPC система — это процесс, находящийся под мониторингом.

\noindent\textbf{Определение 3.3.} На каждом шаге (то есть в каждый момент времени) процесс \(\{Y_t, t = 1, 2, 3, \ldots\}\) может остаться в текущем состоянии или перейти в другое состояние в соответствии с некоторым распределением вероятностей. Изменения состояния называются переходами, а вероятности, связанные с этими переходами, — переходными вероятностями. Эти вероятности обозначаются как
\[
	P(Y_t = \omega_j \mid Y_{t-1} = \omega_i) = p_{i,j}(t), \quad \text{где } \omega_j, \omega_i \in \Omega,
\]
в момент \(t\). Переходные вероятности \(p_{i,j}(t)\), \(1 \le i, j \le m\), могут быть представлены в виде матрицы \(m \times m\)
\[
	M_t = [p_{i,j}(t)] =
	\begin{bmatrix}
		p_{1,1}(t) & p_{1,2}(t) & \cdots & p_{1,m}(t) \\
		p_{2,1}(t) & \cdots     & \cdots & \cdots     \\
		\vdots     & \cdots     & \cdots & \cdots     \\
		p_{m,1}(t) & p_{m,2}(t) & \cdots & p_{m,m}(t)
	\end{bmatrix}_{m \times m}.
\]
Матрицы \(M_t\), \(t = 1, 2, 3, \ldots\), — это одношаговые матрицы переходных вероятностей. Отметим, что эта матрица переходных вероятностей зависит от времени; это означает, что \(p_{i,j}(t)\) не обязательно равна \(p_{i,j}(s)\) при \(t \ne s\).

\noindent\textbf{Определение 3.4.} Цепь Маркова \(\{Y_0, Y_1, Y_2, \ldots\}\) однородна, если переходные вероятности постоянны во времени (то есть не зависят от времени), а именно
\[
	P(Y_t = \omega_j \mid Y_{t-1} = \omega_i) = p_{i,j}
\]
для любых \(\omega_i, \omega_j \in \Omega\) и всех \(t = 1, 2, 3, \ldots\), где переходные вероятности \(p_{i,j}\) не зависят от индекса времени \(t\).

\noindent\textbf{Определение 3.5.} Состояние \(\omega_i \in \Omega\) цепи Маркова называется поглощающим, если, попав в состояние \(\omega_i\), система уже никогда не может его покинуть, то есть \(p_{\omega_i,\omega_i} \equiv 1\) и \(p_{\omega_i,\omega_j} \equiv 0\) для любого \(i \ne j\).

\noindent\textbf{Определение 3.6.} Состояние \(\omega_j \in \Omega\) цепи Маркова называется невозвратным, если при условии, что процесс находится в состоянии \(\omega_j\), существует положительная вероятность того, что процесс никогда не вернётся в состояние \(\omega_j\).

\noindent\textbf{Интерпретация поглощающих и невозвратных состояний на контрольной карте.} Перемещение цепи Маркова между невозвратным(и) состоянием(ями) интерпретируется как (связано с) стохастическое движение статистики для нанесения на карту по контрольной карте без подачи сигнала. Это указывает на то, что процесс находится (остаётся) в состоянии IC, и состояние(я), связанное(ые) с пребыванием процесса в состоянии IC, отождествляется(ются) с невозвратным(и) состоянием(ями). Существует положительная вероятность того, что статистика для нанесения на карту в любой момент времени окажется в области OOC, в результате чего карта подаст сигнал. Последнее эквивалентно утверждению, что стохастический процесс \(\{Y_t\}\) перешёл в поглощающее состояние (после чего процесс уже не может вернуться ни в одно из невозвратных состояний). Это означает, что как только на карте появляется сигнал, причина сигнала выявляется и устраняется, а мониторинг процесса начинается заново из состояния IC.

\noindent\textbf{Доказательство: цепи Маркова.}

Из математической статистики: дифференцируя производящую функцию моментов \(M_X(t)\) любой случайной величины \(X\) по \(t\) и полагая \(t = 0\), можно найти моменты случайной величины \(X\), а именно
\[
	\frac{d}{dt} M_X(t)\Big|_{t=0} = M_X'(0) = E(X), \qquad
	\frac{d^2}{dt^2} M_X(t)\Big|_{t=0} = M_X''(0) = E(X^2).
\]
В контексте SPC одна из ключевых интересующих случайных величин — это длина серии \(N\). Из теоремы 7.4 работы Fu and Lou (2003, с. 143) имеем
\[
	M_N(t) = (e^t - 1)\xi (I - e^t Q)^{-1}\mathbf{1} + 1,
\]
так что
\[
\begin{aligned}
	\frac{d}{dt} M_N(t)
	&= \frac{d}{dt}\left[(e^t - 1)\xi (I - e^t Q)^{-1}\mathbf{1} + 1\right] \\
	&= e^t \xi (I - e^t Q)^{-1}\mathbf{1}
	   + (e^t - 1)\xi (I - e^t Q)^{-2}(e^t Q)\mathbf{1}.
\end{aligned}
\]
% CHECK: в оригинале промежуточные строки первой производной набраны с опечатками; приведён корректный вид, совпадающий с итоговой строкой книги.
Далее
\[
\begin{aligned}
	\frac{d^2}{dt^2} M_N(t)
	&= \frac{d}{dt} M_N'(t) \\
	&= e^t \xi (I - e^t Q)^{-1}\mathbf{1}
	   + e^t \xi (I - e^t Q)^{-2}(e^t Q)\mathbf{1} \\
	&\quad + (2e^{2t} - e^t)\xi (I - e^t Q)^{-2}Q\mathbf{1}
	   + 2(e^{3t} - e^{2t})\xi (I - e^t Q)^{-3}Q^2\mathbf{1},
\end{aligned}
\]
так что при подстановке \(t = 0\)
\[
\begin{aligned}
	M_N'(0)
	&= \xi (I - Q)^{-1}\mathbf{1} + (1 - 1)\xi (I - Q)^{-2}Q\mathbf{1} \\
	&= \xi (I - Q)^{-1}\mathbf{1}.
\end{aligned}
\]
Аналогично
\[
\begin{aligned}
	M_N''(0)
	&= \xi (I - Q)^{-1}\mathbf{1} + \xi (I - Q)^{-2}Q\mathbf{1}
	   + (2 - 1)\xi (I - Q)^{-2}Q\mathbf{1} \\
	&= \xi (I - Q)^{-1}\mathbf{1} + 2\xi (I - Q)^{-2}Q\mathbf{1} \\
	&= \xi\big(I + 2(I - Q)^{-1}Q\big)(I - Q)^{-1}\mathbf{1} \\
	&= \xi (I + Q)(I - Q)^{-2}\mathbf{1}.
\end{aligned}
\]
Следовательно, первые два момента длины серии \(N\) равны
\[
	E(N) = M_N'(0) = \xi (I - Q)^{-1}\mathbf{1}, \qquad
	E(N^2) = M_N''(0) = \xi (I + Q)(I - Q)^{-2}\mathbf{1},
\]
откуда
\[
\begin{aligned}
	\operatorname{var}(N)
	&= E(N^2) - (E(N))^2 \\
	&= \xi (I + Q)(I - Q)^{-2}\mathbf{1} - \big(\xi (I - Q)^{-1}\mathbf{1}\big)^2, \\
	\operatorname{SD}(N)
	&= \sqrt{\operatorname{var}(N)}
	 = \sqrt{\xi (I + Q)(I - Q)^{-2}\mathbf{1} - \big(\xi (I - Q)^{-1}\mathbf{1}\big)^2}.
\end{aligned}
\]

\section*{Приложение 3.3. Некоторые выводы для карт Шухарта по разбросу}

Константы для нанесения на карту для карт Шухарта по дисперсии можно получить в два шага. Сначала, приравняв безусловную среднюю длину серии в состоянии IC некоторому номинальному значению, решаем уравнение относительно вероятности \(\alpha(m,n)\). Затем на втором шаге константы для нанесения на карту (а значит, и контрольные границы) для карт разброса Фазы II получаются из соответствующих процентилей распределения хи-квадрат (точно или приближённо) с \(n - 1\) степенями свободы. Эти границы являются безусловными вероятностными границами. Здесь мы кратко изложим детали на примере \(S\)-карты с \(S_p\) в качестве оценки \(\sigma\) по Фазе I, используя метод CUC. Отметим, что в разделе о распределении длины серии мы показали, что вероятностные границы задаются как \(H_3^*(m,n) = \sqrt{\chi^2_{\alpha(m,n)/2,\,n-1}/(n - 1)}\) и \(H_4^*(m,n) = \sqrt{\chi^2_{1-\alpha(m,n)/2,\,n-1}/(n - 1)}\), где \(\alpha(m,n)\) — вероятность того, что статистика для нанесения на карту \(S_i\) окажется вне нижней или верхней контрольной границы, когда процесс находится в состоянии IC. Таким образом, чтобы получить константы для нанесения на карту \(H_3^*\) и \(H_4^*\), нужны значения \(\alpha(m,n)\).

Для этого прежде всего отметим, что
\[
\begin{aligned}
	\alpha(m,n) &= 1 - P\big(H_3^*(m,n)S_p < S_i < H_4^*(m,n)S_p \mid \text{IC}\big) \\
	&= 1 - P\left(\sqrt{\frac{\chi^2_{\alpha(m,n)/2,\,n-1}}{n-1}}\,S_p < S_i
	< \sqrt{\frac{\chi^2_{1-\alpha(m,n)/2,\,n-1}}{n-1}}\,S_p \;\middle|\; \text{IC}\right),
\end{aligned}
\]
что можно переписать как
\[
\begin{aligned}
	\alpha(m,n)
	&= 1 - P\left(\frac{\chi^2_{\alpha(m,n)/2,\,n-1}\,S_p^2}{n-1} < S_i^2
	< \frac{\chi^2_{1-\alpha(m,n)/2,\,n-1}\,S_p^2}{n-1} \;\middle|\; \text{IC}\right) \\
	&= 1 - P\left(\frac{\chi^2_{\alpha(m,n)/2,\,n-1}\,S_p^2}{(n-1)\sigma^2} < \frac{S_i^2}{\sigma^2}
	< \frac{\chi^2_{1-\alpha(m,n)/2,\,n-1}\,S_p^2}{(n-1)\sigma^2} \;\middle|\; \text{IC}\right) \\
	&= 1 - P\left(\chi^2_{\alpha(m,n)/2,\,n-1} < \frac{(n-1)S_i^2/\sigma^2}{U/m}
	< \chi^2_{1-\alpha(m,n)/2,\,n-1} \;\middle|\; \text{IC}\right) \\
	&= 1 - P\left(\chi^2_{\alpha(m,n)/2,\,n-1} < \frac{V}{U/m}
	< \chi^2_{1-\alpha(m,n)/2,\,n-1}\right),
\end{aligned}
\]
где \(V = \dfrac{(n-1)S_i^2}{\sigma^2}\) и \(U = \dfrac{m(n-1)S_p^2}{\sigma^2}\) — независимые случайные величины, каждая из которых в состоянии IC имеет распределение хи-квадрат с \((n-1)\) и \(m(n-1)\) степенями свободы соответственно.

Последнюю вероятность можно вычислить, обусловливая на \(S_p\) или, что эквивалентно, на \(U\). Это называется условной частотой ложных тревог (CFAR), которая является функцией \(U\), \(m\) и \(n\). Следовательно, CFAR можно записать как
\begin{equation}
\begin{aligned}
	&\mathrm{CFAR}(U, m, n, \alpha(m,n)) \\
	&\quad = 1 - F_{\chi^2_{n-1}}\left(\chi^2_{\alpha(m,n)/2,\,n-1}\frac{U}{m}\right)
	       + F_{\chi^2_{n-1}}\left(\chi^2_{1-\alpha(m,n)/2,\,n-1}\frac{U}{m}\right),
\end{aligned}
\tag{A3.6}
\label{eq:A3.6}
\end{equation}
так что условное распределение длины серии является геометрическим с вероятностью CFAR\((U, m, n, p)\). Таким образом, применяя метод CUC, безусловное распределение длины серии в состоянии IC (pmf) задаётся как
\[
	P(N = k) = \int_0^{\infty} \mathrm{CFAR}(U, m, n, \alpha(m,n))
	\big[1 - \mathrm{CFAR}(U, m, n, \alpha(m,n))\big]^{k-1}
	f_{\chi^2_{m(n-1)}}(u)\,du, \quad k = 1, 2, \ldots
\]
Отметим, что это безусловное распределение длины серии не является геометрическим.

Аналогичным образом, используя метод CUC, условная средняя длина серии в состоянии IC равна \([\mathrm{CFAR}(U, m, n, \alpha(m,n))]^{-1}\), и потому безусловная средняя длина серии в состоянии IC задаётся как
\begin{equation}
	ARL_{\text{IC}} = \int_0^{\infty} \big[\mathrm{CFAR}(U, m, n, \alpha(m,n))\big]^{-1}
	f_{\chi^2_{m(n-1)}}(u)\,du.
	\tag{A3.7}
	\label{eq:A3.7}
\end{equation}
И снова отметим, что \(ARL_{\text{IC}}\) не является обратной величиной к CFAR. Теперь, в безусловном методе, приравнивая выражение для \(ARL_{\text{IC}}\) некоторому номинальному значению, например 370, мы численно решаем уравнение \(ARL_{\text{IC}} = 370\) относительно \(\alpha(m,n)\) для заданных значений \(m\) и \(n\). Затем находим \(100\cdot\alpha(m,n)/2\)-й и \(100\cdot[1 - \alpha(m,n)/2]\)-й процентили распределения хи-квадрат с \(n - 1\) степенями свободы. Искомые константы для нанесения на карту \(H_3^*(m,n)\) и \(H_4^*(m,n)\) получаются делением каждого из этих процентилей на \(n - 1\) и извлечением квадратного корня. Эти значения приведены в таблице H приложения A. Проверку остальной части этой таблицы для других карт и других оценок \(\sigma\) по Фазе I мы оставляем читателю.
% CHECK: в оригинале в первом процентиле, по-видимому, пропущено деление на 2; приведено в соответствии с формулами для H_3^* и H_4^*.

Наконец, отметим, что эта формулировка допускает объём выборки Фазы II, скажем \(n_1\), не равный объёму выборки Фазы I \(m\). Разница в этом случае будет в степенях свободы распределения \(V\), которые теперь будут равны \(n_1 - 1\) в первом распределении хи-квадрат при определении CFAR в уравнении \eqref{eq:A3.6}.
